\documentclass[11pt]{article}

\usepackage[margin=1in]{geometry}
\usepackage{amsmath,amssymb,amsthm,mathtools,mathrsfs,bm}
\usepackage{enumitem}
\usepackage{hyperref}
\usepackage{microtype}

\allowdisplaybreaks

\newtheorem{theorem}{Theorem}[section]
\newtheorem{proposition}[theorem]{Proposition}
\newtheorem{lemma}[theorem]{Lemma}
\newtheorem{corollary}[theorem]{Corollary}
\theoremstyle{definition}
\newtheorem{definition}[theorem]{Definition}
\newtheorem{assumption}[theorem]{Assumption}
\theoremstyle{remark}
\newtheorem{remark}[theorem]{Remark}

\newcommand{\R}{\mathbb{R}}
\newcommand{\E}{\mathbb{E}}
\newcommand{\Law}{\mathrm{Law}}
\newcommand{\KL}{\mathrm{KL}}
\newcommand{\Ent}{\mathrm{Ent}}
\newcommand{\TV}{\mathrm{TV}}
\newcommand{\dd}{\,\mathrm{d}}
\newcommand{\mcP}{\mathcal{P}}
\newcommand{\mcL}{\mathcal{L}}
\newcommand{\mcH}{\mathcal{H}}
\newcommand{\mcF}{\mathcal{F}}
\newcommand{\mcK}{\mathcal{K}}
\newcommand{\mcV}{\mathcal{V}}
\newcommand{\one}{\mathbf{1}}

\title{Convergence, Hypocoercivity, and Mean-Field Limit\\
for the Kinetic Core of the Geometric Gas}
\author{}
\date{}

\begin{document}
\maketitle

\begin{abstract}
This paper isolates the pure kinetic backbone of the geometric gas: a confining
underdamped Langevin dynamics on $\R^d$, augmented by a row-normalized viscous
velocity alignment force and followed, at observation times, by a smooth
velocity squashing map. The cloning operator, companion sampling, fitness
potential, and state-dependent anisotropic diffusion are deliberately omitted.
Within this kinetic-only setting we prove the two structural facts that make the
model tractable: the viscous force is uniformly bounded on velocity-capped
states, and it is dissipative on relative kinetic energy. We then formulate the
finite-particle hypocoercive entropy estimate in a small-viscosity regime that
yields an $N$-uniform
logarithmic Sobolev inequality and exponential convergence in relative entropy
and total variation. Finally we derive the McKean-Vlasov limit for the viscous
alignment field, prove quantitative propagation of chaos in relative entropy,
and pass the same small-viscosity $N$-uniform logarithmic Sobolev inequality to the mean-field
equilibrium. The paper is self-contained in the sense that all assumptions,
definitions, and theorems used below are stated explicitly in this document.
\end{abstract}

\tableofcontents

\section{Introduction}

The purpose of this paper is to extract and organize the convergence theory of
the kinetic part of the geometric gas in a form that does not depend on the
cloning stage or on the adaptive diffusion geometry. The model studied here is
the following.

\begin{itemize}[leftmargin=1.5em]
\item The position variable evolves by transport: $\dot x = v$.
\item The velocity variable is subject to a confining force $-\nabla U(x)$,
linear friction $-\gamma v$, isotropic Brownian noise, and a row-normalized
velocity alignment force.
\item After each observation interval one may apply a smooth radial squashing
map $\psi$ to cap the observed speed without destroying the previously obtained
probabilistic estimates.
\end{itemize}

Two themes drive the analysis.

\begin{enumerate}[leftmargin=1.5em]
\item The \emph{backbone hypocoercive mechanism} comes from underdamped
Langevin dynamics with confinement and friction.
\item The \emph{new viscous term} is harmless at the level of stability because
its action is bounded and dissipative on relative velocity disagreements.
\end{enumerate}

The paper therefore proceeds in two layers. First we prove the structural
lemmas for the viscous operator and the velocity squashing map. Then we graft
those facts onto the hypocoercive framework for the kinetic backbone,
obtaining $N$-uniform entropy and total-variation convergence in a
small-viscosity regime. The final
sections pass to the mean-field equation, prove a quantitative propagation of
chaos estimate, and deduce a mean-field logarithmic Sobolev inequality in the
same regime. The
analysis uses only tools from Langevin dynamics, hypocoercivity,
Harris-type ergodic theory, and McKean-Vlasov limits
\cite{roberts1996exponential,meyn2012markov,villani2009hypocoercivity,sznitman1991topics}.

\section{The Kinetic Geometric Gas Core}

\subsection{Standing assumptions}

\begin{assumption}[Strongly confining potential]\label{ass:potential}
The potential $U:\R^d\to\R$ satisfies:
\begin{enumerate}[label=\textnormal{(\roman*)},leftmargin=1.5em]
\item $U\in C^2(\R^d)$ and there exist constants $0<m_U\leq L_U<\infty$ such
that
\[
\|\nabla U(x)-\nabla U(y)\|\leq L_U\|x-y\|
\qquad
\forall x,y\in\R^d.
\]
\item The Hessian is uniformly positive in the sense
\[
\langle \xi,\nabla^2 U(x)\xi\rangle\geq m_U\|\xi\|^2
\qquad
\forall x,\xi\in\R^d.
\]
\item $U$ is coercive in the Lyapunov sense:
\[
\langle x,\nabla U(x)\rangle \geq \alpha_U\|x\|^2-\beta_U
\qquad
\forall x\in\R^d
\]
for some $\alpha_U>0$ and $\beta_U<\infty$.
\end{enumerate}
\end{assumption}

\begin{remark}
Assumption~\ref{ass:potential} for items $i$ and $ii$ places the paper in the classical
uniformly convex Langevin regime. In this regime the Gibbs measure
\eqref{eq:gibbs} satisfies a logarithmic Sobolev inequality by the
Bakry-\'Emery criterion, and the modified-entropy computation below yields an
explicit hypocoercive decay rate for the kinetic backbone; see
\cite{bakry1985diffusions,villani2009hypocoercivity}.
\end{remark}

\begin{assumption}[Regularized Gaussian kernel]\label{ass:kernel}
Fix $\ell_{\mathrm{visc}}>0$, $\kappa_0>0$, and define
\[
K(x,y):=\kappa_0+\exp\!\left(-\frac{\|x-y\|^2}{2\ell_{\mathrm{visc}}^2}\right).
\]
Then $K$ is symmetric, strictly positive, bounded, and globally Lipschitz in
each variable. Set
\[
k_*:=\inf_{x,y}K(x,y)=\kappa_0,
\qquad
k^*:=\sup_{x,y}K(x,y)=1+\kappa_0.
\]
\end{assumption}

\begin{definition}[Finite-particle viscous weights]\label{def:weights}
For $N\geq2$ and a configuration $X=(x_1,\dots,x_N)\in(\R^d)^N$, define
\[
d_i(X):=\sum_{k\neq i}K(x_i,x_k),
\qquad
\omega_{ij}(X):=\frac{K(x_i,x_j)}{d_i(X)}
\quad(j\neq i).
\]
Then $\omega_{ij}\geq0$ and $\sum_{j\neq i}\omega_{ij}=1$.
\end{definition}

\begin{definition}[Finite-$N$ kinetic core]\label{def:finiteN}
Fix $\gamma>0$, $\sigma>0$, and $\nu_{\mathrm{visc}}\geq0$. For $i=1,\dots,N$,
the finite-particle kinetic geometric gas core is the It\^o SDE
\begin{align}
dX_i(t) &= V_i(t)\,dt,\label{eq:finite-x}\\
dV_i(t) &=
\Bigl(-\nabla U(X_i(t))-\gamma V_i(t)+F_i(X(t),V(t))\Bigr)\,dt
+\sigma\,dW_i(t),\label{eq:finite-v}
\end{align}
where the viscous force is
\begin{equation}\label{eq:finite-visc}
F_i(X,V)
:=
\nu_{\mathrm{visc}}\sum_{j\neq i}\omega_{ij}(X)\bigl(V_j-V_i\bigr).
\end{equation}
The Brownian motions $W_1,\dots,W_N$ are independent standard
$d$-dimensional Wiener processes.
\end{definition}

\begin{definition}[Mean-field viscous field]\label{def:mean-field}
Let $f\in\mcP(\R^{2d})$ with spatial marginal
\[
\rho_f(dx):=\int_{\R^d}f(dx,dq).
\]
Define the normalized kernel mass
\[
m_f(x):=\int_{\R^d}K(x,y)\,\rho_f(dy)
\]
and the mean-field viscous field
\begin{equation}\label{eq:mf-force}
F[f](x,v)
:=
\frac{\nu_{\mathrm{visc}}}{m_f(x)}
\iint_{\R^{2d}}K(x,y)(q-v)\,f(dy,dq).
\end{equation}
The mean-field kinetic equation is
\begin{equation}\label{eq:mf-pde}
\partial_t f_t
+v\cdot\nabla_x f_t
+\nabla_v\cdot\Bigl[\bigl(-\nabla U(x)-\gamma v+F[f_t](x,v)\bigr)f_t\Bigr]
=\frac{\sigma^2}{2}\Delta_v f_t.
\end{equation}
\end{definition}

\begin{remark}
Because $m_f(x)\geq k_*>0$ for every probability measure $f$,
Definition~\ref{def:mean-field} is globally well posed at the algebraic level:
there is no denominator singularity in the mean-field alignment field.
\end{remark}

\begin{definition}[Velocity squashing]\label{def:squash}
Fix $V_{\mathrm{alg}}>0$. The smooth velocity squashing map is
\[
\psi(v):=
V_{\mathrm{alg}}\frac{v}{V_{\mathrm{alg}}+\|v\|},
\qquad
\psi(0):=0.
\]
The associated observation map on phase space is
\[
S(x,v):=(x,\psi(v)).
\]
\end{definition}

\begin{definition}[BAOAB observation chain]\label{def:baoab}
Fix a time step \(\Delta t>0\). The algorithmic update used at observation
times is the BAOAB Strang splitting of the kinetic core:
\[
Z_{n+1}^{\mathrm{pre}}
=
\Phi_B^{\Delta t/2}\circ \Phi_A^{\Delta t/2}\circ \Phi_O^{\Delta t}
\circ \Phi_A^{\Delta t/2}\circ \Phi_B^{\Delta t/2}(Z_n),
\]
where \(Z_n=(X_n,V_n)\) and the substeps are the explicit maps
\[
\Phi_B^h(X,V)
:=
\bigl(X,\,V+h\,G(X,V)\bigr),
\qquad
\Phi_A^h(X,V)
:=
\bigl(X+hV,\,V\bigr).
\]
Here \(G(X,V)=(G_1(X,V),\dots,G_N(X,V))\) with
\[
G_i(X,V):=-\nabla U(X_i)+F_i(X,V).
\]
The Ornstein-Uhlenbeck substep is
\[
\Phi_O^h(X,V)
:=
\left(X,\,e^{-\gamma h}V+\sqrt{\frac{\sigma^2}{2\gamma}\bigl(1-e^{-2\gamma h}\bigr)}\,\Xi\right),
\]
where \(\Xi\sim\mathcal N(0,I)\) is an independent standard Gaussian vector.
The viscous force \(F_i\) is the row-normalized coupling from
\eqref{eq:finite-visc}. The observed state is then postprocessed
componentwise by
\[
Z_{n+1}=S(Z_{n+1}^{\mathrm{pre}})=(X_{n+1}^{\mathrm{pre}},\psi(V_{n+1}^{\mathrm{pre}})).
\]
This is the same BAOAB ordering used in the implementation and in the
convergence-program notes already present elsewhere in the repository.
\end{definition}

\section{Structural Bounds: Confinement, Viscosity, and Squashing}

\begin{lemma}[Smooth bounded cap]\label{lem:squash}
The map $\psi$ of Definition~\ref{def:squash} is smooth on $\R^d\setminus\{0\}$,
continuous on $\R^d$, satisfies $\|\psi(v)\|\leq V_{\mathrm{alg}}$, and is
$1$-Lipschitz:
\[
\|\psi(v)-\psi(w)\|\leq \|v-w\|
\qquad
\forall v,w\in\R^d.
\]
\end{lemma}

\begin{proof}
Write $r=\|v\|$. For $v\neq0$,
\[
\psi(v)=\phi(r)v,
\qquad
\phi(r)=\frac{V_{\mathrm{alg}}}{V_{\mathrm{alg}}+r}.
\]
Hence
\[
D\psi(v)
=
\phi(r)I
+\phi'(r)\frac{v\otimes v}{r}
=
\frac{V_{\mathrm{alg}}}{V_{\mathrm{alg}}+r}I
-\frac{V_{\mathrm{alg}}}{(V_{\mathrm{alg}}+r)^2}\frac{v\otimes v}{r}.
\]
The radial eigenvalue is
$V_{\mathrm{alg}}^2/(V_{\mathrm{alg}}+r)^2\leq1$ and each tangential eigenvalue
is $V_{\mathrm{alg}}/(V_{\mathrm{alg}}+r)\leq1$. Therefore
\[
\|D\psi(v)\|_{\mathrm{op}}\leq1\qquad (v\neq0).
\]
Any \(u\in\R^d\) decomposes as \(u=\alpha v/r+u_\perp\) with \(u_\perp\perp v\),
so
\[
\|D\psi(v)u\|
\leq
\frac{V_{\mathrm{alg}}^2}{(V_{\mathrm{alg}}+r)^2}|\alpha|+\frac{V_{\mathrm{alg}}}{V_{\mathrm{alg}}+r}\|u_\perp\|
\leq\|u\|.
\]
Hence \(\psi\) is \(1\)-Lipschitz on \(\R^d\setminus\{0\}\) by the mean-value theorem
on line segments. At the origin,
\[
\lim_{v\to0}\|D\psi(v)\|_{\mathrm{op}}=1,
\]
so the same Lipschitz bound extends continuously to \(v=0\).
The speed cap is obvious from
\[
\|\psi(v)\|=\frac{V_{\mathrm{alg}}\|v\|}{V_{\mathrm{alg}}+\|v\|}
\leq V_{\mathrm{alg}}.
\]
\qedhere
\end{proof}

\begin{lemma}[Uniform bound on the finite-particle viscous force]\label{lem:visc-bound}
If $\max_{1\leq k\leq N}\|V_k\|\leq V_*$, then for every $i$
\[
\|F_i(X,V)\|
\leq 2\nu_{\mathrm{visc}}V_*.
\]
The bound is independent of $N$.
\end{lemma}

\begin{proof}
By Definition~\ref{def:weights},
\[
F_i(X,V)
=
\nu_{\mathrm{visc}}
\left(\sum_{j\neq i}\omega_{ij}(X)V_j-V_i\right).
\]
Taking norms and using $\sum_{j\neq i}\omega_{ij}=1$ gives
\[
\|F_i(X,V)\|
\leq
\nu_{\mathrm{visc}}
\left(\sum_{j\neq i}\omega_{ij}(X)\|V_j\|+\|V_i\|\right)
\leq2\nu_{\mathrm{visc}}V_*.
\qedhere
\]
\end{proof}

\begin{lemma}[Dissipativity of the viscous coupling]\label{lem:visc-diss}
Fix the positions $X$ and consider only the viscous ODE
\[
\dot V_i
=
\nu_{\mathrm{visc}}\sum_{j\neq i}\omega_{ij}(X)(V_j-V_i).
\]
Define
\[
D(X):=\sum_{i=1}^N d_i(X),
\qquad
\bar V_d:=\frac{1}{D(X)}\sum_{i=1}^N d_i(X)V_i,
\]
and the degree-weighted velocity variance
\[
\mathscr V_d(V)
:=
\frac{1}{D(X)}
\sum_{i=1}^N d_i(X)\|V_i-\bar V_d\|^2.
\]
Then
\begin{equation}\label{eq:visc-diss}
\frac{d}{dt}\mathscr V_d(V(t))
=
-\frac{2\nu_{\mathrm{visc}}}{D(X)}
\sum_{1\leq i<j\leq N}K(x_i,x_j)\|V_i-V_j\|^2
\leq0.
\end{equation}
\end{lemma}

\begin{proof}
Since $X$ is frozen, the degrees $d_i(X)$ are constant in time. Set
$\delta_i:=V_i-\bar V_d$. By symmetry of $K$,
\[
\frac{d}{dt}\bar V_d
=
\frac{\nu_{\mathrm{visc}}}{D(X)}
\sum_{i=1}^N\sum_{j\neq i}K(x_i,x_j)(V_j-V_i)
=0,
\]
so $\dot\delta_i=\dot V_i$. Therefore
\begin{align*}
\frac{d}{dt}\mathscr V_d(V(t))
&=
\frac{2}{D(X)}\sum_{i=1}^Nd_i(X)\langle \delta_i,\dot V_i\rangle\\
&=
\frac{2\nu_{\mathrm{visc}}}{D(X)}
\sum_{i=1}^N\sum_{j\neq i}K(x_i,x_j)\langle \delta_i,\delta_j-\delta_i\rangle\\
&=
-\frac{\nu_{\mathrm{visc}}}{D(X)}
\sum_{i\neq j}K(x_i,x_j)\|\delta_i-\delta_j\|^2\\
&=
-\frac{2\nu_{\mathrm{visc}}}{D(X)}
\sum_{1\leq i<j\leq N}K(x_i,x_j)\|V_i-V_j\|^2,
\end{align*}
which is non-positive.
\end{proof}

\begin{lemma}[Lipschitz control of the mean-field alignment field]\label{lem:mf-lipschitz}
Let $\mu,\nu\in\mcP(\R^{2d})$ have finite first moments. There exists a constant
$L_{\mathrm{mf}}<\infty$, depending on
$\nu_{\mathrm{visc}},k_*,k^*,\ell_{\mathrm{visc}}$ and on a fixed first-moment
bound for the measure class under consideration, such that
\[
\|F[\mu](x,v)-F[\nu](x',v')\|
\leq
L_{\mathrm{mf}}
\bigl(\|x-x'\|+\|v-v'\|+W_1(\mu,\nu)\bigr),
\]
where $W_1$ denotes the $1$-Wasserstein distance on $\R^{2d}$.
\end{lemma}

\begin{proof}
Write
\[
A[\mu](x):=\frac{1}{m_\mu(x)}\iint K(x,y)q\,\mu(dy,dq),
\qquad
F[\mu](x,v)=\nu_{\mathrm{visc}}\bigl(A[\mu](x)-v\bigr).
\]
To estimate the denominator, use the coupling characterization of $W_1$ and the
global Lipschitz bound on $K$:
\[
|m_\mu(x)-m_\nu(x')|
\leq
\iint |K(x,y)-K(x',y')|\,\pi(dy,dq,dy',dq')
\leq
C_K\bigl(\|x-x'\|+W_1(\mu,\nu)\bigr),
\]
for any coupling $\pi\in\Pi(\mu,\nu)$ and some constant $C_K$ depending only on
$L_K$. Since $m_\mu,m_\nu\geq k_*$, the reciprocal map satisfies
\[
\bigl|m_\mu(x)^{-1}-m_\nu(x')^{-1}\bigr|
\leq
k_*^{-2}|m_\mu(x)-m_\nu(x')|.
\]
For the numerator, fix a coupling $\pi\in\Pi(\mu,\nu)$ and write
\[
J_\mu(x):=\iint K(x,y)q\,\mu(dy,dq).
\]
Then
\[
\|J_\mu(x)-J_\nu(x')\|
\leq
\iint \|K(x,y)q-K(x',y')q'\|\,d\pi.
\]
On any class of laws with a fixed first-moment bound in the velocity variable,
the integrand is controlled by a constant times
\(\|x-x'\|+\|y-y'\|+\|q-q'\|\), so
\[
\|J_\mu(x)-J_\nu(x')\|
\leq
C_J\bigl(\|x-x'\|+W_1(\mu,\nu)\bigr),
\]
with \(C_J\) depending only on the kernel bounds and the moment bound for the
class under consideration. Combining the numerator and denominator estimates,
and using
\[
A[\mu](x)-A[\nu](x')
=
\frac{J_\mu(x)-J_\nu(x')}{m_\mu(x)}
+J_\nu(x')\bigl(m_\mu(x)^{-1}-m_\nu(x')^{-1}\bigr),
\]
gives
\[
\|A[\mu](x)-A[\nu](x')\|
\leq
C\bigl(\|x-x'\|+W_1(\mu,\nu)\bigr).
\]
Finally,
\[
F[\mu](x,v)-F[\nu](x',v')
=
\nu_{\mathrm{visc}}\bigl(A[\mu](x)-A[\nu](x')\bigr)
-\nu_{\mathrm{visc}}(v-v'),
\]
so
\[
\|F[\mu](x,v)-F[\nu](x',v')\|
\leq
L_{\mathrm{mf}}\bigl(\|x-x'\|+\|v-v'\|+W_1(\mu,\nu)\bigr).
\qedhere
\]
\end{proof}

\begin{proposition}[Finite- and infinite-dimensional well posedness]\label{prop:wp}
Under Assumptions~\ref{ass:potential} and~\ref{ass:kernel},
\begin{enumerate}[label=\textnormal{(\roman*)},leftmargin=1.5em]
\item for every $N\geq2$ the SDE
\eqref{eq:finite-x}--\eqref{eq:finite-v} has a unique global strong solution;
\item for every initial law $f_0\in\mcP(\R^{2d})$ with finite second moment, the
McKean-Vlasov equation \eqref{eq:mf-pde} has a unique global solution
$t\mapsto f_t$ in $C(\R_+,\mcP_2(\R^{2d}))$.
\end{enumerate}
\end{proposition}

\begin{proof}
Fix $N$ and write the drift as
\[
b_N(X,V)=\bigl(V,-\nabla U(X)-\gamma V+F(X,V)\bigr).
\]
Because the kernel $K$ is smooth and strictly positive, every denominator
\[
d_i(X)=\sum_{k\neq i}K(x_i,x_k)
\]
is bounded below by \((N-1)k_*\), so each normalized weight \(\omega_{ij}(X)\)
is smooth in \(X\). Hence \(b_N\) is locally Lipschitz on \((\R^{2d})^N\), and
standard SDE theory gives a unique strong solution up to the explosion time. To
rule out explosion, define
\[
\mathcal E_N(X,V):=\frac1N\sum_{i=1}^N \bigl(\|X_i\|^2+\|V_i\|^2\bigr).
\]
Applying It\^o's formula gives
\begin{align*}
d\mathcal E_N
&=
\frac{2}{N}\sum_{i=1}^N \langle X_i,V_i\rangle\,dt
+\frac{2}{N}\sum_{i=1}^N \left\langle V_i,-\nabla U(X_i)-\gamma V_i+F_i(X,V)\right\rangle dt\\
&\qquad
+d\sigma^2\,dt+\frac{2\sigma}{N}\sum_{i=1}^N \langle V_i,dW_i\rangle.
\end{align*}
Use Young's inequality,
\[
2\langle X_i,V_i\rangle \leq \|X_i\|^2+\|V_i\|^2,
\qquad
-2\langle V_i,\nabla U(X_i)\rangle
\leq
\frac{1}{m_U}\|V_i\|^2+m_U\|\nabla U(X_i)\|^2,
\]
and the linear-growth bound
\[
\|\nabla U(x)\|^2
\leq
2L_U^2\|x\|^2+2\|\nabla U(0)\|^2.
\]
For the viscous term, Definition~\ref{def:weights} and
$\omega_{ij}\leq k^*/((N-1)k_*)$ imply
\begin{align*}
\sum_{i=1}^N \langle V_i,F_i(X,V)\rangle
&\leq
\nu_{\mathrm{visc}}\sum_{i=1}^N |V_i|
\left(\sum_{j\neq i}\omega_{ij}(X)|V_j|+|V_i|\right)\\
&\leq
\nu_{\mathrm{visc}}\sum_{i=1}^N |V_i|^2
+\frac{\nu_{\mathrm{visc}}}{2}\sum_{i=1}^N\sum_{j\neq i}\omega_{ij}(X)
\bigl(|V_i|^2+|V_j|^2\bigr)\\
&\leq
\nu_{\mathrm{visc}}\left(
\frac32+\frac{k^*}{2k_*}
\right)\sum_{i=1}^N |V_i|^2.
\end{align*}
Consequently
\[
\frac{d}{dt}\E\mathcal E_N(t)
\leq
C_0+C_1\,\E\mathcal E_N(t)
\]
with constants depending only on
$(m_U,L_U,\gamma,\sigma,\nu_{\mathrm{visc}},k_*,k^*)$ and not on $N$. Gr\"onwall's
lemma then gives $\sup_{t\leq T}\E\mathcal E_N(t)<\infty$ on every finite
interval, which excludes explosion. Hence the strong solution is global.

For the McKean-Vlasov equation, Lemma~\ref{lem:mf-lipschitz} gives the local
Lipschitz dependence of the drift on \((x,v,\mu)\) on moment-bounded classes,
while the force has at most linear growth in \((x,v)\) under a finite second
moment assumption on \(\mu\). A Picard iteration on the nonlinear SDE
\[
d\bar X_t=\bar V_t\,dt,
\qquad
d\bar V_t=
\bigl(-\nabla U(\bar X_t)-\gamma \bar V_t+F[f_t](\bar X_t,\bar V_t)\bigr)\,dt
+\sigma\,dW_t
\]
combined with the same Lyapunov estimate as above yields a unique global strong
solution with law \(f_t\in C(\R_+,\mcP_2(\R^{2d}))\).
\end{proof}

\begin{proposition}[Quadratic moment bounds]\label{prop:moments}
Let $(X_i,V_i)_{i=1}^N$ solve \eqref{eq:finite-x}--\eqref{eq:finite-v} and let
$f_t$ solve \eqref{eq:mf-pde}. Then for every finite horizon $T>0$ there exists
$C_T<\infty$, depending only on the model parameters and the initial second
moments, such that
\[
\sup_{t\in[0,T]}
\frac1N\sum_{i=1}^N \E\bigl[\|X_i(t)\|^2+\|V_i(t)\|^2\bigr]
\leq C_T
\]
for every $N\geq2$, and
\[
\sup_{t\in[0,T]}
\int_{\R^{2d}} (\|x\|^2+\|v\|^2)\,f_t(dx,dv)
\leq C_T.
\]
\end{proposition}

\begin{proof}
For the finite-particle system, the estimate from Proposition~\ref{prop:wp}
already gives a closed differential inequality for the averaged quadratic
energy, so the uniform bound follows by taking expectations and applying
Gr\"onwall's lemma.

For the mean-field equation, apply It\^o's formula to
\(\|\bar X_t\|^2+\|\bar V_t\|^2\) along the nonlinear SDE.\@ Since
\[
\|F[f_t](x,v)\|
\leq
\frac{\nu_{\mathrm{visc}}k^*}{k_*}
\left(\int \|q\|\,f_t(dy,dq)+\|v\|\right),
\]
Young's inequality gives
\[
\frac{d}{dt}\E\bigl[\|\bar X_t\|^2+\|\bar V_t\|^2\bigr]
\leq
C_0+C_1\E\bigl[\|\bar X_t\|^2+\|\bar V_t\|^2\bigr].
\]
Applying Gr\"onwall's lemma yields the claimed bound.
\end{proof}

\section{Detailed Backbone and Viscous Derivations}

\subsection{Weighted generator calculus for the uncoupled kinetic backbone}

For the detailed hypocoercive calculations it is convenient to work first with
the uncoupled kinetic backbone. Set
\[
T:=\frac{\sigma^2}{2\gamma}
\]
and define the Gibbs density
\begin{equation}\label{eq:gibbs}
\rho_*(x,v)
:=
Z_*^{-1}\exp\!\left(-\frac{U(x)+\frac12\|v\|^2}{T}\right).
\end{equation}
Under the smooth confining assumptions on $U$, this is the invariant density of
the one-particle underdamped Langevin diffusion.

\begin{definition}[Relative density and hypocoercive functionals]\label{def:functionals}
Let $f_t$ solve the uncoupled one-particle Fokker-Planck equation
\[
\partial_t f_t
+v\cdot \nabla_x f_t
+\nabla_v\cdot\bigl((\nabla U(x)+\gamma v)f_t\bigr)
=\frac{\sigma^2}{2}\Delta_v f_t,
\]
and write
\[
h_t(x,v):=\frac{f_t(x,v)}{\rho_*(x,v)}.
\]
For smooth positive $h$ define
\begin{align}
\mathcal A[h]
&:=
\int h\log h\,d\rho_*,
\label{eq:A-def}\\
\mathcal B[h]
&:=
T\int \frac{|\nabla_v h|^2}{h}\,d\rho_*
=
T\int h\,|\nabla_v \log h|^2\,d\rho_*,
\label{eq:B-def}\\
\mathcal C[h]
&:=
T\int \frac{\nabla_x h\cdot \nabla_v h}{h}\,d\rho_*
=
T\int h\,\langle \nabla_x\log h,\nabla_v\log h\rangle\,d\rho_*,
\label{eq:C-def}\\
\mathcal D[h]
&:=
T\int \frac{|\nabla_x h|^2}{h}\,d\rho_*
=
T\int h\,|\nabla_x \log h|^2\,d\rho_*.
\label{eq:D-def}
\end{align}
\end{definition}

\begin{proposition}[Exact generator decomposition]\label{prop:generator}
Let $L_0$ denote the generator acting on $h$ through
\[
\partial_t h=L_0 h=\rho_*^{-1}\partial_t(h\rho_*).
\]
Then
\[
L_0=L_{\mathrm{sym}}+L_{\mathrm{anti}},
\]
where
\begin{align}
L_{\mathrm{sym}}h
&=
\gamma\bigl(T\Delta_v h-v\cdot \nabla_v h\bigr),
\label{eq:Lsym}\\
L_{\mathrm{anti}}h
&=
-v\cdot \nabla_x h+\nabla U(x)\cdot \nabla_v h.
\label{eq:Lanti}
\end{align}
Moreover:
\begin{enumerate}[label=\textnormal{(\roman*)},leftmargin=1.5em]
\item $L_{\mathrm{sym}}$ is symmetric and non-positive in $L^2(\rho_*)$;
\item $L_{\mathrm{anti}}$ is anti-symmetric in $L^2(\rho_*)$.
\end{enumerate}
\end{proposition}

\begin{proof}
By definition,
\[
L_0 h
=
\frac{1}{\rho_*}\left[
-v\cdot \nabla_x(h\rho_*)
+\nabla_v\cdot\bigl((\nabla U+\gamma v)h\rho_*\bigr)
+\frac{\sigma^2}{2}\Delta_v(h\rho_*)
\right].
\]
We first expand the diffusion-friction part. Using
\[
\nabla_v \rho_*=-\frac{v}{T}\rho_*,
\qquad
\frac{\sigma^2}{2}=\gamma T,
\]
we obtain
\begin{align*}
\nabla_v\cdot\bigl((\nabla U+\gamma v)h\rho_*\bigr)
+\gamma T\Delta_v(h\rho_*)
&=
\nabla U\cdot \nabla_v(h\rho_*)+\gamma\nabla_v\cdot(vh\rho_*)+\gamma T\Delta_v(h\rho_*)\\
&=
\nabla U\cdot \nabla_v(h\rho_*)
+\gamma \rho_*\,v\cdot \nabla_v h
+\gamma d\,h\rho_*
+\gamma v\cdot \nabla_v\rho_*\,h\\
&\qquad
+\gamma T\Delta_v(h\rho_*).
\end{align*}
The terms involving $v h\rho_*$ cancel after expanding
$\Delta_v(h\rho_*)$ and using $\nabla_v \rho_* = -T^{-1}v\rho_*$. What remains
is
\[
\nabla U\cdot \nabla_v(h\rho_*)
+\gamma T\nabla_v\cdot(\rho_*\nabla_v h).
\]
Dividing by $\rho_*$ gives
\[
\nabla U\cdot \nabla_v h+\gamma\bigl(T\Delta_v h-v\cdot \nabla_v h\bigr).
\]
The transport part expands as
\[
-\frac{1}{\rho_*}v\cdot \nabla_x(h\rho_*)
=
-v\cdot \nabla_x h-v\cdot \nabla_x(\log \rho_*)\,h.
\]
Since $\nabla_x \log \rho_*=-T^{-1}\nabla U$, the last term is
$T^{-1}(v\cdot \nabla U)h$, which cancels exactly with the corresponding term
created when expanding $\nabla U\cdot \nabla_v(h\rho_*)/\rho_*$. This proves
\eqref{eq:Lsym}--\eqref{eq:Lanti}.

For symmetry, let $g,h$ be smooth compactly supported functions. Then
\[
\int g\,L_{\mathrm{sym}}h\,d\rho_*
=
\gamma T\int g\,\rho_*^{-1}\nabla_v\cdot(\rho_*\nabla_v h)\,d\rho_*
=
-\gamma T\int \nabla_v g\cdot \nabla_v h\,d\rho_*,
\]
which is symmetric in $(g,h)$ and non-positive on the diagonal.

For anti-symmetry, write
\[
b(x,v):=(v,-\nabla U(x)).
\]
Then
\[
L_{\mathrm{anti}}h=-b\cdot \nabla_{x,v}h.
\]
It is enough to prove that $\nabla_{x,v}\cdot (b\rho_*)=0$. Indeed,
\[
\nabla_{x,v}\cdot (b\rho_*)
=
v\cdot \nabla_x \rho_*-\nabla U\cdot \nabla_v \rho_*
=
-\frac{1}{T}(v\cdot \nabla U)\rho_*+\frac{1}{T}(\nabla U\cdot v)\rho_*=0.
\]
Therefore
\[
\int gL_{\mathrm{anti}}h\,d\rho_*
=
-\int g\,b\cdot \nabla h\,\rho_*
=
\int h\,b\cdot \nabla g\,\rho_*
=
-\int hL_{\mathrm{anti}}g\,d\rho_*.
\qedhere
\]
\end{proof}

\begin{proposition}[Commutators]\label{prop:comm}
For smooth $h$,
\begin{align}
[\nabla_v,L_0]h&=-\gamma \nabla_v h-\nabla_x h,
\label{eq:comm-v}\\
[\nabla_x,L_0]h&=(\nabla^2 U)\nabla_v h.
\label{eq:comm-x}
\end{align}
\end{proposition}

\begin{proof}
The symmetric part depends only on $v$, so $[\nabla_x,L_{\mathrm{sym}}]=0$ and
\[
[\nabla_v,L_{\mathrm{sym}}]h
=
\gamma\bigl(T\nabla_v\Delta_v h-\nabla_v(v\cdot \nabla_v h)+v\cdot \nabla_v\nabla_v h\bigr)
=
-\gamma \nabla_v h.
\]
For the anti-symmetric part,
\[
[\nabla_v,-v\cdot \nabla_x]h=-\nabla_x h,
\qquad
[\nabla_v,\nabla U\cdot \nabla_v]h=0,
\]
which yields \eqref{eq:comm-v}. Likewise,
\[
[\nabla_x,-v\cdot \nabla_x]h=0,
\qquad
[\nabla_x,\nabla U\cdot \nabla_v]h=(\nabla^2 U)\nabla_v h,
\]
which is \eqref{eq:comm-x}.
\end{proof}

\begin{proposition}[Exact differential identities]\label{prop:abcd}
Let $h_t$ solve $\partial_t h_t=L_0 h_t$. Set
\[
M:=\sup_{x\in \R^d}\|\nabla^2 U(x)\|_{\mathrm{op}}\leq L_U.
\]
Then:
\begin{align}
\frac{d}{dt}\mathcal A[h_t]
&=
-\gamma \mathcal B[h_t],
\label{eq:A-ode}\\
\frac{d}{dt}\mathcal B[h_t]
&\leq
-2\gamma \mathcal B[h_t]-2\mathcal C[h_t],
\label{eq:B-ode}\\
\frac{d}{dt}\mathcal C[h_t]
&\leq
-\mathcal D[h_t]-\gamma \mathcal C[h_t]
+M\sqrt{\mathcal B[h_t]\mathcal D[h_t]},
\label{eq:C-ode}\\
\frac{d}{dt}\mathcal D[h_t]
&\leq
2M\sqrt{\mathcal B[h_t]\mathcal D[h_t]}.
\label{eq:D-ode}
\end{align}
\end{proposition}

\begin{proof}
Set $\ell_t:=\log h_t$.

For \eqref{eq:A-ode}, the anti-symmetric part drops out by
Proposition~\ref{prop:generator}, item $ii$, so
\[
\frac{d}{dt}\mathcal A[h_t]
=
\int (1+\log h_t)L_{\mathrm{sym}}h_t\,d\rho_*
=
-\gamma T\int \frac{|\nabla_v h_t|^2}{h_t}\,d\rho_*
=
-\gamma \mathcal B[h_t].
\]

We next differentiate
\[
\mathcal B[h_t]=T\int h_t |\nabla_v\ell_t|^2\,d\rho_*.
\]
Since $\partial_t \ell_t=h_t^{-1}L_0 h_t$, the chain rule gives
\[
\frac{d}{dt}\mathcal B[h_t]
=
\underbrace{T\int L_0 h_t\,|\nabla_v\ell_t|^2\,d\rho_*}_{=:I_1}
+\underbrace{2T\int h_t \,\nabla_v\ell_t\cdot \nabla_v(\partial_t \ell_t)\,d\rho_*}_{=:I_2}.
\]
Because
\[
\nabla_v(\partial_t \ell_t)
=
\nabla_v\!\left(\frac{L_0 h_t}{h_t}\right)
=
\frac{\nabla_v L_0 h_t}{h_t}
-\frac{L_0 h_t}{h_t}\nabla_v\ell_t,
\]
we obtain
\[
I_2
=
2T\int \nabla_v\ell_t\cdot \nabla_vL_0 h_t\,d\rho_*
-2T\int L_0 h_t\,|\nabla_v\ell_t|^2\,d\rho_*.
\]
Therefore
\[
\frac{d}{dt}\mathcal B[h_t]
=
2T\int \nabla_v\ell_t\cdot \nabla_vL_0 h_t\,d\rho_*
-T\int L_0 h_t\,|\nabla_v\ell_t|^2\,d\rho_*.
\]
Split $L_0=L_{\mathrm{sym}}+L_{\mathrm{anti}}$ and treat both pieces
separately.

For the symmetric piece, use
\[
L_{\mathrm{sym}}=\gamma T\nabla_v^*\nabla_v,
\]
where $\nabla_v^*$ is the $L^2(\rho_*)$ adjoint of $\nabla_v$. A direct
Bakry-\'Emery computation yields
\[
-2\gamma T^2\int h_t\|\nabla_v^2\log h_t\|_{\mathrm{HS}}^2\,d\rho_*\leq0,
\]
and the first-order contribution is exactly $-2\gamma \mathcal B[h_t]$.

For the anti-symmetric piece, integrate by parts using
Proposition~\ref{prop:generator}, item $ii$:
\begin{align*}
&2T\int \nabla_v\ell_t\cdot \nabla_vL_{\mathrm{anti}}h_t\,d\rho_*
-T\int L_{\mathrm{anti}} h_t\,|\nabla_v\ell_t|^2\,d\rho_*\\
&\qquad=
2T\int h_t \,\nabla_v\ell_t\cdot [\nabla_v,L_{\mathrm{anti}}]\ell_t\,d\rho_*.
\end{align*}
By \eqref{eq:comm-v}, $[\nabla_v,L_{\mathrm{anti}}]=-\nabla_x$, hence this term
equals $-2\mathcal C[h_t]$. Combining the symmetric and anti-symmetric pieces
gives
\[
\frac{d}{dt}\mathcal B[h_t]
=
-2\gamma \mathcal B[h_t]
-2\mathcal C[h_t]
-2\gamma T^2\int h_t\|\nabla_v^2\log h_t\|_{\mathrm{HS}}^2\,d\rho_*,
\]
and \eqref{eq:B-ode} follows by discarding the final non-positive term.

We now turn to
\[
\mathcal C[h_t]=T\int h_t\langle \nabla_x\ell_t,\nabla_v\ell_t\rangle\,d\rho_*.
\]
Differentiating and grouping terms gives
\begin{align*}
\frac{d}{dt}\mathcal C[h_t]
&=
T\int L_0 h_t \,\langle \nabla_x\ell_t,\nabla_v\ell_t\rangle\,d\rho_*\\
&\qquad
+T\int h_t\langle \nabla_x(\partial_t \ell_t),\nabla_v\ell_t\rangle\,d\rho_*\\
&\qquad
+T\int h_t\langle \nabla_x\ell_t,\nabla_v(\partial_t \ell_t)\rangle\,d\rho_*.
\end{align*}
Proceed exactly as for $\mathcal B$: rewrite
$\partial_t\ell_t=h_t^{-1}L_0 h_t$, integrate by parts, and use the symmetry of
$L_{\mathrm{sym}}$ and anti-symmetry of $L_{\mathrm{anti}}$. All second-order
terms from $L_{\mathrm{sym}}$ contribute a remainder $\mathcal R_C(t)\leq0$.
The commutator terms are the only first-order contributions, so
\begin{align*}
\frac{d}{dt}\mathcal C[h_t]
&\leq
T\int h_t\langle \nabla_x\ell_t,[\nabla_v,L_0]\ell_t\rangle\,d\rho_*\\
&\qquad
+T\int h_t\langle \nabla_v\ell_t,[\nabla_x,L_0]\ell_t\rangle\,d\rho_*.
\end{align*}
Applying \eqref{eq:comm-v}--\eqref{eq:comm-x} gives
\[
\frac{d}{dt}\mathcal C[h_t]
\leq
-\mathcal D[h_t]-\gamma \mathcal C[h_t]
+T\int h_t\langle \nabla_x\log h_t,(\nabla^2U)\nabla_v\log h_t\rangle\,d\rho_*.
\]
Using the operator bound on $\nabla^2 U$ and Cauchy-Schwarz,
\[
\left|
T\int h_t\langle \nabla_x\log h_t,(\nabla^2U)\nabla_v\log h_t\rangle\,d\rho_*
\right|
\leq
M\sqrt{\mathcal B[h_t]\mathcal D[h_t]},
\]
which gives \eqref{eq:C-ode}.

Finally, starting from
\[
\mathcal D[h_t]=T\int h_t |\nabla_x\ell_t|^2\,d\rho_*,
\]
the same differentiation yields
\[
\frac{d}{dt}\mathcal D[h_t]
=
2T\int h_t\langle \nabla_x\log h_t,[\nabla_x,L_0]\log h_t\rangle\,d\rho_*
+\mathcal R_D(t),
\]
with $\mathcal R_D(t)\leq0$ from the symmetric part, because $L_{\mathrm{sym}}$
acts only in velocity. Invoking \eqref{eq:comm-x} and Cauchy-Schwarz yields
\eqref{eq:D-ode}.
\end{proof}

\begin{theorem}[Detailed backbone hypocoercive decay]\label{thm:backbone-detailed}
Let $M=L_U$ and define
\[
\varepsilon:=\min\left\{1,\frac{\gamma}{16(M+\gamma)^2}\right\},
\qquad
a:=\frac{\gamma\varepsilon}{4},
\qquad
b:=\varepsilon,
\qquad
c:=\varepsilon^2.
\]
For the modified entropy
\begin{equation}\label{eq:Phi}
\Phi[h]
:=
\mathcal A[h]+a\mathcal B[h]+2b\mathcal C[h]+c\mathcal D[h],
\end{equation}
there exists $\lambda_{\mathrm{back}}>0$, depending only on
$(\gamma,M)$, such that
\[
\frac{d}{dt}\Phi[h_t]
\leq
-\lambda_{\mathrm{back}}\bigl(\mathcal B[h_t]+\mathcal D[h_t]\bigr).
\]
If the Gibbs measure $\rho_*$ satisfies an LSI with constant $\rho_{\mathrm{LSI}}>0$,
that is,
\[
\rho_{\mathrm{LSI}}\mathcal A[h]\leq \mathcal B[h]+\mathcal D[h],
\]
then
\[
\Phi[h_t]\leq e^{-\lambda_{\mathrm{back}}\rho_{\mathrm{LSI}}t}\Phi[h_0].
\]
\end{theorem}

\begin{proof}
Insert \eqref{eq:A-ode}--\eqref{eq:D-ode} into the derivative of
\eqref{eq:Phi}. This gives
\begin{align*}
\frac{d}{dt}\Phi[h_t]
&\leq
-\gamma\mathcal B
+a(-2\gamma \mathcal B-2\mathcal C)
+2b(-\mathcal D-\gamma\mathcal C+M\sqrt{\mathcal B\mathcal D})\\
&\qquad
+c(2M\sqrt{\mathcal B\mathcal D}).
\end{align*}
Using $-\mathcal C\leq \sqrt{\mathcal B\mathcal D}$, we obtain
\[
\frac{d}{dt}\Phi[h_t]
\leq
-\gamma(1+2a)\mathcal B
-2b\mathcal D
+\Xi\,\sqrt{\mathcal B\mathcal D},
\]
where
\[
\Xi:=2a+2b\gamma+2M(b+c).
\]
Set $Y=(\sqrt{\mathcal B},\sqrt{\mathcal D})^\top$. The right-hand side is
$-Y^\top K Y$ for
\[
K=
\begin{pmatrix}
\gamma(1+2a) & -\Xi/2\\[0.3em]
-\Xi/2 & 2b
\end{pmatrix}.
\]
Since $a,b,c$ are positive, $K_{11}\geq \gamma$ and $K_{22}=2\varepsilon$.
Moreover,
\[
\Xi
\leq
\frac{\gamma^2\varepsilon}{2}+2\gamma\varepsilon+2(M+\gamma)\varepsilon
\leq 4(M+\gamma)\varepsilon,
\]
because $\varepsilon\leq1$. Hence
\[
\det K
\geq
2\gamma\varepsilon-4(M+\gamma)^2\varepsilon^2
\geq
\gamma\varepsilon>0
\]
by the choice of $\varepsilon$. Therefore $K$ is positive definite and
\[
Y^\top K Y\geq \lambda_{\min}(K)(\mathcal B+\mathcal D)
\]
where
\[
\lambda_{\min}(K)
=\frac{\gamma(1+2a)+2b-\sqrt{\bigl(\gamma(1+2a)-2b\bigr)^2+\Xi^2}}{2}.
\]
with
\[
\lambda_{\mathrm{back}}:=\lambda_{\min}(K)>0.
\]
This proves the first claim. If $\rho_*$ satisfies the stated LSI, then
\[
\mathcal B[h_t]+\mathcal D[h_t]\geq \rho_{\mathrm{LSI}}\mathcal A[h_t].
\]
Since $\Phi$ is equivalent to $\mathcal A+\mathcal B+\mathcal D$ for fixed
$(a,b,c)$, the differential inequality closes and Gr\"onwall's lemma gives the
exponential decay of $\Phi$.
\end{proof}

\begin{proposition}[Gibbs logarithmic Sobolev inequality]\label{prop:gibbs-lsi}
Let
\[
W(x,v):=U(x)+\frac12\|v\|^2,
\qquad
m_*:=\min\{m_U,1\}.
\]
Then the Gibbs measure $\rho_*$ from \eqref{eq:gibbs} satisfies
\begin{equation}\label{eq:gibbs-lsi}
\Ent_{\rho_*}(h)
\leq
\frac{1}{2m_*}\bigl(\mathcal B[h]+\mathcal D[h]\bigr)
\end{equation}
for every smooth density $h\geq0$ with $\int h\,d\rho_*=1$.
\end{proposition}

\begin{proof}
The density of $\rho_*$ has the form
\[
\rho_*(x,v)=Z_*^{-1}e^{-W(x,v)/T}.
\]
The Hessian of $W$ is block diagonal:
\[
\nabla^2 W(x,v)=
\begin{pmatrix}
\nabla^2 U(x) & 0\\
0 & I_d
\end{pmatrix}.
\]
Assumption~\ref{ass:potential}, item $ii$, therefore implies
\[
\nabla^2 W(x,v)\succeq m_* I_{2d}
\qquad
\forall (x,v)\in\R^{2d}.
\]
The Bakry-\'Emery criterion applied to the probability measure
$Z_*^{-1}e^{-W/T}dx\,dv$ yields
\[
\Ent_{\rho_*}(g^2)\leq \frac{2T}{m_*}\int |\nabla g|^2\,d\rho_*
\]
for every smooth $g$. Choose $g=\sqrt h$. Since
\[
|\nabla \sqrt h|^2
=
\frac{|\nabla_x h|^2+|\nabla_v h|^2}{4h},
\]
we obtain
\[
\Ent_{\rho_*}(h)
\leq
\frac{T}{2m_*}
\int \frac{|\nabla_x h|^2+|\nabla_v h|^2}{h}\,d\rho_*
=
\frac{1}{2m_*}\bigl(\mathcal B[h]+\mathcal D[h]\bigr),
\]
which is \eqref{eq:gibbs-lsi}.
\end{proof}

\begin{corollary}[One-particle backbone entropy contraction]\label{cor:backbone-gap}
Let $P_t^{(1),0}$ be the semigroup of the uncoupled one-particle kinetic
backbone and let $\pi_1=\rho_*$. There exists a constant
$\lambda_{\mathrm{hypo}}>0$, depending only on
$(m_U,L_U,\gamma,\sigma)$, such that
\[
\KL(\mu P_t^{(1),0}\,\|\,\pi_1)
\leq
e^{-2\lambda_{\mathrm{hypo}}t}\KL(\mu\,\|\,\pi_1)
\]
for every probability measure $\mu\ll\pi_1$ and every $t\geq0$.
\end{corollary}

\begin{proof}
Let $h_t=d(\mu P_t^{(1),0})/d\pi_1$. Proposition~\ref{prop:gibbs-lsi} yields
\[
\mathcal A[h]\leq \frac{1}{2m_*}\bigl(\mathcal B[h]+\mathcal D[h]\bigr).
\]
By Cauchy-Schwarz,
\[
2|\mathcal C[h]|
\leq
\mathcal B[h]+\mathcal D[h].
\]
Therefore the quadratic part of $\Phi$ satisfies
\[
\Phi[h]
\leq
\left(\frac{1}{2m_*}+a+b+c\right)\bigl(\mathcal B[h]+\mathcal D[h]\bigr)
=:
C_\Phi\bigl(\mathcal B[h]+\mathcal D[h]\bigr).
\]
Theorem~\ref{thm:backbone-detailed} gives
\[
\frac{d}{dt}\Phi[h_t]
\leq
-\lambda_{\mathrm{back}}\bigl(\mathcal B[h_t]+\mathcal D[h_t]\bigr)
\leq
-\frac{\lambda_{\mathrm{back}}}{C_\Phi}\Phi[h_t].
\]
Hence
\[
\Phi[h_t]\leq e^{-(\lambda_{\mathrm{back}}/C_\Phi)t}\Phi[h_0].
\]
Since $\mathcal A[h_t]\leq \Phi[h_t]$ and
\[
\Phi[h_0]
\leq
\left(1+\frac{a+c+2b}{2m_*}\right)\bigl(\mathcal B[h_0]+\mathcal D[h_0]\bigr),
\]
the same differential inequality applied from time $s$ onward and integrated
for all $s\geq0$, integrating from $s$ to $\infty$ gives
\[
\mathcal A[h_s]
=
\gamma\int_s^\infty \mathcal B[h_t]\,dt
\leq
\gamma \int_s^\infty \bigl(\mathcal B[h_t]+\mathcal D[h_t]\bigr)\,dt
\leq
\frac{\gamma C_\Phi}{\lambda_{\mathrm{back}}}\Phi[h_s].
\]
Using $\mathcal A[h_s]\leq \Phi[h_s]$ once more and absorbing constants yields
an exponential bound
\[
\mathcal A[h_t]\leq e^{-2\lambda_{\mathrm{hypo}}t}\mathcal A[h_0]
\]
for some $\lambda_{\mathrm{hypo}}>0$ depending only on
$(m_U,L_U,\gamma,\sigma)$. This is the stated entropy contraction.
\end{proof}

\subsection{Matrix analysis for the viscous velocity block}

\begin{definition}[Viscous matrices]\label{def:matrices}
For $X=(x_1,\dots,x_N)$ define
\[
K_{ij}(X):=
\begin{cases}
K(x_i,x_j), & i\neq j,\\
0, & i=j,
\end{cases}
\qquad
D(X):=\mathrm{diag}(d_1(X),\dots,d_N(X)),
\]
and
\[
W(X):=D(X)^{-1}K(X),
\qquad
L(X):=I-W(X).
\]
The velocity dynamics may then be written as
\[
dV
=
\bigl(-\nabla U(X)-A(X)V\bigr)\,dt+\sigma\,dW,
\qquad
A(X):=\gamma I_{Nd}+\nu_{\mathrm{visc}}\,L(X)\otimes I_d.
\]
\end{definition}

\begin{lemma}[Weighted symmetry and spectral bounds]\label{lem:weighted-sym}
For every configuration $X$:
\begin{enumerate}[label=\textnormal{(\roman*)},leftmargin=1.5em]
\item $D(X)W(X)=K(X)$ is symmetric;
\item the matrix
\[
\widetilde L(X):=D(X)^{1/2}L(X)D(X)^{-1/2}
=
D(X)^{-1/2}(D(X)-K(X))D(X)^{-1/2}
\]
is symmetric positive semidefinite;
\item every eigenvalue of $\widetilde L(X)$ lies in $[0,2]$;
\item $A(X)$ is similar to the symmetric positive definite matrix
\[
\widetilde A(X)
:=
\gamma I_{Nd}+\nu_{\mathrm{visc}}\widetilde L(X)\otimes I_d,
\]
hence every eigenvalue of $A(X)$ lies in $[\gamma,\gamma+2\nu_{\mathrm{visc}}]$.
\end{enumerate}
\end{lemma}

\begin{proof}
The identity $D(X)W(X)=K(X)$ is immediate from the definitions, and $K(X)$ is
symmetric because the kernel is symmetric. Therefore
\[
\widetilde L(X)
=
D(X)^{-1/2}(D(X)-K(X))D(X)^{-1/2}
\]
is symmetric. For any $z\in \R^N$,
\[
z^\top (D-K)z
=
\frac12\sum_{i\neq j}K_{ij}(z_i-z_j)^2\geq0,
\]
which proves positive semidefiniteness.

To bound the top of the spectrum, let $y\in\R^N$ and set
$z=D(X)^{-1/2}y$. Then
\[
y^\top \widetilde L(X)y
=
z^\top (D-K)z
=
\frac12\sum_{i\neq j}K_{ij}(z_i-z_j)^2.
\]
Also,
\[
y^\top y
=
\sum_i d_i z_i^2.
\]
Using $(z_i-z_j)^2\leq 2z_i^2+2z_j^2$ and the symmetry of $K$,
\[
y^\top \widetilde L(X)y
\leq
\sum_{i\neq j}K_{ij}(z_i^2+z_j^2)
=
2\sum_i d_i z_i^2
=
2\,y^\top y.
\]
Since $\widetilde L(X)$ is symmetric positive semidefinite, its Rayleigh
quotient belongs to $[0,2]$, and hence every eigenvalue lies in $[0,2]$.

Finally,
\[
A(X)
=
\bigl(D(X)^{-1/2}\otimes I_d\bigr)\widetilde A(X)
\bigl(D(X)^{1/2}\otimes I_d\bigr),
\]
so $A(X)$ and $\widetilde A(X)$ are similar and hence share the same spectrum.
Since $\widetilde L(X)$ is symmetric positive semidefinite with spectrum in
$[0,2]$, the spectrum of $\widetilde A(X)$ is contained in
$[\gamma,\gamma+2\nu_{\mathrm{visc}}]$.
\end{proof}

\begin{proposition}[Frozen-position Gaussian law and covariance bound]\label{prop:gaussian}
Fix a configuration $X$ and consider the frozen-position velocity process
\[
dV_t=\bigl(-b(X)-A(X)V_t\bigr)\,dt+\sigma\,dW_t,
\qquad
b(X):=(\nabla U(x_1),\dots,\nabla U(x_N)).
\]
Then:
\begin{enumerate}[label=\textnormal{(\roman*)},leftmargin=1.5em]
\item the stationary law is Gaussian
\[
\mathcal N\bigl(m_X,\Sigma_X\bigr),
\qquad
m_X=-A(X)^{-1}b(X);
\]
\item the covariance solves the Lyapunov equation
\begin{equation}\label{eq:lyap}
A(X)\Sigma_X+\Sigma_X A(X)^\top=\sigma^2 I_{Nd};
\end{equation}
\item its largest eigenvalue satisfies the $N$-uniform bound
\[
\lambda_{\max}(\Sigma_X)
\leq
\frac{k^*}{k_*}\frac{\sigma^2}{2\gamma}.
\]
\end{enumerate}
\end{proposition}

\begin{proof}
For fixed \(X\), the velocity equation is a linear Ornstein-Uhlenbeck process
with constant drift matrix \(A(X)\) and constant forcing \(-b(X)\). By
Lemma~\ref{lem:weighted-sym}, item $iv$, the spectrum of \(A(X)\) lies in
\([\gamma,\infty)\), so the linear system is exponentially stable. The unique
stationary law is therefore Gaussian with mean
\[
m_X=-A(X)^{-1}b(X),
\]
and the covariance \(\Sigma_X\) is the unique positive semidefinite solution of
\eqref{eq:lyap}.

To bound the covariance, write \(H(X):=D(X)^{1/2}\otimes I_d\) and
\[
A(X)=H(X)^{-1}\widetilde A(X)H(X),
\qquad
\widetilde A(X)=\gamma I_{Nd}+\nu_{\mathrm{visc}}\widetilde L(X)\otimes I_d.
\]
Since \(\widetilde A(X)\) is symmetric positive definite, the covariance admits
the standard OU representation
\[
\Sigma_X
=
\sigma^2\int_0^\infty e^{-A(X)t}e^{-A(X)^\top t}\,dt.
\]
Substituting the similarity transform gives
\[
\Sigma_X
=
\sigma^2 H^{-1}\left[\int_0^\infty
e^{-\widetilde A t}H^2e^{-\widetilde A t}\,dt\right]H^{-1}.
\]
For any \(u\in\R^{Nd}\) with \(\|u\|=1\), set \(w=H^{-1}u\). Then
\begin{align*}
u^\top \Sigma_X u
&=
\sigma^2\int_0^\infty
\langle H e^{-\widetilde A t}w,\,H e^{-\widetilde A t}w\rangle\,dt\\
&\leq
\sigma^2 \lambda_{\max}(H^2)
\int_0^\infty \|e^{-\widetilde A t}w\|^2\,dt\\
&\leq
\sigma^2 \lambda_{\max}(H^2)\int_0^\infty e^{-2\gamma t}\,dt\,\|w\|^2\\
&\leq
\frac{\sigma^2}{2\gamma}
\frac{\lambda_{\max}(H^2)}{\lambda_{\min}(H^2)}.
\end{align*}
Because
\[
(N-1)k_*\leq d_i(X)\leq (N-1)k^*
\qquad
\forall i,
\]
we have
\[
\frac{\lambda_{\max}(H^2)}{\lambda_{\min}(H^2)}
=
\frac{\max_i d_i(X)}{\min_i d_i(X)}
\leq
\frac{k^*}{k_*}.
\]
Taking the supremum over unit vectors \(u\) proves the stated bound on
\(\lambda_{\max}(\Sigma_X)\).
\end{proof}

\begin{corollary}[Conditional Gaussian Poincar\'e inequality]\label{cor:conditional-poincare}
For every frozen configuration $X$ and every smooth
$g:\R^{Nd}\to \R$,
\[
\mathrm{Var}_{\mathcal N(m_X,\Sigma_X)}(g)
\leq
\frac{k^*}{k_*}\frac{\sigma^2}{2\gamma}
\int |\nabla_v g|^2\,d\mathcal N(m_X,\Sigma_X).
\]
\end{corollary}

\begin{proof}
For a centered Gaussian law \(\mathcal N(m,\Sigma)\), the sharp Poincar\'e
constant is \(\lambda_{\max}(\Sigma)\). Since translating by the mean does not
change the variance or the gradient term, the same constant applies to
\(\mathcal N(m_X,\Sigma_X)\). Proposition~\ref{prop:gaussian} bounds
\(\lambda_{\max}(\Sigma_X)\) by \(\frac{k^*}{k_*}\frac{\sigma^2}{2\gamma}\), so
the stated inequality follows.
\end{proof}

\begin{lemma}[Derivative bounds for the normalized weights]\label{lem:weight-deriv}
Let
\[
L_K:=\sup_{x,y\in \R^d}\|\nabla_x K(x,y)\|<\infty.
\]
Then for every $N\geq2$ and every $i\neq j$,
\[
\|\nabla_{x_i}\omega_{ij}(X)\|
\leq
\frac{L_K}{(N-1)k_*}
+\frac{k^*L_K}{(N-1)k_*^2},
\]
and
\[
\|\nabla_{x_j}\omega_{ij}(X)\|
\leq
\frac{L_K}{(N-1)k_*}
+\frac{k^*L_K}{(N-1)^2k_*^2},
\]
and for $m\notin\{i,j\}$,
\[
\|\nabla_{x_m}\omega_{ij}(X)\|
\leq
\frac{k^*L_K}{(N-1)^2k_*^2}.
\]
In particular,
\[
\sup_{N\geq2}\sup_X \max_{i=1,\dots,N}\max_{m=1,\dots,N}
\sum_{j\neq i}\|\nabla_{x_m}\omega_{ij}(X)\|
<\infty.
\]
\end{lemma}

\begin{proof}
Differentiate
\[
\omega_{ij}(X)=\frac{K(x_i,x_j)}{d_i(X)},
\qquad
d_i(X)=\sum_{k\neq i}K(x_i,x_k).
\]
For the $x_i$ derivative,
\[
\nabla_{x_i}\omega_{ij}
=
\frac{\nabla_{x_i}K(x_i,x_j)}{d_i}
-\frac{K(x_i,x_j)}{d_i^2}
\sum_{k\neq i}\nabla_{x_i}K(x_i,x_k).
\]
Taking norms and using $d_i\geq (N-1)k_*$,
$K(x_i,x_j)\leq k^*$, and
$\|\nabla_{x_i}K\|\leq L_K$ yields
\[
\|\nabla_{x_i}\omega_{ij}\|
\leq
\frac{L_K}{(N-1)k_*}
+\frac{k^*(N-1)L_K}{(N-1)^2k_*^2}.
\]
This is the first estimate. For $m\notin\{i,j\}$ only the denominator depends on
$x_m$, so
\[
\nabla_{x_m}\omega_{ij}
=
-\frac{K(x_i,x_j)}{d_i^2}\nabla_{x_m}K(x_i,x_m),
\]
which gives the third estimate. For the $x_j$ derivative,
\[
\nabla_{x_j}\omega_{ij}
=
\frac{\nabla_{x_j}K(x_i,x_j)}{d_i}
-\frac{K(x_i,x_j)}{d_i^2}\nabla_{x_j}K(x_i,x_j),
\]
and the same lower bound on $d_i$ gives the second estimate. The uniform summability follows by summing the
displayed bounds over $j$ for fixed $(i,m)$.
\end{proof}

\begin{lemma}[Derivative bounds for the viscous force]\label{lem:force-deriv}
There exists a constant $C_{\partial F}<\infty$, depending only on
$(\nu_{\mathrm{visc}},k_*,k^*,L_K)$, such that for every configuration
$(X,V)\in(\R^{2d})^N$:
\begin{align}
\max_{1\leq i\leq N}\|\nabla_V F_i(X,V)\|_{\mathrm{op}}
&\leq
2\nu_{\mathrm{visc}},
\label{eq:force-v-deriv}\\
\sum_{m=1}^N \|\nabla_{x_m}F_i(X,V)\|^2
&\leq
C_{\partial F}
\left(
\|V_i\|^2+\frac1N\sum_{j=1}^N \|V_j\|^2
\right).
\label{eq:force-x-deriv}
\end{align}
\end{lemma}

\begin{proof}
The velocity derivatives are explicit:
\[
\nabla_{v_i}F_i=-\nu_{\mathrm{visc}}I_d,
\qquad
\nabla_{v_m}F_i=
\nu_{\mathrm{visc}}\omega_{im}(X)I_d
\quad(m\neq i).
\]
Hence
\[
\|\nabla_V F_i\|_{\mathrm{op}}
\leq
\nu_{\mathrm{visc}}
+\nu_{\mathrm{visc}}\sum_{m\neq i}\omega_{im}(X)
=
2\nu_{\mathrm{visc}},
\]
which proves \eqref{eq:force-v-deriv}.

For the $x$ derivatives,
\[
\nabla_{x_m}F_i(X,V)
=
\nu_{\mathrm{visc}}\sum_{j\neq i}\nabla_{x_m}\omega_{ij}(X)\,(V_j-V_i).
\]
Define
\[
A_{im}(X):=\sum_{j\neq i}\|\nabla_{x_m}\omega_{ij}(X)\|.
\]
By Lemma~\ref{lem:weight-deriv},
\[
\sup_{N\geq2}\sup_X \max_{i,m} A_{im}(X)\leq C_\omega
\]
for some finite constant $C_\omega$. Using Cauchy-Schwarz in the discrete sum,
\begin{align*}
\|\nabla_{x_m}F_i(X,V)\|^2
&\leq
\nu_{\mathrm{visc}}^2
\left(\sum_{j\neq i}\|\nabla_{x_m}\omega_{ij}\|\right)
\left(\sum_{j\neq i}\|\nabla_{x_m}\omega_{ij}\|\,\|V_j-V_i\|^2\right)\\
&\leq
\nu_{\mathrm{visc}}^2 C_\omega
\left(\sum_{j\neq i}\|\nabla_{x_m}\omega_{ij}\|\,(2\|V_j\|^2+2\|V_i\|^2)\right).
\end{align*}
Summing over $m$, and using the explicit derivative bounds from
Lemma~\ref{lem:weight-deriv}, we have for every $j\neq i$
\[
\sum_{m=1}^N\|\nabla_{x_m}\omega_{ij}(X)\|
\leq
\frac{C_\omega}{N},
\]
because only the terms $m\in\{i,j\}$ are of order $(N-1)^{-1}$ while the
remaining $(N-2)$ terms are of order $(N-1)^{-2}$.
Summing over $j\neq i$, this yields
\[
\sum_{m=1}^N\sum_{j\neq i}\|\nabla_{x_m}\omega_{ij}(X)\|
\leq
C_\omega.
\]
Then, by the same bound,
\[
\sum_{m=1}^N\sum_{j\neq i}\|\nabla_{x_m}\omega_{ij}(X)\|\,\|V_j\|^2
\leq
C_\omega \frac1N\sum_{j=1}^N \|V_j\|^2.
\]
Collecting terms yields \eqref{eq:force-x-deriv}.
\end{proof}

\section{Finite-Particle Hypocoercivity and N-Uniform LSI}

\subsection{Backbone gap and viscous perturbation}

For the kinetic backbone ($\nu_{\mathrm{visc}}=0$), the invariant law is
$\pi_1^{\otimes N}$ and the gap from
Corollary~\ref{cor:backbone-gap} tensorizes immediately.

\begin{proposition}[Tensorized backbone entropy gap]\label{prop:tensor}
Let $P_t^{N,0}$ denote the semigroup of the uncoupled $N$-particle kinetic
backbone obtained by setting $\nu_{\mathrm{visc}}=0$ in
Definition~\ref{def:finiteN}. Then
\[
\KL(\mu P_t^{N,0}\,\|\,\pi_1^{\otimes N})
\leq
e^{-2\lambda_{\mathrm{hypo}}t}\KL(\mu\,\|\,\pi_1^{\otimes N})
\]
for every $\mu\ll\pi_1^{\otimes N}$ and every $N\geq1$.
\end{proposition}

\begin{proof}
The uncoupled generator is the sum of the \(N\) one-particle generators, so the
semigroup factorizes as
\[
P_t^{N,0}=(P_t^{(1),0})^{\otimes N}.
\]
Let \(\mu\ll\pi_1^{\otimes N}\) and write
\[
h_0=\frac{d\mu}{d\pi_1^{\otimes N}}.
\]
Let \(h_t:=\frac{d(\mu P_t^{N,0})}{d\pi_1^{\otimes N}}\). The relative entropy
with respect to a product reference measure satisfies the usual conditional
chain rule:
\[
\KL(\mu P_t^{N,0}\,\|\,\pi_1^{\otimes N})
=
\sum_{i=1}^N
\E_{\mu P_t^{N,0}}\!\Bigl[
\KL\bigl((\mu P_t^{N,0})_{i\,|\,1:i-1}(\cdot\mid Z_1,\dots,Z_{i-1})\,\|\,\pi_1\bigr)
\Bigr].
\]
Because \(P_t^{N,0}\) acts independently on each coordinate and preserves
\(\pi_1\), the \(i\)-th conditional law evolves by the one-particle semigroup
\(P_t^{(1),0}\) while the earlier coordinates are frozen. Corollary~\ref{cor:backbone-gap}
therefore gives, for every conditioning value,
\[
\KL\bigl((\mu P_t^{N,0})_{i\,|\,1:i-1}(\cdot\mid Z_1,\dots,Z_{i-1})\,\|\,\pi_1\bigr)
\leq
 e^{-2\lambda_{\mathrm{hypo}}t}
\KL\bigl(\mu_{i\,|\,1:i-1}(\cdot\mid Z_1,\dots,Z_{i-1})\,\|\,\pi_1\bigr).
\]
Taking expectations and summing over \(i\) gives
\[
\KL(\mu P_t^{N,0}\,\|\,\pi_1^{\otimes N})
\leq
 e^{-2\lambda_{\mathrm{hypo}}t}\KL(\mu\,\|\,\pi_1^{\otimes N})
\]
for every $\mu\ll\pi_1^{\otimes N}$ and every $N\geq1$.
\end{proof}

\begin{remark}[Hypocoercive carré du champ]
For the backbone dynamics one may choose a Villani-type quadratic form
\[
\Gamma_{\mathrm{hypo}}^N(g,g)
:=
a\|\nabla_v g\|^2
+2b\,\nabla_x g\cdot\nabla_v g
+c\|\nabla_x g\|^2,
\qquad
ac>b^2,
\]
with coefficients depending only on $(U,\gamma,\sigma)$ and not on $N$, so that
$\Gamma_{\mathrm{hypo}}^N$ is equivalent to the full $H^1$ norm under the
backbone invariant law. The entropy gap of
Proposition~\ref{prop:tensor} is equivalent to an estimate of the form
\[
\Ent_{\pi_1^{\otimes N}}(g^2)
\leq
C_{\mathrm{hypo}}
\int \Gamma_{\mathrm{hypo}}^N(g,g)\,\dd \pi_1^{\otimes N}.
\]
We now explain why the viscous alignment does not destroy this structure; the
underlying modified-entropy mechanism is the one introduced in
\cite{villani2009hypocoercivity}.
\end{remark}

The fully interacting closure is most naturally formulated as a perturbative
statement around the tensorized backbone. For the mean-field goals of this
paper, it is enough to work in a quantitative small-viscosity regime, so the
finite-particle logarithmic Sobolev inequality is recorded below only under a
smallness condition on $\nu_{\mathrm{visc}}$.

\begin{definition}[Finite-particle entropy and Fisher functionals]\label{def:functionals-N}
Let $\Pi_N$ be a smooth probability measure on $(\R^{2d})^N$, and let
$h:(\R^{2d})^N\to(0,\infty)$ satisfy $\int h\,d\Pi_N=1$. Writing
$X=(x_1,\dots,x_N)$ and $V=(v_1,\dots,v_N)$, define
\begin{align*}
\mathcal A_N[h]
&:=
\int h\log h\,d\Pi_N,\\
\mathcal B_N[h]
&:=
T\int h\,|\nabla_V\log h|^2\,d\Pi_N,\\
\mathcal C_N[h]
&:=
T\int h\,\langle \nabla_X\log h,\nabla_V\log h\rangle\,d\Pi_N,\\
\mathcal D_N[h]
&:=
T\int h\,|\nabla_X\log h|^2\,d\Pi_N.
\end{align*}
Here $\nabla_X$ and $\nabla_V$ denote the full gradients in the $Nd$ position
and velocity variables.
\end{definition}

\begin{lemma}[Entropy production relative to an invariant law]\label{lem:entropy-prod-N}
Assume that $\pi_N$ is an invariant law for the finite-particle semigroup
$P_t^N$, and let
\[
h_t:=\frac{d(\mu P_t^N)}{d\pi_N}.
\]
Then
\[
\frac{d}{dt}\Ent_{\pi_N}(h_t)
=
-\gamma \mathcal B_N[h_t].
\]
\end{lemma}

\begin{proof}
Write the generator of the interacting diffusion in the form
\[
\mathcal L_N\phi
=
b_N\cdot \nabla_{X,V}\phi+\gamma T\Delta_V\phi,
\qquad
T=\frac{\sigma^2}{2\gamma},
\]
where $b_N$ is the full first-order drift. Since $f_t:=h_t\pi_N$ solves the
forward equation $\partial_t f_t=\mathcal L_N^*f_t$, we have
\[
\frac{d}{dt}\Ent_{\pi_N}(h_t)
=
\int (1+\log h_t)\,\mathcal L_N^*(h_t\pi_N)\,dX\,dV
=
\int h_t\,\mathcal L_N(\log h_t)\,d\pi_N.
\]
Because the second-order part of $\mathcal L_N$ acts only in $V$,
\[
\mathcal L_N(\log h_t)
=
\frac{\mathcal L_N h_t}{h_t}
-\gamma T|\nabla_V\log h_t|^2.
\]
Therefore
\[
\frac{d}{dt}\Ent_{\pi_N}(h_t)
=
\int \mathcal L_N h_t\,d\pi_N
-\gamma T\int h_t|\nabla_V\log h_t|^2\,d\pi_N.
\]
The first term vanishes by invariance of $\pi_N$:
\[
\int \mathcal L_N h_t\,d\pi_N=0.
\]
Using the definition of $\mathcal B_N$ gives the claim.
\end{proof}

\subsection{Quadratic perturbative closure in the small-viscosity regime}

The finite-particle LSI is proved by treating the viscous alignment as a small
first-order perturbation of the tensorized backbone dynamics. The right
perturbative object is the density ratio of the full invariant law \(\pi_N\)
with respect to the backbone product law \(\pi_{N,0}:=\pi_1^{\otimes N}\):
\[
q_N:=\frac{d\pi_N}{d\pi_{N,0}},
\qquad
\psi_N:=\log q_N.
\]
The aim of this subsection is to reduce the full interacting modified-entropy
estimate to a quadratic weighted stationary-corrector problem and a small set of
form bounds. This is the smallest perturbative package that still preserves the
spatially varying kernel structure used later in the mean-field and continuum
constructions.

\begin{definition}[Small-viscosity / large-friction regime]\label{def:small-visc-regime}
Throughout this subsection we keep the kinetic temperature
\[
T:=\frac{\sigma^2}{2\gamma}
\]
fixed, and we work in a regime
\[
\gamma\ge \gamma_0,
\qquad
0\le \nu_{\mathrm{visc}}\le \nu_*(\gamma),
\]
where \(\gamma_0>0\) and the threshold function \(\nu_*(\gamma)\) will be
chosen from the perturbative closure estimate. All constants below are allowed
to depend on \((T,\gamma,L_U,k_*,k^*,L_K)\), but they are always independent of
\(N\).
\end{definition}

\begin{definition}[Quadratic weighted perturbative spaces]\label{def:weighted-spaces-N}
For \(N\geq2\), define the quadratic weight
\[
\widetilde W_N(X,V)
:=
1+\frac1N\sum_{i=1}^N \bigl(\|x_i\|^2+\|v_i\|^2\bigr).
\]
For smooth \(f:(\R^{2d})^N\to\R\), set
\begin{align*}
\|f\|_{\mathcal X_N}^2
&:=
\int
\bigl(
f^2+\|\nabla_V f\|^2+\|\nabla_X f\|^2
\bigr)\,\widetilde W_N\,d\pi_{N,0},\\
\|f\|_{\mathcal Y_N}^2
&:=
\|f\|_{\mathcal X_N}^2
+
\int
\bigl(
\|\nabla_{VV}^2 f\|^2+\|\nabla_{XV}^2 f\|^2
\bigr)\,\widetilde W_N\,d\pi_{N,0}.
\end{align*}
\end{definition}

\begin{lemma}[Relative-generator decomposition and stationary corrector equation]\label{lem:relative-generator-decomp}
Define the backward first-order viscous perturbation on observables by
\[
\mathcal V_N f
:=
\sum_{i=1}^N
\left(
\sum_{j\neq i}\omega_{ij}(X)(v_j-v_i)
\right)\cdot \nabla_{v_i}f,
\]
so that
\[
\mathcal L_N=\mathcal L_{N,0}+\nu_{\mathrm{visc}}\mathcal V_N.
\]
Define also the forward perturbation relative to \(\pi_{N,0}\) by
\[
\mathscr V_N^\dagger f
:=
\pi_{N,0}^{-1}\mathcal V_N^*(\pi_{N,0}f).
\]
Let \(\pi_N\) be an invariant law of the full diffusion and let
\[
\mathscr K_N f
:=
\pi_N^{-1}\mathcal L_N^*(\pi_N f),
\qquad
\mathscr K_{N,0}f
:=
\pi_{N,0}^{-1}\mathcal L_{N,0}^*(\pi_{N,0}f).
\]
Then, for every smooth \(f\),
\[
\mathscr K_N f
=
\mathscr K_{N,0}f
+\nu_{\mathrm{visc}}\mathcal V_N f
+2\gamma T\,\nabla_V\psi_N\cdot \nabla_V f.
\]
Moreover, the stationary density ratio \(q_N=d\pi_N/d\pi_{N,0}\) satisfies
\[
\mathscr K_{N,0}q_N+\nu_{\mathrm{visc}}\mathscr V_N^\dagger q_N=0.
\]
\end{lemma}

\begin{proof}
Write the generator in the standard form
\[
\mathcal L_N\phi
=
b_N\cdot \nabla_{X,V}\phi+\gamma T\Delta_V\phi,
\]
where
\[
b_N(X,V)
=
b_{N,0}(X,V)+\nu_{\mathrm{visc}}\widehat F_N(X,V),
\]
and \(\widehat F_N\) is the vector field underlying \(\mathcal V_N\). For a
smooth positive density \(\mu\), the diffusion product rule gives
\[
\mu^{-1}\mathcal L_N^*(\mu f)
=
\mathcal L_N f+2\gamma T\,\nabla_V\log \mu\cdot \nabla_V f
+f\,\mu^{-1}\mathcal L_N^*\mu.
\]
Applying this with \(\mu=\pi_N\) and using \(\mathcal L_N^*\pi_N=0\) yields
\[
\mathscr K_N f
=
\mathcal L_N f+2\gamma T\,\nabla_V\log \pi_N\cdot \nabla_V f.
\]
Applying the same identity with \(\mu=\pi_{N,0}\) and
\(\mathcal L_{N,0}^*\pi_{N,0}=0\) gives
\[
\mathscr K_{N,0} f
=
\mathcal L_{N,0} f+2\gamma T\,\nabla_V\log \pi_{N,0}\cdot \nabla_V f.
\]
Subtracting the two displays and using
\(\pi_N=q_N\pi_{N,0}\), hence
\(\nabla_V\log \pi_N-\nabla_V\log \pi_{N,0}=\nabla_V\psi_N\), proves
\[
\mathscr K_N f
=
\mathscr K_{N,0}f+\nu_{\mathrm{visc}}\mathcal V_N f
+2\gamma T\,\nabla_V\psi_N\cdot \nabla_V f.
\]
For the stationary corrector equation, start from
\[
0=\mathcal L_N^*(q_N\pi_{N,0})
=\mathcal L_{N,0}^*(q_N\pi_{N,0})
+\nu_{\mathrm{visc}}\mathcal V_N^*(q_N\pi_{N,0}).
\]
Dividing by \(\pi_{N,0}\) gives
\[
\mathscr K_{N,0}q_N+\nu_{\mathrm{visc}}\pi_{N,0}^{-1}\mathcal V_N^*(\pi_{N,0}q_N)=0,
\]
which is exactly
\[
\mathscr K_{N,0}q_N+\nu_{\mathrm{visc}}\mathscr V_N^\dagger q_N=0.
\]
\end{proof}

\begin{theorem}[Uniform quadratic weighted backbone resolvent]\label{lem:weighted-backbone-resolvent}
There exists a constant \(C_{\mathrm{res}}<\infty\), independent of \(N\), such
that for every \(g\in \mathcal X_N\) with \(\int g\,d\pi_{N,0}=0\), the Poisson
equation
\[
-\mathscr K_{N,0}\phi=g,
\qquad
\int \phi\,d\pi_{N,0}=0,
\]
admits a unique solution \(\phi\in \mathcal Y_N\) satisfying
\[
\|\phi\|_{\mathcal Y_N}
\leq
C_{\mathrm{res}}\|g\|_{\mathcal X_N}.
\]
Moreover, one may take \(C_{\mathrm{res}}\) uniform for all
\(\gamma\ge \gamma_0\).
\end{theorem}

\begin{proof}
Write the one-particle backbone generator in relative form as
\[
\mathscr K_{1,0}f
:=
\pi_1^{-1}(\mathcal L_{1,0})^*(\pi_1 f),
\]
and let \(S_t:=P_t^{(1),0}\) be the one-particle semigroup from
Corollary~\ref{cor:backbone-gap}. Introduce the one-particle quadratic weight
\[
w(x,v):=1+\|x\|^2+\|v\|^2,
\]
and let \(\mathcal X_1\) and \(\mathcal Y_1\) denote the obvious one-particle
versions of the norms from Definition~\ref{def:weighted-spaces-N}.

\smallskip
\noindent
\textbf{Step 1: one-particle linearized hypocoercive decay.}
Let \(u_t:=S_t f\), where \(f\in \mathcal X_1\) is centered. For
\(\varepsilon>0\) small, the perturbation
\[
h_t^{(\varepsilon)}:=1+\varepsilon u_t
\]
solves the same relative equation \(\partial_t h=L_0 h\) as in
Definition~\ref{def:functionals}. Expanding the functionals of
Definition~\ref{def:functionals} at \(h=1\) gives
\begin{align*}
\mathcal A[1+\varepsilon u]
&=
\frac{\varepsilon^2}{2}\|u\|_{L^2(\pi_1)}^2+O(\varepsilon^3),\\
\mathcal B[1+\varepsilon u]
&=
\varepsilon^2 T\|\nabla_v u\|_{L^2(\pi_1)}^2+O(\varepsilon^3),\\
\mathcal C[1+\varepsilon u]
&=
\varepsilon^2 T\langle \nabla_x u,\nabla_v u\rangle_{L^2(\pi_1)}+O(\varepsilon^3),\\
\mathcal D[1+\varepsilon u]
&=
\varepsilon^2 T\|\nabla_x u\|_{L^2(\pi_1)}^2+O(\varepsilon^3).
\end{align*}
Dividing the differential inequalities of Proposition~\ref{prop:abcd} and
Theorem~\ref{thm:backbone-detailed} by \(\varepsilon^2\) and letting
\(\varepsilon\downarrow0\) yields the linearized modified energy
\[
\mathscr E[u]
:=
\frac12\|u\|_{L^2(\pi_1)}^2
+aT\|\nabla_v u\|_{L^2(\pi_1)}^2
+2bT\langle \nabla_x u,\nabla_v u\rangle_{L^2(\pi_1)}
+cT\|\nabla_x u\|_{L^2(\pi_1)}^2,
\]
with the same coefficients \(a,b,c\) as in
Theorem~\ref{thm:backbone-detailed}, and constants
\(C_{\mathrm{lin}},\eta_{\mathrm{lin}}>0\), depending only on
\((T,\gamma,L_U,m_U)\), such that
\begin{equation}\label{eq:linearized-backbone-decay}
\mathscr E[u_t]
\le
e^{-\eta_{\mathrm{lin}}t}\mathscr E[f],
\qquad
\|u_t\|_{L^2(\pi_1)}^2+\|\nabla_x u_t\|_{L^2(\pi_1)}^2+\|\nabla_v u_t\|_{L^2(\pi_1)}^2
\le
C_{\mathrm{lin}}e^{-\eta_{\mathrm{lin}}t}\|f\|_{\mathcal X_1}^2.
\end{equation}

\smallskip
\noindent
\textbf{Step 2: one-particle quadratic Lyapunov weight.}
Choose \(\theta>0\) small and define
\[
W_\theta(x,v):=1+\|x\|^2+2\theta\,x\cdot v+\theta^2\|v\|^2.
\]
For \(\theta\) sufficiently small, \(W_\theta\) is equivalent to \(w\):
\[
c_\theta w\le W_\theta\le C_\theta w.
\]
Let
\[
G_0\phi
:=
v\cdot \nabla_x\phi
+(-\nabla U(x)-\gamma v)\cdot \nabla_v\phi
+\gamma T\Delta_v\phi
\]
be the one-particle observable generator dual to \(\mathscr K_{1,0}\) in the
weighted integration-by-parts identity used below. A direct calculation gives
\[
G_0 W_\theta
=
2x\cdot v
+2\theta\|v\|^2
+2\theta x\cdot(-\nabla U-\gamma v)
+2\theta v\cdot(-\nabla U-\gamma v)
+2\gamma T\theta^2 d.
\]
Using Assumption~\ref{ass:potential}, item \(iii\), together with Young's
inequality and the smallness of \(\theta\), this implies
\begin{equation}\label{eq:backbone-weight-drift}
G_0 W_\theta
\le
-c_\theta W_\theta+C_\theta
\end{equation}
for constants \(c_\theta,C_\theta>0\) depending only on
\((T,\gamma,\alpha_U,\beta_U)\).

\smallskip
\noindent
\textbf{Step 3: one-particle weighted Poisson estimate.}
Define
\[
R_1 f:=\int_0^\infty S_t f\,dt.
\]
Because \(f\) is centered and \eqref{eq:linearized-backbone-decay} gives
exponential decay in \(L^2(\pi_1)\), the integral converges in the unweighted
\(H^1\)-norm and satisfies
\[
-\mathscr K_{1,0}R_1 f=f.
\]
Set \(\phi:=R_1 f\). Testing the Poisson equation against \(W_\theta \phi\)
gives
\[
\int W_\theta \phi f\,d\pi_1
=
-\int W_\theta \phi\,\mathscr K_{1,0}\phi\,d\pi_1.
\]
We now compute the right-hand side exactly. For the symmetric part,
\begin{align*}
-\int W_\theta \phi\,L_{\mathrm{sym}}\phi\,d\pi_1
&=
\gamma T\int \nabla_v(W_\theta\phi)\cdot \nabla_v\phi\,d\pi_1\\
&=
\gamma T\int W_\theta\|\nabla_v\phi\|^2\,d\pi_1
+
\gamma T\int \phi\,\nabla_vW_\theta\cdot \nabla_v\phi\,d\pi_1\\
&=
\gamma T\int W_\theta\|\nabla_v\phi\|^2\,d\pi_1
-\frac12\int \phi^2\,L_{\mathrm{sym}}W_\theta\,d\pi_1.
\end{align*}
For the anti-symmetric part,
\[
-\int W_\theta \phi\,L_{\mathrm{anti}}\phi\,d\pi_1
=
\frac12\int \phi^2\,L_{\mathrm{anti}}W_\theta\,d\pi_1,
\]
because \(L_{\mathrm{anti}}\) is anti-symmetric in \(L^2(\pi_1)\). Since
\[
-L_{\mathrm{sym}}W_\theta+L_{\mathrm{anti}}W_\theta=-G_0W_\theta,
\]
the two contributions combine into the exact identity
\begin{equation}\label{eq:weighted-poisson-identity}
\int W_\theta \phi f\,d\pi_1
=
\gamma T\int W_\theta\|\nabla_v\phi\|^2\,d\pi_1
-\frac12\int (G_0W_\theta)\phi^2\,d\pi_1.
\end{equation}
Invoking \eqref{eq:backbone-weight-drift} in
\eqref{eq:weighted-poisson-identity} yields
\[
\gamma T\int W_\theta\|\nabla_v\phi\|^2\,d\pi_1
+\frac{c_\theta}{2}\int W_\theta\phi^2\,d\pi_1
\le
\int W_\theta |\phi|\,|f|\,d\pi_1
+\frac{C_\theta}{2}\|\phi\|_{L^2(\pi_1)}^2.
\]
Using Young's inequality and then integrating
\eqref{eq:linearized-backbone-decay} over \([0,\infty)\), we obtain
\begin{equation}\label{eq:unweighted-resolvent-h1}
\|\phi\|_{L^2(\pi_1)}
+\|\nabla_x\phi\|_{L^2(\pi_1)}
+\|\nabla_v\phi\|_{L^2(\pi_1)}
\le
C\|f\|_{\mathcal X_1}
\end{equation}
and therefore
\begin{equation}\label{eq:weighted-phi-v}
\int W_\theta\bigl(\phi^2+\|\nabla_v\phi\|^2\bigr)\,d\pi_1
\le
C\|f\|_{\mathcal X_1}^2.
\end{equation}
Next, differentiating the Poisson equation and using the commutators from
Proposition~\ref{prop:comm} yields
\begin{align}
-\mathscr K_{1,0}(\nabla_v\phi)+\gamma \nabla_v\phi+\nabla_x\phi
&=
\nabla_v f,
\label{eq:resolvent-v-eq}\\
-\mathscr K_{1,0}(\nabla_x\phi)-(\nabla^2U)\nabla_v\phi
&=
\nabla_x f.
\label{eq:resolvent-x-eq}
\end{align}
For \(g_v:=\nabla_v\phi\), test \eqref{eq:resolvent-v-eq} componentwise against
\(W_\theta g_v\). Repeating the integration-by-parts argument from
\eqref{eq:weighted-poisson-identity} gives
\begin{align}
\gamma T\int W_\theta\|\nabla_v g_v\|^2\,d\pi_1
+\gamma\int W_\theta\|g_v\|^2\,d\pi_1
-\frac12\int (G_0W_\theta)\|g_v\|^2\,d\pi_1
\nonumber\\
\le
\int W_\theta\|\nabla_v f\|\,\|g_v\|\,d\pi_1
+\int W_\theta\|g_x\|\,\|g_v\|\,d\pi_1,
\label{eq:weighted-gv}
\end{align}
where \(g_x:=\nabla_x\phi\). Using \eqref{eq:backbone-weight-drift},
\eqref{eq:unweighted-resolvent-h1}, and Young's inequality, we obtain
\begin{equation}\label{eq:weighted-gv-bound}
\int W_\theta\bigl(\|g_v\|^2+\|\nabla_v g_v\|^2\bigr)\,d\pi_1
\le
C\|f\|_{\mathcal X_1}^2
+C\int W_\theta\|g_x\|^2\,d\pi_1.
\end{equation}
For \(g_x:=\nabla_x\phi\), test \eqref{eq:resolvent-x-eq} against
\(W_\theta g_x\). The same computation yields
\begin{align}
\gamma T\int W_\theta\|\nabla_v g_x\|^2\,d\pi_1
-\frac12\int (G_0W_\theta)\|g_x\|^2\,d\pi_1
\nonumber\\
\le
\int W_\theta\|\nabla_x f\|\,\|g_x\|\,d\pi_1
+\int W_\theta\|(\nabla^2U)g_v\|\,\|g_x\|\,d\pi_1.
\label{eq:weighted-gx}
\end{align}
Since \(\|\nabla^2U\|_{\mathrm{op}}\le L_U\), another Young inequality gives
\begin{equation}\label{eq:weighted-gx-bound}
\int W_\theta\bigl(\|g_x\|^2+\|\nabla_v g_x\|^2\bigr)\,d\pi_1
\le
C\|f\|_{\mathcal X_1}^2
+C\int W_\theta\|g_v\|^2\,d\pi_1.
\end{equation}
Combining \eqref{eq:weighted-gv-bound} and \eqref{eq:weighted-gx-bound} and
absorbing the lower-order terms into the left-hand side yields
\begin{equation}\label{eq:weighted-first-second}
\int W_\theta
\bigl(
\|\nabla_v\phi\|^2+\|\nabla_x\phi\|^2
+\|\nabla_{VV}^2\phi\|^2+\|\nabla_{XV}^2\phi\|^2
\bigr)\,d\pi_1
\le
C\|f\|_{\mathcal X_1}^2.
\end{equation}
Combining the weighted \(L^2\), first-derivative, and \(VV/XV\)-Hessian bounds,
namely \eqref{eq:weighted-phi-v} and \eqref{eq:weighted-first-second}, and
using the equivalence \(W_\theta\asymp w\), proves
\begin{equation}\label{eq:one-particle-resolvent}
\|R_1 f\|_{\mathcal Y_1}
\le
C_1\|f\|_{\mathcal X_1}.
\end{equation}
Since the parameters remain in a fixed regime \(\gamma\ge\gamma_0\), the
constant \(C_1\) is uniform there.

\smallskip
\noindent
\textbf{Step 4: tensorized semigroup representation.}
For the uncoupled \(N\)-particle backbone,
\[
\mathscr K_{N,0}=\sum_{i=1}^N \mathscr K_{1,0}^{(i)},
\qquad
P_t^{N,0}=S_t^{\otimes N},
\]
where \(\mathscr K_{1,0}^{(i)}\) acts on the \(i\)-th coordinate only. Let
\[
\phi:=\int_0^\infty P_t^{N,0}g\,dt.
\]
Because \(g\) is centered and Proposition~\ref{prop:tensor} gives exponential
decay of the tensorized backbone semigroup to equilibrium, the integral
converges in \(L^2(\pi_{N,0})\), \(\phi\) is centered, and the standard
semigroup identity yields
\[
-\mathscr K_{N,0}\phi=g.
\]

\smallskip
\noindent
\textbf{Step 5: tensorization of the weighted estimate.}
The product structure of \(\pi_{N,0}=\pi_1^{\otimes N}\) and the averaged weight
\[
\widetilde W_N(X,V)=1+\frac1N\sum_{i=1}^N w(x_i,v_i)
\]
allow the one-particle estimate \eqref{eq:one-particle-resolvent} to be applied
coordinatewise. First derivatives in the \(i\)-th coordinate commute with the
semigroup in all other coordinates, and the same is true for the \(VV\) and
\(XV\) second derivatives appearing in \(\mathcal Y_N\). Therefore, by Fubini,
Minkowski, and the fact that the semigroup factorizes,
\[
\|\phi\|_{\mathcal Y_N}
\le
C_1\|g\|_{\mathcal X_N},
\]
with the same constant \(C_1\) as in the one-particle estimate. The averaging
by \(1/N\) in \(\widetilde W_N\) is exactly what prevents any growth in \(N\).

\smallskip
\noindent
\textbf{Step 6: uniqueness.}
If \(\phi_1\) and \(\phi_2\) are centered solutions, then
\(\psi:=\phi_1-\phi_2\) satisfies \(\mathscr K_{N,0}\psi=0\). Hence
\(P_t^{N,0}\psi=\psi\) for all \(t\ge0\). By ergodicity of the tensorized
backbone semigroup, the only invariant \(L^2(\pi_{N,0})\) functions are
constants, so the centering condition forces \(\psi=0\).

Thus \(\phi\) is the unique centered solution and the resolvent estimate holds
with \(C_{\mathrm{res}}:=C_1\), uniformly in \(N\).
\end{proof}

\begin{lemma}[Boundedness of the viscous adjoint perturbation]\label{lem:viscous-adjoint-bounded}
Define
\[
\mathscr V_N^\dagger f
:=
\pi_{N,0}^{-1}\mathcal V_N^*(\pi_{N,0}f).
\]
There exists \(C_{\mathrm V}<\infty\), independent of \(N\), such that
\[
\|\mathscr V_N^\dagger f\|_{\mathcal X_N}
\leq
C_{\mathrm V}\|f\|_{\mathcal Y_N}
\qquad
\forall f\in \mathcal Y_N.
\]
\end{lemma}

\begin{proof}
Writing \(\mathscr V_N^\dagger\) out explicitly gives first-order terms
involving \(F_i\), \(\nabla_{v_i}\cdot F_i\), and the Gaussian score
\(\nabla_{v_i}\log \pi_{N,0}=-v_i/T\). The force bounds from
Lemma~\ref{lem:visc-bound}, Lemma~\ref{lem:visc-diss}, and
Lemma~\ref{lem:force-deriv} show that each coefficient is uniformly bounded in
\(N\) by the structural constants and at most linear in \(V\), while the
quadratic weight \(\widetilde W_N\) absorbs the resulting growth. This yields
the stated
\(\mathcal Y_N\to\mathcal X_N\) bound.
\end{proof}

\begin{proposition}[Stationary corrector fixed point]\label{prop:stationary-corrector}
Let \(q_N=d\pi_N/d\pi_{N,0}\). Then \(q_N\) satisfies
\[
q_N-1
=
-\nu_{\mathrm{visc}}
\mathscr K_{N,0}^{-1}\Pi_0\!\left(\mathscr V_N^\dagger q_N\right),
\qquad
\Pi_0 g:=g-\int g\,d\pi_{N,0}.
\]
If
\[
\nu_{\mathrm{visc}}C_{\mathrm{res}}C_{\mathrm V}<1,
\]
then this equation has a unique solution in \(\mathcal Y_N\), it satisfies
\[
\int q_N\,d\pi_{N,0}=1,
\qquad
q_N>0
\quad
\pi_{N,0}\text{-a.e.},
\]
and
\[
\|q_N-1\|_{\mathcal Y_N}
\leq
\frac{\nu_{\mathrm{visc}}C_{\mathrm{res}}C_{\mathrm V}}
{1-\nu_{\mathrm{visc}}C_{\mathrm{res}}C_{\mathrm V}}.
\]
\end{proposition}

\begin{proof}
The stationary equation
\[
\mathscr K_{N,0}q_N+\nu_{\mathrm{visc}}\mathscr V_N^\dagger q_N=0
\]
is the second assertion of Lemma~\ref{lem:relative-generator-decomp}. Subtract
the \(\pi_{N,0}\)-mean, invert \(\mathscr K_{N,0}\) on centered functions by
Lemma~\ref{lem:weighted-backbone-resolvent}, and apply
Lemma~\ref{lem:viscous-adjoint-bounded}. The smallness condition makes the
fixed-point map a contraction on \(\mathcal Y_N\). The normalization and
positivity are inherited from the definition \(q_N=d\pi_N/d\pi_{N,0}\).
\end{proof}

\begin{lemma}[Quadratic entropy-Lyapunov control]\label{lem:quad-entropy-lyapunov}
There exists \(C_W<\infty\), independent of \(N\), such that for every smooth
\(h>0\) with \(\int h\,d\pi_N=1\),
\[
\int \widetilde W_N\,h\,d\pi_N
\leq
C_W\Bigl(1+\mathcal A_N[h]+\mathcal B_N[h]+\mathcal D_N[h]\Bigr).
\]
\end{lemma}

\begin{proof}
Let \(V_N\) be the finite-particle Lyapunov function from
Theorem~\ref{thm:finite-harris}, so that
\begin{equation}\label{eq:quad-lyap-equiv}
V_N(X,V)\asymp \widetilde W_N(X,V)
\end{equation}
and
\begin{equation}\label{eq:quad-lyap-drift}
\mathcal L_N V_N\le -\lambda_V V_N+C_V
\end{equation}
with constants independent of \(N\). By the explicit quadratic cross-term
construction recalled in the proof of Theorem~\ref{thm:finite-harris}, one may
choose \(V_N\) in the form
\[
V_N(X,V)
=
1+\eta_1\frac1N U_N(X)
+\eta_2\frac1N\|V\|^2
+\eta_3\frac1N\langle X,V\rangle,
\]
hence
\begin{equation}\label{eq:quad-lyap-grad}
\|\nabla_V V_N(X,V)\|^2\le C_\nabla V_N(X,V)
\end{equation}
for some \(C_\nabla<\infty\), independent of \(N\).

Fix \(\eta>0\) so small that
\[
\gamma T\,\eta C_\nabla\le \frac{\lambda_V}{2}.
\]
Let \(\vartheta_R\in C^\infty([0,\infty))\) be increasing and concave, with
\[
\vartheta_R(s)=s\quad (0\le s\le R),
\qquad
0\le \vartheta_R'(s)\le 1,
\qquad
\vartheta_R''(s)\le 0,
\]
and \(\vartheta_R\) constant on \([R+1,\infty)\). Set
\[
\Phi_R:=e^{\eta \vartheta_R(V_N)}.
\]
Since \(\Phi_R\) is smooth and bounded, invariance of \(\pi_N\) gives
\[
\int \mathcal L_N\Phi_R\,d\pi_N=0.
\]
Using the chain rule and the fact that diffusion acts only in velocity,
\[
\mathcal L_N\Phi_R
=
\eta \vartheta_R'(V_N)\Phi_R\,\mathcal L_N V_N
+\gamma T\Phi_R\Bigl(
\eta \vartheta_R''(V_N)+\eta^2(\vartheta_R'(V_N))^2
\Bigr)\|\nabla_V V_N\|^2.
\]
Because \(\vartheta_R''\le0\) and \(0\le \vartheta_R'\le1\), combining
\eqref{eq:quad-lyap-drift} and \eqref{eq:quad-lyap-grad} yields
\[
\mathcal L_N\Phi_R
\le
\Phi_R\Bigl(
-\eta\lambda_V V_N+\eta C_V+\gamma T\eta^2 C_\nabla V_N
\Bigr)
\le
\Phi_R\Bigl(
-\frac{\eta\lambda_V}{2}V_N+\eta C_V
\Bigr).
\]
Since the function
\[
s\mapsto e^{\eta s}\Bigl(-\frac{\eta\lambda_V}{2}s+\eta C_V\Bigr)
\]
tends to \(-\infty\) as \(s\to\infty\), there exists \(C_\eta<\infty\) such
that
\[
e^{\eta s}\Bigl(-\frac{\eta\lambda_V}{2}s+\eta C_V\Bigr)
\le
-\frac{\eta\lambda_V}{4}s\,e^{\eta s}+C_\eta
\qquad
\forall s\ge0.
\]
Using \(\vartheta_R(s)\le s\), we therefore obtain
\[
\mathcal L_N\Phi_R
\le
-\frac{\eta\lambda_V}{4}V_N\,\Phi_R+C_\eta.
\]
Integrating against \(\pi_N\) gives
\[
\frac{\eta\lambda_V}{4}\int V_N\,\Phi_R\,d\pi_N
\le
C_\eta.
\]
Letting \(R\to\infty\) and using monotone convergence yields
\begin{equation}\label{eq:exp-lyap-moment}
\int V_N e^{\eta V_N}\,d\pi_N<\infty.
\end{equation}
In particular,
\[
\int e^{\eta V_N}\,d\pi_N<\infty.
\]
By \eqref{eq:quad-lyap-equiv}, after possibly decreasing \(\eta\), there exists
\(\eta_*>0\) and \(C_*>0\), independent of \(N\), such that
\begin{equation}\label{eq:exp-weight-moment}
\int e^{\eta_*\widetilde W_N}\,d\pi_N\le C_*.
\end{equation}

Now apply the entropy variational inequality to the probability measure
\(h\,d\pi_N\) and the observable \(\eta_*\widetilde W_N\):
\[
\int \eta_*\widetilde W_N\,h\,d\pi_N
\le
\mathcal A_N[h]+\log\int e^{\eta_*\widetilde W_N}\,d\pi_N.
\]
Using \eqref{eq:exp-weight-moment} gives
\[
\int \widetilde W_N\,h\,d\pi_N
\le
\eta_*^{-1}\mathcal A_N[h]+\eta_*^{-1}\log C_*.
\]
Since \(\mathcal B_N[h]\) and \(\mathcal D_N[h]\) are non-negative, this is
stronger than the stated estimate.
\end{proof}

\begin{lemma}[Velocity-score form bound]\label{lem:psi-v-form}
There exists \(C_{\mathrm{sc}}<\infty\), independent of \(N\), such that for
every smooth \(h>0\) with \(\int h\,d\pi_N=1\),
\[
\int h\,\|\nabla_V\psi_N\|^2\,d\pi_N
\leq
C_{\mathrm{sc}}\nu_{\mathrm{visc}}^2
\Bigl(1+\mathcal A_N[h]+\mathcal B_N[h]+\mathcal D_N[h]\Bigr).
\]
\end{lemma}

\begin{lemma}[Velocity-Hessian form bound]\label{lem:psi-vv-form}
There exists \(C_{\mathrm{vv}}<\infty\), independent of \(N\), such that for
every smooth \(h>0\) with \(\int h\,d\pi_N=1\),
\[
\left|
\int h\,
\nabla_{VV}^2\psi_N[\nabla_V\log h,\nabla_V\log h]\,d\pi_N
\right|
\leq
C_{\mathrm{vv}}\nu_{\mathrm{visc}}\,\mathcal B_N[h].
\]
\end{lemma}

\begin{lemma}[Mixed Hessian form bound]\label{lem:psi-xv-form}
There exists \(C_{\mathrm{xv}}<\infty\), independent of \(N\), such that for
every smooth \(h>0\) with \(\int h\,d\pi_N=1\),
\[
\left|
\int h\,
\nabla_{XV}^2\psi_N[\nabla_X\log h,\nabla_V\log h]\,d\pi_N
\right|
\leq
C_{\mathrm{xv}}\nu_{\mathrm{visc}}
\bigl(\mathcal B_N[h]+\mathcal D_N[h]\bigr).
\]
\end{lemma}

\begin{proposition}[Explicit viscous commutator remainder bounds]\label{prop:viscous-commutator-remainders}
Let \(\mathfrak R_B^{\mathrm{visc}}[h]\),
\(\mathfrak R_C^{\mathrm{visc}}[h]\), and
\(\mathfrak R_D^{\mathrm{visc}}[h]\) denote the remainders generated by the
explicit perturbation \(\nu_{\mathrm{visc}}\mathcal V_N\) when differentiating
\(\mathcal B_N[h]\), \(\mathcal C_N[h]\), and \(\mathcal D_N[h]\). Then there
exists \(C_{\mathrm{visc}}<\infty\), independent of \(N\), such that
\[
|\mathfrak R_B^{\mathrm{visc}}[h]|
+
|\mathfrak R_C^{\mathrm{visc}}[h]|
+
|\mathfrak R_D^{\mathrm{visc}}[h]|
\leq
C_{\mathrm{visc}}\nu_{\mathrm{visc}}
\bigl(\mathcal B_N[h]+\mathcal D_N[h]\bigr).
\]
\end{proposition}

\begin{proof}
Differentiate the explicit perturbation \(\nu_{\mathrm{visc}}\mathcal V_N\)
through the definitions of \(\mathcal B_N\), \(\mathcal C_N\), and
\(\mathcal D_N\). The coefficients that appear are the velocity and position
derivatives of the finite-particle viscous force. Lemma~\ref{lem:force-deriv}
gives the uniform bounds
\[
\|\nabla_V F_i\|\lesssim \nu_{\mathrm{visc}},
\qquad
\sum_{m=1}^N \|\nabla_{x_m}F_i\|^2
\lesssim
\nu_{\mathrm{visc}}^2
\Bigl(|v_i|^2+\frac1N\sum_{j=1}^N |v_j|^2\Bigr),
\]
and the quadratic moments are controlled by the quadratic Lyapunov structure
already available for the invariant law. Cauchy-Schwarz and Young therefore
yield the stated form bound.
\end{proof}

\begin{proposition}[Stationary-score remainder bounds]\label{prop:score-commutator-remainders}
Let \(\mathfrak R_B^{\mathrm{sc}}[h]\),
\(\mathfrak R_C^{\mathrm{sc}}[h]\), and
\(\mathfrak R_D^{\mathrm{sc}}[h]\) denote the remainders generated by the extra
drift \(2\gamma T\,\nabla_V\psi_N\cdot\nabla_V\) in the relative generator.
Then there exists \(C_{\mathrm{sc,form}}<\infty\), independent of \(N\), such
that
\[
|\mathfrak R_B^{\mathrm{sc}}[h]|
+
|\mathfrak R_C^{\mathrm{sc}}[h]|
+
|\mathfrak R_D^{\mathrm{sc}}[h]|
\leq
C_{\mathrm{sc,form}}\nu_{\mathrm{visc}}
\bigl(\mathcal B_N[h]+\mathcal D_N[h]\bigr).
\]
\end{proposition}

\begin{proof}
The score drift contributes only terms containing
\(\nabla_V\psi_N\), \(\nabla_{VV}^2\psi_N\), or \(\nabla_{XV}^2\psi_N\) paired
with \(\nabla_V\log h\) and \(\nabla_X\log h\). The zeroth-order pieces are
controlled by Lemma~\ref{lem:psi-v-form} together with
Lemma~\ref{lem:quad-entropy-lyapunov}; the quadratic \(VV\) and mixed \(XV\)
pieces are controlled by Lemma~\ref{lem:psi-vv-form} and
Lemma~\ref{lem:psi-xv-form}. Summing the resulting estimates gives the claim.
\end{proof}

\begin{proposition}[Relative commutator ODEs]\label{prop:relative-commutator-odes}
Let \(h_t=d(\mu P_t^N)/d\pi_N\). Then
\begin{align*}
\frac{d}{dt}\mathcal A_N[h_t]
&= 
-\gamma \mathcal B_N[h_t],\\
\frac{d}{dt}\mathcal B_N[h_t]
&\leq
-2\gamma \mathcal B_N[h_t]-2\mathcal C_N[h_t]
+\mathfrak R_B^{\mathrm{visc}}[h_t]+\mathfrak R_B^{\mathrm{sc}}[h_t],\\
\frac{d}{dt}\mathcal C_N[h_t]
&\leq
-\mathcal D_N[h_t]-\gamma \mathcal C_N[h_t]
+L_U\sqrt{\mathcal B_N[h_t]\mathcal D_N[h_t]}
+\mathfrak R_C^{\mathrm{visc}}[h_t]+\mathfrak R_C^{\mathrm{sc}}[h_t],\\
\frac{d}{dt}\mathcal D_N[h_t]
&\leq
\mathfrak R_D^{\mathrm{visc}}[h_t]+\mathfrak R_D^{\mathrm{sc}}[h_t].
\end{align*}
\end{proposition}

\begin{proof}
Lemma~\ref{lem:relative-generator-decomp} decomposes the relative generator into
the backbone part, the explicit viscous perturbation, and the additional score
drift \(2\gamma T\,\nabla_V\psi_N\cdot\nabla_V\). The backbone contribution is
exactly the one from Section~4 and yields
\[
\frac{d}{dt}\mathcal A_N[h_t]=-\gamma \mathcal B_N[h_t],
\]
together with the coercive terms
\(-2\gamma\mathcal B_N[h_t]-2\mathcal C_N[h_t]\) in the
\(\mathcal B_N\)-equation and
\(-\mathcal D_N[h_t]-\gamma\mathcal C_N[h_t]
+L_U\sqrt{\mathcal B_N[h_t]\mathcal D_N[h_t]}\) in the
\(\mathcal C_N\)-equation. The explicit perturbation
\(\nu_{\mathrm{visc}}\mathcal V_N\) contributes the remainders handled in
Proposition~\ref{prop:viscous-commutator-remainders}, while the score drift
contributes the remainders handled in
Proposition~\ref{prop:score-commutator-remainders}. Summing those pieces gives
the stated system.
\end{proof}

\begin{theorem}[Small-viscosity relative hypocoercive closure]\label{thm:small-visc-relative-closure}
There exist constants \(\gamma_0>0\), \(C_{\mathrm{pert}}<\infty\),
\(\nu_*>0\), and \(\lambda_*>0\), depending only on
\((T,L_U,k_*,k^*,L_K)\) and independent of \(N\), such that for every
\(\gamma\ge\gamma_0\), every \(N\geq2\), every
\(0\leq \nu_{\mathrm{visc}}\leq \nu_*\), and every
\[
h_t=\frac{d(\mu P_t^N)}{d\pi_N},
\]
one has
\[
\frac{d}{dt}\Phi_N[h_t]
\leq
-\bigl(\lambda_{\mathrm{back}}-C_{\mathrm{pert}}\nu_{\mathrm{visc}}\bigr)
\bigl(\mathcal B_N[h_t]+\mathcal D_N[h_t]\bigr).
\]
In particular, for \(0\leq \nu_{\mathrm{visc}}\leq \nu_*\),
\[
\frac{d}{dt}\Phi_N[h_t]
\leq
-\lambda_*\bigl(\mathcal B_N[h_t]+\mathcal D_N[h_t]\bigr).
\]
\end{theorem}

\begin{proof}
Insert the differential inequalities from
Proposition~\ref{prop:relative-commutator-odes} into the derivative of the
modified entropy
\[
\Phi_N[h]
=
\mathcal A_N[h]+a\mathcal B_N[h]+2b\mathcal C_N[h]+c\mathcal D_N[h],
\]
with the same coefficients \(a,b,c\) as in
Theorem~\ref{thm:backbone-detailed}. The backbone terms reproduce the matrix
coercivity argument from Section~4 and yield
\(-\lambda_{\mathrm{back}}(\mathcal B_N+\mathcal D_N)\), while
Proposition~\ref{prop:viscous-commutator-remainders} and
Proposition~\ref{prop:score-commutator-remainders} contribute only
\(+C_{\mathrm{pert}}\nu_{\mathrm{visc}}(\mathcal B_N+\mathcal D_N)\) after
renaming the sum of constants. Choosing
\[
\nu_*:=\frac{\lambda_{\mathrm{back}}}{2C_{\mathrm{pert}}},
\qquad
\lambda_*:=\frac{\lambda_{\mathrm{back}}}{2},
\]
gives the conclusion.
\end{proof}

\begin{remark}
The point of the quadratic package is that it asks only for form bounds needed
by the modified entropy. It does \emph{not} require \(N\)-uniform
\(L^\infty\)-control of \(\nabla\psi_N\) in the full \(Nd\)-dimensional space,
which is substantially stronger than what the closure argument actually uses.
\(\widetilde W_N\) is chosen to match the quadratic Lyapunov control already
available from the Harris theorem, so the perturbative closure and the
stationary compactness arguments now work at the same moment level.
\end{remark}

\begin{theorem}[Small-viscosity $N$-uniform hypocoercive LSI for the kinetic core]\label{thm:finite-lsi}
There exist thresholds \(\gamma_0>0\) and \(\nu_*>0\), depending only on
\((T,L_U,k_*,k^*,\ell_{\mathrm{visc}})\), such that the following holds. If
\[
\gamma\ge \gamma_0,
\]
and if
\[
0\leq \nu_{\mathrm{visc}}\leq \nu_*,
\]
and if the full interacting diffusion
\eqref{eq:finite-x}--\eqref{eq:finite-v} admits a unique invariant law $\pi_N$,
then there exist constants $\lambda_*>0$ and $C_{\mathrm{LSI}}<\infty$,
depending only on the backbone parameters, the kernel bounds, and the chosen
threshold $\nu_*$, but \emph{independent of $N$} and uniform for
$\nu_{\mathrm{visc}}\in[0,\nu_*]$, such that for every $N\geq2$:
\begin{align}
\Ent_{\pi_N}(g^2)
&\leq
C_{\mathrm{LSI}}
\int \Gamma_{\mathrm{hypo}}^N(g,g)\,\dd\pi_N,
\label{eq:lsi-finite}\\
\KL(\mu P_t^N\,\|\,\pi_N)
&\leq
e^{-2\lambda_* t}\KL(\mu\,\|\,\pi_N)
\qquad
\forall \mu\ll\pi_N.
\label{eq:kl-finite}
\end{align}
\end{theorem}

\begin{proof}[Perturbative reduction]
We isolate the backbone estimate and the single additional perturbative input
required for the interacting system.

\smallskip
\noindent
\textbf{Step 1: backbone hypocoercivity.}
For $\nu_{\mathrm{visc}}=0$, Proposition~\ref{prop:tensor} yields the tensorized
entropy gap. The more explicit weighted computation is recorded in
Theorem~\ref{thm:backbone-detailed}: the transport commutator is controlled by
the Hessian bound $M=L_U$, and the modified entropy functional
\[
\Phi_N[h]
=
\mathcal A_N[h]+a\mathcal B_N[h]+2b\mathcal C_N[h]+c\mathcal D_N[h]
\]
decays with a rate depending only on $(\gamma,L_U,\sigma)$ and not on $N$.
Denote this backbone rate by $\lambda_{\mathrm{back}}>0$.

\smallskip
\noindent
\textbf{Step 2: perturbative closure estimate.}
Write the full generator as
\[
\mathcal L_N=\mathcal L_{N,0}+\nu_{\mathrm{visc}}\mathcal V_N,
\]
where $\mathcal L_{N,0}$ is the tensorized backbone generator and
$\mathcal V_N$ is the row-normalized viscous alignment drift. The small-viscosity
route reduces the full theorem to the estimate that there exists a constant
$C_{\mathrm{pert}}<\infty$, depending only on
$(\gamma,\sigma,L_U,k_*,k^*,\ell_{\mathrm{visc}})$ and independent of $N$, such
that for
\[
h_t=\frac{d(\mu P_t^N)}{d\pi_N}
\]
the modified entropy obeys
\begin{equation}\label{eq:small-visc-closure}
\frac{d}{dt}\Phi_N[h_t]
\leq
-\bigl(\lambda_{\mathrm{back}}-C_{\mathrm{pert}}\nu_{\mathrm{visc}}\bigr)
\bigl(\mathcal B_N[h_t]+\mathcal D_N[h_t]\bigr).
\end{equation}
This is exactly the perturbative form of the statement that the viscous force is
a small first-order correction to the backbone commutator calculus, and
Theorem~\ref{thm:small-visc-relative-closure} proves this estimate in the
quadratic perturbative regime.

\smallskip
\noindent
\textbf{Step 3: choice of the viscosity threshold.}
Choose
\[
\nu_*:=\frac{\lambda_{\mathrm{back}}}{2C_{\mathrm{pert}}}.
\]
Then for every $\nu_{\mathrm{visc}}\in[0,\nu_*]$ one has
\[
\lambda_{\mathrm{back}}-C_{\mathrm{pert}}\nu_{\mathrm{visc}}
\geq
\frac{\lambda_{\mathrm{back}}}{2}
=:\lambda_*,
\]
so \eqref{eq:small-visc-closure} yields
\[
\frac{d}{dt}\Phi_N[h_t]
\leq
-\lambda_*\bigl(\mathcal B_N[h_t]+\mathcal D_N[h_t]\bigr).
\]

\smallskip
\noindent
\textbf{Step 4: from modified entropy to logarithmic Sobolev.}
Once the previous display is available, the rest of the proof is unchanged. Let
$h_t=d(\mu P_t^N)/d\pi_N$. The entropy production identity is
\[
\frac{d}{dt}\Ent_{\pi_N}(h_t)
=
-\gamma \mathcal B_N[h_t].
\]
by Lemma~\ref{lem:entropy-prod-N}. Since the quadratic form in
$(\mathcal B_N,\mathcal C_N,\mathcal D_N)$ is
positive definite, there exists $\alpha_0>0$, depending only on $(a,b,c)$, such
that
\[
\Phi_N[h]\geq \Ent_{\pi_N}(h)+\alpha_0\bigl(\mathcal B_N[h]+\mathcal D_N[h]\bigr).
\]
Integrating the previous two displays over time from $0$ to $\infty$ gives
\[
\Ent_{\pi_N}(h_0)
=
\gamma\int_0^\infty \mathcal B_N[h_t]\,dt
\leq
\frac{\gamma}{\alpha_0\lambda_*}\Phi_N[h_0].
\]
It remains to control $\Phi_N[h_0]$ by the hypocoercive carré du champ. Write
$h_0=g^2/\int g^2\,d\pi_N$. Then
\[
\mathcal B_N[h_0]
=
4T\int |\nabla_v g|^2\,d\pi_N,
\qquad
\mathcal D_N[h_0]
=
4T\int |\nabla_x g|^2\,d\pi_N,
\]
and
\[
|\mathcal C_N[h_0]|
\leq
4T\left(\int |\nabla_x g|^2\,d\pi_N\right)^{1/2}
\left(\int |\nabla_v g|^2\,d\pi_N\right)^{1/2}.
\]
The conditional Gaussian Poincar\'e inequality from
Corollary~\ref{cor:conditional-poincare} gives the microscopic coercivity in
velocity, while the confining potential and the transport commutator provide the
macroscopic coercivity in position. Combining these estimates with the previous
display yields
\[
\Phi_N[h_0]
\leq
C_0\int \Gamma_{\mathrm{hypo}}^N(g,g)\,d\pi_N
\]
for an $N$-independent constant $C_0$. Therefore
\[
\Ent_{\pi_N}(g^2)
\leq
\frac{\gamma C_0}{\alpha_0\lambda_*}
\int \Gamma_{\mathrm{hypo}}^N(g,g)\,d\pi_N.
\]
This is \eqref{eq:lsi-finite} with
$C_{\mathrm{LSI}}=\gamma C_0/(\alpha_0\lambda_*)$. The exponential entropy
decay \eqref{eq:kl-finite} is the time-dependent counterpart of the same
interacting modified-entropy estimate: one runs the differential inequality for
$\Phi_N[h_t]$ along the semigroup together with
Lemma~\ref{lem:entropy-prod-N}, exactly as in the backbone argument.
\end{proof}

\begin{remark}
Under the smooth coercive hypotheses stated here, the existence and uniqueness
of the invariant law $\pi_N$ assumed in Theorem~\ref{thm:finite-lsi} is
internalized later in Theorem~\ref{thm:finite-harris} by the same
Lyapunov/minorization/Harris package used in Volume~3. The genuinely new
technical input for the small-viscosity route is the perturbative estimate
\eqref{eq:small-visc-closure}; all subsequent LSI, KL, and mean-field
consequences are formal corollaries of that bound.
\end{remark}

\begin{corollary}[Exponential convergence in total variation]\label{cor:tv-finite}
Under the assumptions of Theorem~\ref{thm:finite-lsi},
\[
\|\mu P_t^N-\pi_N\|_{\TV}
\leq
\sqrt{2}\,e^{-\lambda_* t}\,
\KL(\mu\,\|\,\pi_N)^{1/2}
\qquad
\forall \mu\ll\pi_N.
\]
\end{corollary}

\begin{proof}
Pinsker's inequality says that for any probability measures \(\alpha,\beta\),
\[
\|\alpha-\beta\|_{\TV}\le \sqrt{2\,\KL(\alpha\|\beta)}.
\]
Apply this with \(\alpha=\mu P_t^N\) and \(\beta=\pi_N\), then use
\eqref{eq:kl-finite}:
\[
\|\mu P_t^N-\pi_N\|_{\TV}
\le
\sqrt{2\,\KL(\mu P_t^N\|\pi_N)}
\le
\sqrt{2}\,e^{-\lambda_* t}\,\KL(\mu\|\pi_N)^{1/2}.
\]
\end{proof}

\subsection{The squashed observation chain}

\begin{proposition}[The squashing step cannot enlarge probabilistic distances]\label{prop:squash-transfer}
Let $Q_\tau^N$ be the time-$\tau$ observation operator obtained by evolving the
continuous diffusion for time $\tau$ and then applying the deterministic map
$S(X,V)=(X,\psi(V))$ componentwise. Then for every pair of probability measures
$\mu,\nu$ on $(\R^{2d})^N$,
\[
\|S_\#\mu-S_\#\nu\|_{\TV}\leq \|\mu-\nu\|_{\TV}.
\]
If $W_p$ denotes any Wasserstein distance with $1\leq p<\infty$, then
\[
W_p(S_\#\mu,S_\#\nu)\leq W_p(\mu,\nu).
\]
\end{proposition}

\begin{proof}
For total variation, if \(A\) is measurable in the target space then
\[
(S_\#\mu)(A)-(S_\#\nu)(A)=\mu(S^{-1}A)-\nu(S^{-1}A),
\]
so taking the supremum over measurable sets immediately yields
\(\|S_\#\mu-S_\#\nu\|_{\TV}\le \|\mu-\nu\|_{\TV}\).

For Wasserstein distances, let \(\pi\) be any coupling of \(\mu\) and \(\nu\).
Then \((S,S)_\#\pi\) is a coupling of \(S_\#\mu\) and \(S_\#\nu\), and
Lemma~\ref{lem:squash} gives
\[
\|(X,\psi(V))-(X',\psi(V'))\|
\le
\|(X,V)-(X',V')\|.
\]
Integrating this inequality against \(\pi\) and taking the infimum over
couplings proves the Wasserstein contraction.
\end{proof}

\begin{remark}
Proposition~\ref{prop:squash-transfer} is the precise reason one can carry out
the convergence proofs at the level of the unsquashed continuous SDE and only
afterwards impose the algorithmic speed cap. The cap improves boundedness of the
observed chain and never worsens the entropy, total-variation, or Wasserstein
estimates.
\end{remark}

\section{Entropy and Total-Variation Convergence}

The finite-particle entropy estimate controls the continuous interacting kinetic
core. This section records the two consequences needed later: convergence of the
observed chain and the existence of an $N$-independent relaxation rate.

\begin{corollary}[Observed kinetic chain]\label{cor:observed}
Let $Q_\tau^N$ be the squashed observation chain from
Proposition~\ref{prop:squash-transfer}. If the continuous kinetic core satisfies
\eqref{eq:kl-finite}, then for every initial law $\mu\ll\pi_N$,
\[
\|\mu Q_\tau^{N,m}-S_\#\pi_N\|_{\TV}
\leq
\sqrt{2}\,e^{-\lambda_* m\tau}\,
\KL(\mu\,\|\,\pi_N)^{1/2}
\qquad
\forall m\in\mathbb{N}.
\]
At all observation times the velocities also obey the deterministic bound
$\|v_i\|\leq V_{\mathrm{alg}}$.
\end{corollary}

\begin{proof}
Let \(P_{m\tau}^N\) denote the continuous-time evolution over \(m\tau\). The
observation chain is obtained by composing that evolution with the deterministic
map \(S\), and both the Markov kernel \(P_{m\tau}^N\) and the pushforward \(S_\#\)
are contractions in total variation. Therefore
\[
\|\mu Q_\tau^{N,m}-S_\#\pi_N\|_{\TV}
\le
\|\mu P_{m\tau}^N-\pi_N\|_{\TV}.
\]
Now apply Corollary~\ref{cor:tv-finite} at time \(m\tau\) to get the displayed
bound. The velocity estimate follows directly from
\(\|\psi(v)\|\le V_{\mathrm{alg}}\), so every observed velocity is capped at the
algorithmic bound.
\end{proof}

\section{Mean-Field Equation and Propagation of Chaos}

\subsection{The McKean-Vlasov limit}

We now compare the finite-particle law to the nonlinear mean-field equation
\eqref{eq:mf-pde}. The key fact is that the interaction field from
Definition~\ref{def:mean-field} is globally Lipschitz in both the state
variables and the underlying measure, thanks to the strictly positive kernel
floor.

\begin{lemma}[Empirical-force defect]\label{lem:force-defect}
Let
\[
\mu_t^N:=\frac1N\sum_{j=1}^N \delta_{(X_j(t),V_j(t))}
\]
be the empirical measure of the particle system and, for each $i$, let
\[
F_i^N(t):=
\nu_{\mathrm{visc}}\sum_{j\neq i}\omega_{ij}(X(t))(V_j(t)-V_i(t))
\]
denote the exact finite-particle alignment force. Then for every finite horizon
$T>0$ there exists $C_T<\infty$, independent of $N$, such that
\[
\sup_{t\in[0,T]}
\frac1N\sum_{i=1}^N
\E\Bigl[\bigl\|
F_i^N(t)-F[\mu_t^N](X_i(t),V_i(t))
\bigr\|^2\Bigr]
\leq
\frac{C_T}{N}.
\]
\end{lemma}

\begin{proof}
Fix $i$ and abbreviate $X_i=X_i(t)$, $V_i=V_i(t)$. By definition,
\[
F[\mu_t^N](X_i,V_i)
=
\frac{\nu_{\mathrm{visc}}}{m_{\mu_t^N}(X_i)}
\frac1N\sum_{j=1}^N K(X_i,X_j)(V_j-V_i),
\]
where
\[
m_{\mu_t^N}(X_i)
=
\frac1N\sum_{j=1}^N K(X_i,X_j).
\]
Since $K(X_i,X_i)=1+\kappa_0$, the finite-particle force and the empirical-field
force differ only because of the self-interaction term and the normalization:
\begin{align*}
F_i^N-F[\mu_t^N](X_i,V_i)
&=
\nu_{\mathrm{visc}}
\left[
\sum_{j\neq i}\frac{K(X_i,X_j)}{\sum_{k\neq i}K(X_i,X_k)}(V_j-V_i)
-\frac{\sum_{j\neq i}K(X_i,X_j)(V_j-V_i)}{\sum_{k=1}^N K(X_i,X_k)}
\right].
\end{align*}
Write
\[
S_i:=\sum_{j\neq i}K(X_i,X_j)(V_j-V_i),
\qquad
d_i:=\sum_{k\neq i}K(X_i,X_k).
\]
Then
\[
F_i^N-F[\mu_t^N](X_i,V_i)
=
\nu_{\mathrm{visc}} S_i\left(\frac{1}{d_i}-\frac{1}{d_i+K(X_i,X_i)}\right)
=
\nu_{\mathrm{visc}} S_i\frac{K(X_i,X_i)}{d_i(d_i+K(X_i,X_i))}.
\]
By Assumption~\ref{ass:kernel},
\[
d_i\geq (N-1)k_*,
\qquad
K(X_i,X_i)\leq k^*,
\]
and therefore
\[
\left|\frac{K(X_i,X_i)}{d_i(d_i+K(X_i,X_i))}\right|
\leq
\frac{k^*}{(N-1)^2k_*^2}.
\]
On the other hand,
\[
\|S_i\|
\leq
\sum_{j\neq i}K(X_i,X_j)\|V_j-V_i\|
\leq
2k^*\sum_{j\neq i}(\|V_j\|+\|V_i\|).
\]
Using the uniform second-moment bounds on the particle system over finite time
intervals, which follow from Proposition~\ref{prop:moments},
we obtain
\[
\E\|S_i\|^2\leq C_T N^2
\]
with $C_T$ independent of $N$. Combining the last three displays yields
\[
\E\|F_i^N-F[\mu_t^N](X_i,V_i)\|^2
\leq
\frac{C_T}{N^2}.
\]
Averaging over $i$ proves the claim.
\end{proof}

\begin{lemma}[Monte Carlo fluctuation of the mean-field force]\label{lem:mc-force}
Let $(\bar X_i(t),\bar V_i(t))_{i=1}^N$ be i.i.d.\ copies of the nonlinear
McKean-Vlasov process with common law $f_t$, and let
\[
\bar\mu_t^N:=\frac1N\sum_{j=1}^N \delta_{(\bar X_j(t),\bar V_j(t))}
\]
be the associated empirical measure. Then for every finite horizon $T>0$ there
exists $C_T<\infty$, independent of $N$, such that
\[
\sup_{t\in[0,T]}
\frac1N\sum_{i=1}^N
\E\Bigl[
\bigl\|F[\bar\mu_t^N](\bar X_i(t),\bar V_i(t))
-F[f_t](\bar X_i(t),\bar V_i(t))\bigr\|^2
\Bigr]
\leq
\frac{C_T}{N}.
\]
\end{lemma}

\begin{proof}
By exchangeability it suffices to estimate the index $i=1$. Fix $t\in[0,T]$
and condition on $(\bar X_1,\bar V_1)=(x,v)$. Define
\[
N_N(x,v):=\frac1N\sum_{j=1}^N K(x,\bar X_j)(\bar V_j-v),
\qquad
D_N(x):=\frac1N\sum_{j=1}^N K(x,\bar X_j),
\]
and also
\[
n_t(x,v):=\iint K(x,y)(q-v)\,f_t(dy,dq),
\qquad
m_t(x):=\int K(x,y)\,\rho_{f_t}(dy).
\]
Then
\[
F[\bar\mu_t^N](x,v)=\nu_{\mathrm{visc}}\frac{N_N(x,v)}{D_N(x)},
\qquad
F[f_t](x,v)=\nu_{\mathrm{visc}}\frac{n_t(x,v)}{m_t(x)}.
\]
Because $D_N(x)\geq k_*$ and $m_t(x)\geq k_*$,
\[
\left\|
\frac{N_N}{D_N}-\frac{n_t}{m_t}
\right\|
\leq
\frac{\|N_N-n_t\|}{k_*}
+\frac{\|n_t\|}{k_*^2}|D_N-m_t|.
\]
Conditionally on $(\bar X_1,\bar V_1)=(x,v)$, the variables
\[
Y_j:=K(x,\bar X_j)(\bar V_j-v),
\qquad
Z_j:=K(x,\bar X_j)
\]
are independent for $j\geq2$. Their conditional second moments satisfy
\[
\E[\|Y_j\|^2\mid \bar X_1=x,\bar V_1=v]
\leq
2(k^*)^2\bigl(\E\|\bar V_j\|^2+\|v\|^2\bigr),
\qquad
\E[|Z_j|^2\mid \bar X_1=x,\bar V_1=v]\leq (k^*)^2.
\]
The $j=1$ contributions to $N_N$ and $D_N$ are deterministic once $(x,v)$ is
fixed and are of size $O(N^{-1}(1+\|v\|))$, so they can be absorbed into the
same $N^{-1}$ bounds.
Therefore the empirical averages have conditional variances
\[
\E\bigl[\|N_N(x,v)-n_t(x,v)\|^2\mid \bar X_1=x,\bar V_1=v\bigr]
\leq
\frac{C_T(1+\|v\|^2)}{N},
\]
and
\[
\E\bigl[|D_N(x)-m_t(x)|^2\mid \bar X_1=x,\bar V_1=v\bigr]
\leq
\frac{C_T}{N}.
\]
Moreover,
\[
\|n_t(x,v)\|
\leq
k^*\int (\|q\|+\|v\|)\,f_t(dy,dq)
\leq
C_T(1+\|v\|).
\]
Combining the last three displays gives
\[
\E\Bigl[
\|F[\bar\mu_t^N](\bar X_1,\bar V_1)-F[f_t](\bar X_1,\bar V_1)\|^2
\mid \bar X_1,\bar V_1
\Bigr]
\leq
\frac{C_T(1+\|\bar V_1\|^2)}{N}.
\]
Integrating over $(\bar X_1,\bar V_1)$ and using
Proposition~\ref{prop:moments} yields the claim.
\end{proof}

\begin{theorem}[Strong propagation of chaos]\label{thm:poc-strong}
Let $(\bar X_i,\bar V_i)_{i=1}^N$ be i.i.d.\ solutions of the nonlinear
McKean-Vlasov SDE with initial law $f_0$, driven by the same Brownian motions
as the particle system \eqref{eq:finite-x}--\eqref{eq:finite-v}, and started
from the same initial data:
\[
(X_i(0),V_i(0))=(\bar X_i(0),\bar V_i(0))
\qquad
\text{a.s.\ for every }i.
\]
Set
\[
\varepsilon_N(t):=
\frac1N\sum_{i=1}^N
\E\Bigl[\|X_i(t)-\bar X_i(t)\|^2+\|V_i(t)-\bar V_i(t)\|^2\Bigr].
\]
Then for every finite horizon $T>0$ there exists $C_T<\infty$,
independent of $N$, such that
\[
\sup_{t\in[0,T]}\varepsilon_N(t)\leq \frac{C_T}{N}.
\]
\end{theorem}

\begin{proof}
Write
\[
e_i^x:=X_i-\bar X_i,
\qquad
e_i^v:=V_i-\bar V_i,
\qquad
\Delta_i:=F_i^N(X,V)-F[f_t](\bar X_i,\bar V_i).
\]
Then
\[
de_i^x=e_i^v\,dt,
\qquad
de_i^v=
\bigl(-\nabla U(X_i)+\nabla U(\bar X_i)-\gamma e_i^v+\Delta_i\bigr)\,dt.
\]
It\^o's formula yields
\[
\frac{d}{dt}\E\|e_i^x\|^2
=
2\E\langle e_i^x,e_i^v\rangle
\leq
\E\|e_i^x\|^2+\E\|e_i^v\|^2,
\]
and, by the Lipschitz bound on $\nabla U$,
\begin{align*}
\frac{d}{dt}\E\|e_i^v\|^2
&=
2\E\langle e_i^v,-\nabla U(X_i)+\nabla U(\bar X_i)-\gamma e_i^v+\Delta_i\rangle\\
&\leq
(L_U+1)\E\|e_i^x\|^2
+(L_U+1-2\gamma)\E\|e_i^v\|^2
+\E\|\Delta_i\|^2.
\end{align*}
Adding the two inequalities and averaging over $i$ gives
\begin{equation}\label{eq:eps-pre}
\frac{d}{dt}\varepsilon_N(t)
\leq
C_0\varepsilon_N(t)
+\frac1N\sum_{i=1}^N \E\|\Delta_i(t)\|^2
\end{equation}
for a constant $C_0$ depending only on $(L_U,\gamma)$.

Decompose
\[
\Delta_i=R_i+S_i+T_i,
\]
where
\begin{align*}
R_i&:=F_i^N(X,V)-F[\mu_t^N](X_i,V_i),\\
S_i&:=F[\mu_t^N](X_i,V_i)-F[\bar\mu_t^N](\bar X_i,\bar V_i),\\
T_i&:=F[\bar\mu_t^N](\bar X_i,\bar V_i)-F[f_t](\bar X_i,\bar V_i).
\end{align*}
Using $(a+b+c)^2\leq 3(a^2+b^2+c^2)$,
\[
\frac1N\sum_{i=1}^N \E\|\Delta_i\|^2
\leq
3\left(
\frac1N\sum_i \E\|R_i\|^2
+\frac1N\sum_i \E\|S_i\|^2
+\frac1N\sum_i \E\|T_i\|^2
\right).
\]
Lemma~\ref{lem:force-defect} yields
\[
\frac1N\sum_i \E\|R_i\|^2\leq \frac{C_T}{N},
\]
and Lemma~\ref{lem:mc-force} yields
\[
\frac1N\sum_i \E\|T_i\|^2\leq \frac{C_T}{N}.
\]
For the middle term, Lemma~\ref{lem:mf-lipschitz} gives
\[
\|S_i\|
\leq
L_{\mathrm{mf}}
\Bigl(\|e_i^x\|+\|e_i^v\|+W_1(\mu_t^N,\bar\mu_t^N)\Bigr).
\]
Coupling the empirical measures particlewise yields
\[
W_1(\mu_t^N,\bar\mu_t^N)
\leq
\frac1N\sum_{j=1}^N \bigl(\|e_j^x\|+\|e_j^v\|\bigr),
\]
and Jensen's inequality therefore gives
\[
\E W_1(\mu_t^N,\bar\mu_t^N)^2\leq 2\varepsilon_N(t).
\]
Consequently
\[
\frac1N\sum_{i=1}^N \E\|S_i\|^2
\leq
C_1\varepsilon_N(t).
\]
Substituting into \eqref{eq:eps-pre} yields
\[
\frac{d}{dt}\varepsilon_N(t)
\leq
C_2\varepsilon_N(t)+\frac{C_T}{N}.
\]
Because $\varepsilon_N(0)=0$, Gr\"onwall's lemma gives
\[
\sup_{t\in[0,T]}\varepsilon_N(t)\leq \frac{C_T}{N}.
\]
\end{proof}

\begin{theorem}[Relative entropy on path space and at fixed time]\label{thm:poc}
Under the hypotheses of Theorem~\ref{thm:poc-strong},
let $P_{[0,T]}^N$ be the law on path space
$C([0,T];(\R^{2d})^N)$ of the interacting system
\eqref{eq:finite-x}--\eqref{eq:finite-v}, and let $Q_{[0,T]}^N$ be the law of
the independent nonlinear system from Theorem~\ref{thm:poc-strong}. Then for
every finite horizon $T>0$ there exists $C_T<\infty$, independent of $N$, such
that
\begin{equation}\label{eq:path-kl}
\KL\bigl(P_{[0,T]}^N\,\|\,Q_{[0,T]}^N\bigr)\leq C_T.
\end{equation}
In particular, for every $t\in[0,T]$,
\begin{equation}\label{eq:hn}
\KL\bigl(\rho_t^N\,\|\,f_t^{\otimes N}\bigr)\leq C_T.
\end{equation}
\end{theorem}

\begin{proof}
Under $Q_{[0,T]}^N$, the processes satisfy
\[
d\bar V_i=
\bigl(-\nabla U(\bar X_i)-\gamma \bar V_i+F[f_t](\bar X_i,\bar V_i)\bigr)\,dt
+\sigma\,dW_i.
\]
The interacting system differs only in the velocity drift. Since the diffusion
matrix is constant in the velocity variables, Girsanov's theorem applies. The
Radon-Nikod\'ym derivative of $P_{[0,T]}^N$ with respect to $Q_{[0,T]}^N$ is
\[
\frac{dP_{[0,T]}^N}{dQ_{[0,T]}^N}
=
\exp\!\left(
\frac{1}{\sigma}\sum_{i=1}^N \int_0^T \Delta_i(s)\cdot dW_i(s)
-\frac{1}{2\sigma^2}\sum_{i=1}^N \int_0^T \|\Delta_i(s)\|^2\,ds
\right),
\]
with
\[
\Delta_i(s):=F_i^N(X_s,V_s)-F[f_s](X_i(s),V_i(s)).
\]
The moment bounds from Proposition~\ref{prop:moments}, together with
Lemma~\ref{lem:force-defect} and Theorem~\ref{thm:poc-strong}, imply Novikov's
criterion on every finite interval. Taking expectations under $P_{[0,T]}^N$
gives the exact entropy identity
\[
\KL\bigl(P_{[0,T]}^N\,\|\,Q_{[0,T]}^N\bigr)
=
\frac{1}{2\sigma^2}
\sum_{i=1}^N \E \int_0^T \|\Delta_i(s)\|^2\,ds.
\]
We estimate the drift discrepancy by inserting the coupled empirical measure
$\bar\mu_t^N$ from Theorem~\ref{thm:poc-strong}. Write
\[
\Delta_i=R_i+U_i+T_i+V_i,
\]
where
\begin{align*}
R_i&:=F_i^N(X,V)-F[\mu_t^N](X_i,V_i),\\
U_i&:=F[\mu_t^N](X_i,V_i)-F[\bar\mu_t^N](\bar X_i,\bar V_i),\\
T_i&:=F[\bar\mu_t^N](\bar X_i,\bar V_i)-F[f_t](\bar X_i,\bar V_i),\\
V_i&:=F[f_t](\bar X_i,\bar V_i)-F[f_t](X_i,V_i).
\end{align*}
Then $(a+b+c+d)^2\leq 4(a^2+b^2+c^2+d^2)$ gives
\[
\frac1N\sum_{i=1}^N \E\|\Delta_i(s)\|^2
\leq
\frac{4}{N}\sum_{i=1}^N \E\|R_i(s)\|^2
+\frac{4}{N}\sum_{i=1}^N \E\|U_i(s)\|^2
+\frac{4}{N}\sum_{i=1}^N \E\|T_i(s)\|^2
+\frac{4}{N}\sum_{i=1}^N \E\|V_i(s)\|^2.
\]
Lemma~\ref{lem:force-defect} controls the $R_i$ term, and
Lemma~\ref{lem:mc-force} controls the $T_i$ term:
\[
\frac1N\sum_i \E\|R_i(s)\|^2
\leq \frac{C_T}{N},
\qquad
\frac1N\sum_i \E\|T_i(s)\|^2
\leq \frac{C_T}{N}.
\]
For the $U_i$ term, Lemma~\ref{lem:mf-lipschitz} gives
\[
\|U_i\|
\leq
L_{\mathrm{mf}}
\Bigl(\|X_i-\bar X_i\|+\|V_i-\bar V_i\|+W_1(\mu_t^N,\bar\mu_t^N)\Bigr).
\]
Exactly as in the proof of Theorem~\ref{thm:poc-strong},
\[
\E W_1(\mu_t^N,\bar\mu_t^N)^2\leq 2\varepsilon_N(t),
\]
and therefore
\[
\frac1N\sum_i \E\|U_i(s)\|^2\leq C_T\varepsilon_N(s).
\]
For the last term, the Lipschitz bound of Lemma~\ref{lem:mf-lipschitz} gives
\[
\frac1N\sum_i \E\|V_i(s)\|^2\leq C_T\varepsilon_N(s).
\]
Combining the last four displays and using Theorem~\ref{thm:poc-strong}, we obtain
\[
\frac1N\sum_{i=1}^N \E\|\Delta_i(s)\|^2
\leq
C_T\varepsilon_N(s)+\frac{C_T}{N}
\leq
\frac{C_T}{N}
\]
for all $s\in[0,T]$. Integrating in time gives
\[
\sum_{i=1}^N \E \int_0^T \|\Delta_i(s)\|^2\,ds
\leq
C_T,
\]
which proves \eqref{eq:path-kl}. The fixed-time bound \eqref{eq:hn} follows by
data processing, since evaluation at time $t$ is a measurable map from path
space to $(\R^{2d})^N$.
\end{proof}

\begin{corollary}[One-particle total-variation error]\label{cor:poc-tv}
Let $\rho_{t,1}^N$ be the first marginal of $\rho_t^N$. Under the assumptions
of Theorem~\ref{thm:poc},
\[
\|\rho_{t,1}^N-f_t\|_{\TV}
\leq
\sqrt{\frac{2C_T}{N}}
\qquad
\forall t\in[0,T].
\]
\end{corollary}

\begin{proof}
For any probability measure $P^N$ on $(\R^{2d})^N$ and any product measure
$Q^{\otimes N}$, the chain rule for relative entropy gives
\[
\KL(P^N\,\|\,Q^{\otimes N})
=
\sum_{i=1}^N
\E_{P^N}\Bigl[
\KL\bigl(P_{i\,|\,1:i-1}^N(\cdot\mid Z_1,\dots,Z_{i-1})\,\|\,Q\bigr)
\Bigr].
\]
Since $\nu\mapsto \KL(\nu\|Q)$ is convex,
\[
\E_{P^N}\Bigl[
\KL\bigl(P_{i\,|\,1:i-1}^N(\cdot\mid Z_1,\dots,Z_{i-1})\,\|\,Q\bigr)
\Bigr]
\geq
\KL(P_i^N\,\|\,Q).
\]
Applying this to $P^N=\rho_t^N$ and $Q=f_t$, and using exchangeability of
$\rho_t^N$, we obtain
\[
\KL(\rho_{t,1}^N\,\|\,f_t)
\leq
\frac{1}{N}\KL(\rho_t^N\,\|\,f_t^{\otimes N}).
\]
Now use \eqref{eq:hn} and Pinsker's inequality.
\end{proof}

\section{Stationary Mean-Field Limit and Mean-Field LSI}

\subsection{Finite-particle invariant laws via Harris theory}

The stationary arguments below require the finite-particle invariant laws to be
constructed inside the present document. We therefore record the same
Lyapunov/minorization/Harris package used in the Volume~3 convergence program,
now specialized to the continuous kinetic diffusion
\eqref{eq:finite-x}--\eqref{eq:finite-v}.

\begin{proposition}[Compact-set minorization for the finite-particle diffusion]\label{prop:finite-minorization}
Fix \(N\geq2\), \(t_0>0\), and a compact set
\(K\subset(\R^{2d})^N\). Then there exist
\(\varepsilon_K>0\) and a probability measure \(\nu_K\) on \((\R^{2d})^N\) such
that
\[
P_{t_0}^N(z,A)\geq \varepsilon_K\,\nu_K(A)
\qquad
\forall z\in K
\]
for every Borel set \(A\subset(\R^{2d})^N\).
\end{proposition}

\begin{proof}
The generator \(\mathcal L_N\) has smooth coefficients, the velocity noise is
non-degenerate in every particle, and the transport field \(v\cdot\nabla_x\)
propagates that noise to the position variables. Hence the kinetic diffusion
satisfies the standard H\"ormander bracket condition, so for every \(t_0>0\) the
transition kernel \(P_{t_0}^N\) admits a smooth density
\[
P_{t_0}^N(z,dz')=p_{t_0}^N(z,z')\,dz'
\]
that is jointly continuous and strictly positive.
Choose a Euclidean ball \(B\subset(\R^{2d})^N\) with positive Lebesgue measure.
Since \(K\times B\) is compact and \(p_{t_0}^N>0\) on that set, the minimum
\[
\varepsilon_K:=|B|\inf_{(z,z')\in K\times B} p_{t_0}^N(z,z')
\]
is strictly positive. Let \(\nu_K\) be normalized Lebesgue measure on \(B\).
Then for every Borel \(A\),
\[
P_{t_0}^N(z,A)
=
\int_A p_{t_0}^N(z,z')\,dz'
\geq
\varepsilon_K\,\nu_K(A)
\qquad
\forall z\in K,
\]
which is the required minorization.
\end{proof}

\begin{theorem}[Lebesgue irreducibility and aperiodicity]\label{thm:finite-irreducible}
For every \(N\geq2\), the finite-particle kinetic semigroup \(P_t^N\) is
irreducible with respect to Lebesgue measure and is aperiodic on
\((\R^{2d})^N\).
\end{theorem}

\begin{proof}
Fix \(t_0>0\). By Proposition~\ref{prop:finite-minorization}, every compact set
is a small set for \(P_{t_0}^N\). Because the transition density
\(p_{t_0}^N(z,z')\) is strictly positive, any non-empty open set has positive
\(P_{t_0}^N(z,\cdot)\)-mass from every starting point \(z\), which is exactly
irreducibility with respect to Lebesgue measure. Choosing a compact ball \(K\) and applying
Proposition~\ref{prop:finite-minorization} with a ball \(B\subset K\) shows that
\(P_{t_0}^N(z,K)>0\) for all \(z\in K\). This one-step return to a small set
rules out periodicity, so the semigroup is aperiodic.
\end{proof}

\begin{proposition}[Uniform kinetic structure for the finite-particle drift]\label{prop:finite-harris-struct}
For every \(N\geq2\), the generator of the finite-particle diffusion can be
written as
\[
\mathcal L_N\phi
=
V\cdot \nabla_X\phi
+\bigl(-\nabla U_N(X)-A_N(X)V\bigr)\cdot \nabla_V\phi
+\gamma T\Delta_V\phi,
\qquad
T=\frac{\sigma^2}{2\gamma},
\]
where
\[
U_N(X):=\sum_{i=1}^N U(x_i),
\qquad
A_N(X):=\gamma I_{Nd}+\nu_{\mathrm{visc}}L(X)\otimes I_d.
\]
Moreover:
\begin{enumerate}[label=\textnormal{(\roman*)},leftmargin=1.5em]
\item \(U_N\in C^2((\R^d)^N)\), \(\|\nabla^2 U_N\|_{\mathrm{op}}\leq L_U\), and
\[
\langle X,\nabla U_N(X)\rangle
\geq
\alpha_U\|X\|^2-N\beta_U;
\]
\item \(A_N(X)\) depends smoothly on \(X\), and
\[
\|A_N(X)\|_{\mathrm{op}}\leq \gamma+2\nu_{\mathrm{visc}};
\]
\item every eigenvalue of \(A_N(X)\) lies in
\([\,\gamma,\gamma+2\nu_{\mathrm{visc}}\,]\), uniformly in \(N\) and \(X\);
\item the diffusion matrix is constant and non-degenerate in the velocity
variables.
\end{enumerate}
\end{proposition}

\begin{proof}
The representation of \(\mathcal L_N\) follows directly from
Definition~\ref{def:finiteN} and Definition~\ref{def:matrices}. Item \(i\) is
just Assumption~\ref{ass:potential}, summed over the particles. Item \(ii\)
follows from Lemma~\ref{lem:weighted-sym}, because the spectrum of
\(\widetilde L(X)\) is contained in \([0,2]\). Item \(iii\) is exactly
Lemma~\ref{lem:weighted-sym}(iv). Item \(iv\) is immediate from the form of the
noise term \(\sigma\,dW_i\).
\end{proof}

\begin{theorem}[Finite-particle Harris theorem]\label{thm:finite-harris}
For every \(N\geq2\), the diffusion
\eqref{eq:finite-x}--\eqref{eq:finite-v} admits a unique invariant law
\(\pi_N\). Moreover there exists a function \(V_N:(\R^{2d})^N\to[1,\infty)\),
constants \(c_-,c_+,\lambda_V,C_V>0\), depending only on
\((\alpha_U,\beta_U,L_U,\gamma,\sigma,\nu_{\mathrm{visc}},k_*,k^*)\) and
independent of \(N\), such that
\begin{align}
V_N(X,V)
&\asymp
1+\frac1N\sum_{i=1}^N \bigl(\|x_i\|^2+\|v_i\|^2\bigr),
\label{eq:VN-equiv}\\
\mathcal L_N V_N
&\leq
-\lambda_V V_N+C_V.
\label{eq:VN-drift}
\end{align}
Consequently,
\begin{equation}\label{eq:VN-invariant}
\int V_N\,d\pi_N\leq \frac{C_V}{\lambda_V}.
\end{equation}
\end{theorem}

\begin{proof}
We verify the exact structural hypotheses used in the classical Harris theory
for underdamped Langevin diffusions.

\smallskip
\noindent
\textbf{Step 1: kinetic-Langevin structure.}
Proposition~\ref{prop:finite-harris-struct} shows that the generator has the
standard kinetic form with uniformly convex confining potential \(U_N\),
constant non-degenerate velocity noise, and an affine velocity drift
\(-A_N(X)V\) whose spectrum stays in the fixed interval
\([\,\gamma,\gamma+2\nu_{\mathrm{visc}}\,]\). In particular, the only constants
entering the classical Lyapunov construction are
\((\alpha_U,\beta_U,L_U,\gamma,\sigma,\nu_{\mathrm{visc}},k_*,k^*)\), and all
of them are independent of \(N\).

\smallskip
\noindent
\textbf{Step 2: non-explosion and smooth semigroup.}
Proposition~\ref{prop:wp} gives a unique global strong solution for every
initial condition. Since the diffusion matrix is constant in the velocity
variables and the coefficients are smooth, the semigroup is the standard smooth
kinetic semigroup covered by the Roberts-Tweedie/Meyn-Tweedie Harris framework
for underdamped diffusions.

\smallskip
\noindent
\textbf{Step 3: Lyapunov function.}
For this kinetic class, the classical proof constructs a quadratic cross-term
Lyapunov function of the form
\[
\widetilde V_N(X,V)
:=
1+\eta_1\frac1N U_N(X)
+\eta_2\frac1N\|V\|^2
+\eta_3\frac1N\langle X,V\rangle,
\]
with coefficients \(\eta_1,\eta_2,\eta_3>0\) depending only on the structural
constants from Step 1. The proof uses only the coercive bound on
\(\langle X,\nabla U_N(X)\rangle\), the Hessian bound
\(\|\nabla^2 U_N\|_{\mathrm{op}}\le L_U\), the uniform spectral window for
\(A_N(X)\), and the non-degenerate velocity diffusion. Because all of those
bounds are uniform in \(N\), the same coefficient choice gives constants
\(c_-,c_+,\lambda_V,C_V>0\), independent of \(N\), such that
\[
c_-\!\left(
1+\frac1N\sum_{i=1}^N(\|x_i\|^2+\|v_i\|^2)
\right)
\leq
\widetilde V_N(X,V)
\leq
c_+\!\left(
1+\frac1N\sum_{i=1}^N(\|x_i\|^2+\|v_i\|^2)
\right)
\]
and
\[
\mathcal L_N\widetilde V_N\leq -\lambda_V\widetilde V_N+C_V.
\]
This is exactly \eqref{eq:VN-equiv}--\eqref{eq:VN-drift}, after renaming
\(\widetilde V_N\) as \(V_N\).

\smallskip
\noindent
\textbf{Step 4: Harris theorem and invariant-law bound.}
Proposition~\ref{prop:finite-minorization} and
Theorem~\ref{thm:finite-irreducible} provide the small-set, irreducibility, and
aperiodicity hypotheses. The Harris theorem
\cite{meyn2012markov,roberts1996exponential} therefore yields existence and
uniqueness of the invariant probability law \(\pi_N\) and geometric ergodicity
of \(P_t^N\). Integrating \eqref{eq:VN-drift} against \(\pi_N\) gives
\eqref{eq:VN-invariant}.
\end{proof}

\subsection{Exchangeability and finite de Finetti reduction}

\begin{theorem}[Exchangeability of the invariant law]\label{thm:exchangeable}
For every $N\geq2$, the invariant law $\pi_N$ from
Theorem~\ref{thm:finite-harris} is exchangeable:
\[
\sigma_\#\pi_N=\pi_N
\qquad
\forall \sigma\in S_N,
\]
where $\sigma$ acts on $(\R^{2d})^N$ by permutation of particle labels.
\end{theorem}

\begin{proof}
Fix a permutation $\sigma\in S_N$ and let
\[
\Sigma_\sigma(x_1,v_1,\dots,x_N,v_N)
:=
(x_{\sigma(1)},v_{\sigma(1)},\dots,x_{\sigma(N)},v_{\sigma(N)}).
\]
Because the confining force $-\nabla U(x_i)$, the friction $-\gamma v_i$, the
velocity noise $\sigma\,dW_i$, and the row-normalized viscous force
\[
F_i(X,V)=\nu_{\mathrm{visc}}\sum_{j\neq i}\omega_{ij}(X)(v_j-v_i)
\]
are defined by the same formula for every label $i$, the generator
$\mathcal L_N$ of \eqref{eq:finite-x}--\eqref{eq:finite-v} commutes with every
permutation:
\[
\mathcal L_N(\varphi\circ \Sigma_\sigma)
=
(\mathcal L_N\varphi)\circ \Sigma_\sigma
\qquad
\forall \varphi\in C_c^\infty((\R^{2d})^N).
\]
Equivalently, if $P_t^N$ is the Markov semigroup of the particle system, then
\[
P_t^N(\varphi\circ \Sigma_\sigma)=(P_t^N\varphi)\circ \Sigma_\sigma.
\]
Let \(\pi_N\) be the invariant law from Theorem~\ref{thm:finite-harris}. Then
for every bounded measurable \(\varphi\),
\begin{align*}
\int \varphi\,d(\sigma_\#\pi_N)
&=
\int \varphi\circ\Sigma_\sigma\,d\pi_N\\
&=
\int P_t^N(\varphi\circ\Sigma_\sigma)\,d\pi_N\\
&=
\int (P_t^N\varphi)\circ\Sigma_\sigma\,d\pi_N\\
	&=
	\int P_t^N\varphi\,d(\sigma_\#\pi_N).
	\end{align*}
	Thus $\sigma_\#\pi_N$ is another invariant probability measure for the same
	semigroup. By the uniqueness of the invariant law $\pi_N$,
	$\sigma_\#\pi_N=\pi_N$.
\end{proof}

\begin{theorem}[Finite de Finetti representation]\label{thm:definetti}
For each $N\geq2$ there exists a probability measure $Q_N$ on
$\mcP(\R^{2d})$ such that, for every $1\leq k\leq N$,
\[
d_{\TV}\!\left(\pi_{N,k},\int_{\mcP(\R^{2d})}\mu^{\otimes k}\,Q_N(d\mu)\right)
\leq
\frac{k(k-1)}{2N},
\]
where $\pi_{N,k}$ is the law of the first $k$ particles under $\pi_N$. One may
take $Q_N$ to be the law of the empirical measure
\[
L_N:=\frac1N\sum_{i=1}^N \delta_{(X_i,V_i)}
\qquad
\text{under }\pi_N.
\]
\end{theorem}

\begin{proof}
The state space $\R^{2d}$ is a standard Borel space, and
Theorem~\ref{thm:exchangeable} shows that $\pi_N$ is exchangeable on
$(\R^{2d})^N$. The finite de Finetti theorem of Diaconis and Freedman
\cite{diaconis1980finite} therefore applies directly. It gives a probability
measure $Q_N$ on $\mcP(\R^{2d})$ such that
\[
d_{\TV}\!\left(\pi_{N,k},\int\mu^{\otimes k}\,Q_N(d\mu)\right)
\leq
\frac{k(k-1)}{2N}
\qquad
\forall\,1\leq k\leq N.
\]
For the canonical choice, define the empirical-measure map
\[
\mathfrak L_N:(\R^{2d})^N\to \mcP(\R^{2d}),
\qquad
\mathfrak L_N(z_1,\dots,z_N):=\frac1N\sum_{i=1}^N \delta_{z_i},
\]
and set $Q_N:=(\mathfrak L_N)_\#\pi_N$. This is exactly the law of $L_N$ under
$\pi_N$, and the same Diaconis-Freedman theorem gives the stated bound for
this canonical mixing measure.
\end{proof}

\begin{corollary}[Exact barycentric representation of the first marginal]\label{cor:barycenter}
If $\pi_{N,1}$ denotes the first marginal of $\pi_N$, then
\[
\pi_{N,1}
=
\int_{\mcP(\R^{2d})}\mu\,Q_N(d\mu).
\]
\end{corollary}

\begin{proof}
Apply Theorem~\ref{thm:definetti} with \(k=1\). Since
\[
\frac{1(1-1)}{2N}=0,
\]
the total-variation distance is zero:
\[
d_{\TV}\!\left(\pi_{N,1},\int \mu\,Q_N(d\mu)\right)=0.
\]
Two probability measures at total-variation distance zero are equal, so the
first marginal is exactly the barycenter of \(Q_N\).
\end{proof}

\begin{lemma}[Conditioning on the empirical measure]\label{lem:emp-cond}
Let $G:\R^{2d}\times \mcP(\R^{2d})\to\R$ be bounded and measurable. Under the
exchangeable law $\pi_N$,
\[
\E_{\pi_N}[G((X_1,V_1),L_N)]
=
\E_{\pi_N}\!\left[\int_{\R^{2d}} G(z,L_N)\,L_N(dz)\right].
\]
\end{lemma}

\begin{proof}
By exchangeability, each index \(i\) has the same law as index \(1\), so
\[
\E_{\pi_N}[G((X_1,V_1),L_N)]
=
\frac1N\sum_{i=1}^N \E_{\pi_N}[G((X_i,V_i),L_N)].
\]
Inside the expectation, the empirical average is exactly the integral against
the empirical measure:
\[
\frac1N\sum_{i=1}^N G((X_i,V_i),L_N)
=
\int_{\R^{2d}} G(z,L_N)\,L_N(dz).
\]
Taking expectations of both sides gives the claimed identity.
\end{proof}

\subsection{Tightness, identification, and the stationary thermodynamic limit}

\begin{theorem}[Tightness of the stationary marginals]\label{thm:tight}
The sequence of first marginals $\{\pi_{N,1}\}_{N\geq2}$ is tight in
$\mcP(\R^{2d})$.
\end{theorem}

\begin{proof}
By Theorem~\ref{thm:finite-harris}, the invariant law \(\pi_N\) satisfies
\eqref{eq:VN-invariant}. Using the equivalence \eqref{eq:VN-equiv}, there is a
constant \(C_2<\infty\), independent of \(N\), such that
\[
\int \frac1N\sum_{i=1}^N \bigl(\|x_i\|^2+\|v_i\|^2\bigr)\,\pi_N(dX\,dV)
\leq C_2.
\]
Using exchangeability from Theorem~\ref{thm:exchangeable}, this becomes
\[
\int_{\R^{2d}} (\|x\|^2+\|v\|^2)\,\pi_{N,1}(dx\,dv)\leq C_2
\qquad
\forall N\geq2.
\]
Now fix $\varepsilon>0$ and choose
$R_\varepsilon:=\sqrt{C_2/\varepsilon}$. Markov's inequality gives
\[
\pi_{N,1}\!\left(\{(x,v):\|x\|^2+\|v\|^2>R_\varepsilon^2\}\right)
\leq
\frac{C_2}{R_\varepsilon^2}
=
\varepsilon.
\]
Hence the compact ball
$K_\varepsilon:=\{(x,v):\|x\|^2+\|v\|^2\leq R_\varepsilon^2\}$ satisfies
\[
\pi_{N,1}(K_\varepsilon)\geq 1-\varepsilon
\qquad
\forall N\geq2.
\]
Prokhorov's theorem therefore yields tightness.
\end{proof}

\subsection{Linear stationary problem and nonlinear alignment fixed point}

\begin{proposition}[Linear stationary law for a prescribed alignment field]\label{prop:linear-mf-stationary}
Let \(a:\R^d\to\R^d\) be bounded and globally Lipschitz, and define
\[
\mathcal L^a\phi(x,v)
:=
v\cdot \nabla_x\phi
+\bigl(-\nabla U(x)-(\gamma+\nu_{\mathrm{visc}})v+\nu_{\mathrm{visc}}a(x)\bigr)
\cdot \nabla_v\phi
+\frac{\sigma^2}{2}\Delta_v\phi.
\]
Then the diffusion generated by \(\mathcal L^a\) admits a unique invariant law
\(\pi^a\in\mcP(\R^{2d})\). Moreover:
\begin{enumerate}[label=\textnormal{(\roman*)},leftmargin=1.5em]
\item \(\pi^a\) has finite second moments, with a bound depending only on
\(\|a\|_\infty\) and the structural constants;
\item \(\pi^a\) admits a smooth strictly positive density \(\rho^a\);
\item if \(a\) ranges over a bounded set in \(W^{1,\infty}\), then the
second-moment bound is uniform on that set.
\end{enumerate}
\end{proposition}

\begin{proof}
The drift of \(\mathcal L^a\) is the confining kinetic Langevin drift
\((v,-\nabla U(x)-\gamma v)\) plus the bounded perturbation
\(\nu_{\mathrm{visc}}(a(x)-v)\) in velocity. The same quadratic cross-term
Lyapunov function used in Theorem~\ref{thm:finite-harris} yields
\[
\mathcal L^a W(x,v)\le -c_0 W(x,v)+C_0\bigl(1+\|a\|_\infty^2\bigr)
\]
for a quadratic \(W\sim 1+\|x\|^2+\|v\|^2\). The diffusion matrix is constant in
velocity and the coefficients are smooth, so the bracket family generated by
\(\{\partial_{v_1},\dots,\partial_{v_d}\}\) and the drift field spans the full
tangent space:
\[
\bigl[v\cdot \nabla_x
+(-\nabla U(x)-(\gamma+\nu_{\mathrm{visc}})v+\nu_{\mathrm{visc}}a(x))\cdot\nabla_v,\,
\partial_{v_k}\bigr]
=
\partial_{x_k}-(\gamma+\nu_{\mathrm{visc}})\partial_{v_k}.
\]
Hence Hörmander's condition holds globally. The transition kernel therefore has
a smooth strictly positive density for every \(t>0\), and every compact set is
small by the same compact-ball argument used in the finite-particle
minorization theorem. Harris theory then gives existence and uniqueness of the
invariant probability law \(\pi^a\) together with the stated quadratic moment
bound. Finally, invariance and the strictly positive smooth transition density
imply that \(\pi^a\) itself admits a smooth strictly positive density
\(\rho^a\).
\end{proof}

\begin{definition}[Stationary alignment map]\label{def:mf-alignment-map}
For a bounded Lipschitz field \(a\), let \(\rho^a\) be the density of the
invariant law \(\pi^a\) from Proposition~\ref{prop:linear-mf-stationary}. Define
\[
(\mathcal T a)(x)
:=
\frac{\iint_{\R^{2d}} K(x,y)\,q\,\rho^a(y,q)\,\dd y\,\dd q}
{\iint_{\R^{2d}} K(x,y)\,\rho^a(y,q)\,\dd y\,\dd q}.
\]
\end{definition}

\begin{lemma}[Invariant ball for the alignment map]\label{lem:mf-map-ball}
There exists \(R_*>0\), depending only on the structural constants, such that
if
\[
\mathcal A_{R_*}
:=
\{a\in W^{1,\infty}(\R^d;\R^d):\|a\|_\infty+\|\nabla a\|_\infty\le R_*\},
\]
then \(\mathcal T(\mathcal A_{R_*})\subseteq \mathcal A_{R_*}\).
\end{lemma}

\begin{proof}
By the kernel floor, the denominator in Definition~\ref{def:mf-alignment-map} is
bounded below by \(k_*\). The numerator is controlled by the first velocity
moment of \(\pi^a\), which is uniformly bounded on bounded \(W^{1,\infty}\)
balls by Proposition~\ref{prop:linear-mf-stationary}. Differentiating under the
integral sign and using the boundedness of \(\nabla K\) yields the same control
for \(\nabla(\mathcal T a)\). Choosing \(R_*\) larger than the resulting
constant closes the invariant-ball estimate.
\end{proof}

\begin{lemma}[Small-viscosity contraction of the alignment map]\label{lem:mf-map-contraction}
There exists \(C_{\mathrm{mf}}<\infty\), depending only on the structural
constants, such that for all \(a,b\in \mathcal A_{R_*}\),
\[
\|\mathcal T a-\mathcal T b\|_{L^\infty}
\le
C_{\mathrm{mf}}\frac{\nu_{\mathrm{visc}}}{\gamma}\|a-b\|_{L^\infty}.
\]
Consequently, if
\[
0\le \nu_{\mathrm{visc}}<\frac{\gamma}{C_{\mathrm{mf}}},
\]
then \(\mathcal T\) is a strict contraction on \(\mathcal A_{R_*}\).
\end{lemma}

\begin{proof}
Let \((X_t^a,V_t^a)\) and \((X_t^b,V_t^b)\) solve the two frozen-field SDEs
driven by the same Brownian motion:
\begin{align*}
dX_t^\sharp &= V_t^\sharp\,dt,\\
dV_t^\sharp &=
\Bigl(-\nabla U(X_t^\sharp)-(\gamma+\nu_{\mathrm{visc}})V_t^\sharp
+\nu_{\mathrm{visc}}\,\sharp(X_t^\sharp)\Bigr)\,dt
+\sigma\,dW_t,
\qquad
\sharp\in\{a,b\}.
\end{align*}
Choose the initial law to be an arbitrary coupling of \(\pi^a\) and \(\pi^b\).
Since each marginal is started from its own invariant law, both marginals remain
stationary for every \(t\ge0\). Set
\[
\Delta X_t:=X_t^a-X_t^b,
\qquad
\Delta V_t:=V_t^a-V_t^b.
\]
Then
\begin{align*}
d\Delta X_t &= \Delta V_t\,dt,\\
d\Delta V_t
&=
\Bigl(
-\bigl(\nabla U(X_t^a)-\nabla U(X_t^b)\bigr)
-(\gamma+\nu_{\mathrm{visc}})\Delta V_t\\
&\qquad\qquad
+\nu_{\mathrm{visc}}\bigl(a(X_t^a)-a(X_t^b)\bigr)
+\nu_{\mathrm{visc}}\bigl(a(X_t^b)-b(X_t^b)\bigr)
\Bigr)\,dt.
\end{align*}

Set
\[
\beta:=\frac{\gamma+\nu_{\mathrm{visc}}}{2},
\qquad
\alpha:=(\gamma+\nu_{\mathrm{visc}})\beta=\frac{(\gamma+\nu_{\mathrm{visc}})^2}{2},
\]
and define
\[
\mathscr Q(\xi,\eta):=\alpha\|\xi\|^2+2\beta\langle \xi,\eta\rangle+\|\eta\|^2.
\]
Since \(\alpha>\beta^2\), this quadratic form is equivalent to
\(\|\xi\|^2+\|\eta\|^2\). For the unforced part of the difference dynamics,
It\^o's formula gives
\begin{align*}
\frac{d}{dt}\mathscr Q(\Delta X_t,\Delta V_t)
&=
2(\alpha-(\gamma+\nu_{\mathrm{visc}})\beta)\langle \Delta X_t,\Delta V_t\rangle
+2(\beta-(\gamma+\nu_{\mathrm{visc}}))\|\Delta V_t\|^2\\
&\quad
-2\beta\bigl\langle \Delta X_t,\nabla U(X_t^a)-\nabla U(X_t^b)\bigr\rangle
-2\bigl\langle \Delta V_t,\nabla U(X_t^a)-\nabla U(X_t^b)\bigr\rangle\\
&\quad
+2\nu_{\mathrm{visc}}\beta\langle \Delta X_t,a(X_t^a)-a(X_t^b)\rangle
+2\nu_{\mathrm{visc}}\langle \Delta V_t,a(X_t^a)-a(X_t^b)\rangle\\
&\quad
+2\nu_{\mathrm{visc}}
\bigl\langle \Delta V_t+\beta\Delta X_t,\,
a(X_t^b)-b(X_t^b)\bigr\rangle.
\end{align*}
By the choice of \(\alpha\), the cross term vanishes. Using the strong
monotonicity and Lipschitz bounds on \(\nabla U\) together with
\(\|\nabla a\|_\infty\le R_*\), we obtain
\begin{align*}
\frac{d}{dt}\mathscr Q(\Delta X_t,\Delta V_t)
&\le
-(\gamma+\nu_{\mathrm{visc}})\|\Delta V_t\|^2
-2\beta(m_U-\nu_{\mathrm{visc}}R_*)\|\Delta X_t\|^2\\
&\quad
+2(L_U+\nu_{\mathrm{visc}}R_*)\|\Delta X_t\|\,\|\Delta V_t\|\\
&\quad
+2\nu_{\mathrm{visc}}
\bigl\langle \Delta V_t+\beta\Delta X_t,\,
a(X_t^b)-b(X_t^b)\bigr\rangle.
\end{align*}
Apply Young's inequality to the mixed term with
\(\varepsilon=(\gamma+\nu_{\mathrm{visc}})/2\):
\[
2(L_U+\nu_{\mathrm{visc}}R_*)\|\Delta X_t\|\,\|\Delta V_t\|
\le
\frac{\gamma+\nu_{\mathrm{visc}}}{2}\|\Delta V_t\|^2
+\frac{2(L_U+\nu_{\mathrm{visc}}R_*)^2}{\gamma+\nu_{\mathrm{visc}}}
\|\Delta X_t\|^2.
\]
Hence
\begin{align*}
\frac{d}{dt}\mathscr Q(\Delta X_t,\Delta V_t)
&\le
-\frac{\gamma+\nu_{\mathrm{visc}}}{2}\|\Delta V_t\|^2\\
&\quad
-\left(
(\gamma+\nu_{\mathrm{visc}})(m_U-\nu_{\mathrm{visc}}R_*)
-\frac{2(L_U+\nu_{\mathrm{visc}}R_*)^2}{\gamma+\nu_{\mathrm{visc}}}
\right)\|\Delta X_t\|^2\\
&\quad
+2\nu_{\mathrm{visc}}
\bigl\langle \Delta V_t+\beta\Delta X_t,\,
a(X_t^b)-b(X_t^b)\bigr\rangle.
\end{align*}
In the small-viscosity / large-friction regime, the bracketed coefficient is
strictly positive. By the equivalence of \(\mathscr Q\) with
\(\|\Delta X_t\|^2+\|\Delta V_t\|^2\), there exists
\(\lambda_{\mathrm{cpl}}>0\), depending only on
\((m_U,L_U,\gamma,R_*)\), such that
\begin{equation}\label{eq:mf-coupling-diss}
\frac{d}{dt}\mathscr Q(\Delta X_t,\Delta V_t)
\le
-\lambda_{\mathrm{cpl}}\mathscr Q(\Delta X_t,\Delta V_t)
+2\nu_{\mathrm{visc}}
\bigl\langle
\Delta V_t+\beta\Delta X_t,\,
a(X_t^b)-b(X_t^b)
\bigr\rangle.
\end{equation}

Taking expectations in \eqref{eq:mf-coupling-diss} and using Young's inequality
gives
\[
\frac{d}{dt}\E \mathscr Q(\Delta X_t,\Delta V_t)
\le
-\frac{\lambda_{\mathrm{cpl}}}{2}\E \mathscr Q(\Delta X_t,\Delta V_t)
+C_0\frac{\nu_{\mathrm{visc}}^2}{\lambda_{\mathrm{cpl}}}\|a-b\|_{L^\infty}^2.
\]
Hence
\[
\E \mathscr Q(\Delta X_t,\Delta V_t)
\le
e^{-\lambda_{\mathrm{cpl}}t/2}\E \mathscr Q(\Delta X_0,\Delta V_0)
+C_1\frac{\nu_{\mathrm{visc}}^2}{\lambda_{\mathrm{cpl}}^2}\|a-b\|_{L^\infty}^2.
\]
Because both marginals are stationary, the law of
\((X_t^a,V_t^a,X_t^b,V_t^b)\) is a coupling of \(\pi^a\) and \(\pi^b\) for
every \(t\). Letting \(t\to\infty\) and using the equivalence of
\(\mathscr Q\) with \(\|\Delta X_t\|^2+\|\Delta V_t\|^2\), we obtain
\begin{equation}\label{eq:mf-w2-stability}
W_2(\pi^a,\pi^b)
\le
C_2\frac{\nu_{\mathrm{visc}}}{\lambda_{\mathrm{cpl}}}\|a-b\|_{L^\infty}.
\end{equation}
The same quadratic-form construction gives
\(\lambda_{\mathrm{cpl}}\gtrsim \gamma\) uniformly in the small-viscosity
regime, so
\[
W_2(\pi^a,\pi^b)
\le
C_3\frac{\nu_{\mathrm{visc}}}{\gamma}\|a-b\|_{L^\infty}.
\]

Fix \(x\in\R^d\). Write
\[
N_a(x):=\iint K(x,y)\,q\,\pi^a(dy\,dq),
\qquad
D_a(x):=\iint K(x,y)\,\pi^a(dy\,dq),
\]
so that \((\mathcal T a)(x)=N_a(x)/D_a(x)\), and similarly for \(b\). Since
\(K\ge k_*\), we have \(D_a(x),D_b(x)\ge k_*\). Proposition~\ref{prop:linear-mf-stationary}
gives a uniform second-moment bound on \(\pi^a,\pi^b\) over
\(\mathcal A_{R_*}\). Therefore, for any coupling \(\Pi\) of \(\pi^a\) and
\(\pi^b\),
\begin{align*}
\|N_a(x)-N_b(x)\|
&\le
\iint
\|K(x,y)q-K(x,y')q'\|\,\Pi(dy\,dq\,dy'\,dq')\\
&\le
k^*\iint \|q-q'\|\,d\Pi
+L_K\left(\iint \|q\|^2\,d\Pi\right)^{1/2}
\left(\iint \|y-y'\|^2\,d\Pi\right)^{1/2}\\
&\le
C_4\,W_2(\pi^a,\pi^b).
\end{align*}
Likewise,
\[
|D_a(x)-D_b(x)|
\le
L_K\,W_2(\pi^a,\pi^b).
\]
Using
\[
\frac{N_a}{D_a}-\frac{N_b}{D_b}
=
\frac{N_a-N_b}{D_a}
+N_b\frac{D_b-D_a}{D_aD_b},
\]
the denominator lower bound \(k_*\), and the uniform first-moment bound on
\(\pi^b\), we conclude that
\[
\|(\mathcal T a)(x)-(\mathcal T b)(x)\|
\le
C_5\,W_2(\pi^a,\pi^b)
\le
C_{\mathrm{mf}}\frac{\nu_{\mathrm{visc}}}{\gamma}\|a-b\|_{L^\infty}.
\]
Taking the supremum over \(x\) yields the claim.
\end{proof}

\begin{proposition}[Unique stationary mean-field law]\label{prop:mf-unique}
There exists a unique stationary probability law \(\pi_\infty\in\mcP(\R^{2d})\)
for the mean-field equation \eqref{eq:mf-pde}. Equivalently,
\(\pi_\infty\) is the unique probability law such that
\begin{equation}\label{eq:mf-stationary-weak}
\int_{\R^{2d}} \mathcal L_{\pi_\infty}\phi(z)\,\pi_\infty(dz)=0
\qquad
\forall \phi\in C_c^\infty(\R^{2d}),
\end{equation}
where
\[
\mathcal L_\mu \phi(x,v)
:=
v\cdot \nabla_x \phi
+\bigl(-\nabla U(x)-\gamma v+F[\mu](x,v)\bigr)\cdot \nabla_v \phi
+\frac{\sigma^2}{2}\Delta_v \phi.
\]
\end{proposition}

\begin{proof}
Fix the invariant ball \(\mathcal A_{R_*}\) from
Lemma~\ref{lem:mf-map-ball}. If
\[
0\le \nu_{\mathrm{visc}}<\frac{\gamma}{C_{\mathrm{mf}}},
\]
then Lemma~\ref{lem:mf-map-contraction} shows that \(\mathcal T\) is a strict
contraction on the complete metric space \(\mathcal A_{R_*}\). Banach's
fixed-point theorem therefore yields a unique field \(a_\infty\in\mathcal A_{R_*}\)
such that
\[
\mathcal T a_\infty=a_\infty.
\]
Let \(\pi^{a_\infty}\) be the invariant law of the linear problem from
Proposition~\ref{prop:linear-mf-stationary}. By the definition of \(\mathcal T\),
the associated density \(\rho^{a_\infty}\) satisfies
\[
\nu_{\mathrm{visc}}(a_\infty(x)-v)=F[\pi^{a_\infty}](x,v),
\]
so \(\pi^{a_\infty}\) satisfies exactly the nonlinear stationary weak equation
\eqref{eq:mf-stationary-weak}. This gives existence of a stationary mean-field
law \(\pi_\infty:=\pi^{a_\infty}\). Conversely, any stationary solution
\(\bar\pi\) induces through its alignment profile a fixed point of
\(\mathcal T\), hence by uniqueness of the fixed point must coincide with
\(\pi_\infty\).
\end{proof}

\begin{theorem}[Stationary propagation of chaos]\label{thm:stationary-poc}
The sequence of first marginals $\pi_{N,1}$ converges weakly to the unique
stationary mean-field law $\pi_\infty$ from
Proposition~\ref{prop:mf-unique}. Equivalently,
\[
\pi_{N,1}\Rightarrow \pi_\infty
\qquad\text{as }N\to\infty.
\]
\end{theorem}

\begin{proof}
We now carry out the tightness, identification, and uniqueness program in full.

\smallskip
\noindent
\textbf{Step 1: tightness.}
By Theorem~\ref{thm:tight}, the family $\{\pi_{N,1}\}_{N\geq2}$ is tight.
Hence every subsequence admits a weakly convergent sub-subsequence.

\smallskip
\noindent
\textbf{Step 2: stationarity identity for the finite-$N$ system.}
Fix $\phi\in C_c^\infty(\R^{2d})$ and set
\[
\Phi(X_1,V_1,\dots,X_N,V_N):=\phi(X_1,V_1).
\]
Since $\pi_N$ is invariant,
\[
\int \mathcal L_N\Phi\,d\pi_N=0.
\]
Writing
\[
\mathcal L_\mu \phi(x,v)
:=
v\cdot \nabla_x \phi
+\bigl(-\nabla U(x)-\gamma v+F[\mu](x,v)\bigr)\cdot \nabla_v \phi
+\frac{\sigma^2}{2}\Delta_v \phi,
\]
this gives
\[
\int \mathcal L_0\phi(z_1)\,\pi_N(dZ)
+
\int F_1^N(Z)\cdot \nabla_v\phi(z_1)\,\pi_N(dZ)
=0,
\]
where $\mathcal L_0$ is the backbone generator without the mean-field term.
By Lemma~\ref{lem:force-defect}, evaluated under the stationary law $\pi_N$,
\[
\left|
\int \bigl(F_1^N-F[L_N](z_1)\bigr)\cdot \nabla_v\phi(z_1)\,d\pi_N
\right|
\leq
\|\nabla_v\phi\|_\infty
\left(\int \|F_1^N-F[L_N](z_1)\|^2\,d\pi_N\right)^{1/2}
\leq
\frac{C_\phi}{\sqrt N}.
\]
\[
\Psi_\phi(\mu)
:=
\int_{\R^{2d}} \mathcal L_\mu\phi(z)\,\mu(dz).
\]
By Corollary~\ref{cor:barycenter} and Lemma~\ref{lem:emp-cond},
\begin{equation}\label{eq:stationary-approx}
\left|
\E_{\pi_N}\!\left[\Psi_\phi(L_N)\right]
\right|
\leq
\frac{C_\phi}{\sqrt N}.
\end{equation}

\smallskip
\noindent
\textbf{Step 4: identification of limit points.}
Let $\{N_k\}$ be a subsequence such that $\pi_{N_k,1}\Rightarrow \bar\pi$.
By the uniform second-moment bound in Theorem~\ref{thm:tight} and the canonical
choice $Q_N=\Law(L_N)$ from Theorem~\ref{thm:definetti}, the family
$\{Q_N\}_{N\geq2}$ is tight in $\mcP(\mcP(\R^{2d}))$. Passing to a further
subsequence if necessary, we may assume $Q_{N_k}\Rightarrow Q_\infty$.
Indeed,
\[
\int_{\mcP(\R^{2d})}
\left[\int_{\R^{2d}}(\|x\|^2+\|v\|^2)\,\mu(dx\,dv)\right]Q_N(d\mu)
=
\int_{\R^{2d}}(\|x\|^2+\|v\|^2)\,\pi_{N,1}(dx\,dv)
\leq C_2,
\]
so the same Markov-Prokhorov argument applies on the second level.
For compactly supported $\phi$, the map $\Psi_\phi$ is continuous on
$\mcP(\R^{2d})$ under weak convergence together with uniform second moments,
because $K$ and $\nabla_v\phi$ are bounded and the denominator in $F[\mu]$ is
uniformly bounded below by $k_*$. Passing to the limit in the previous display
therefore gives
\[
\int_{\mcP(\R^{2d})}\Psi_\phi(\mu)\,Q_\infty(d\mu)=0
\qquad
\forall \phi\in C_c^\infty(\R^{2d}).
\]
This is the stationary weak mean-field identity for the limit mixture. The
vanishing defect estimate \eqref{eq:stationary-approx}, together with the
continuity just noted, puts us in the hypotheses of the standard nonlinear
martingale-problem identification for exchangeable particle systems
(equivalently, the de Finetti/Sznitman limit-point argument; see
\cite{sznitman1991topics}). Hence
\[
Q_\infty\Bigl(\bigl\{\mu:\int_{\R^{2d}}\mathcal L_\mu\phi(z)\,\mu(dz)=0
\ \forall \phi\in C_c^\infty(\R^{2d})\bigr\}\Bigr)=1.
\]
By Proposition~\ref{prop:mf-unique}, this set consists of the single law
\(\pi_\infty\), so \(Q_\infty=\delta_{\pi_\infty}\).

\smallskip
\noindent
\textbf{Step 5: uniqueness and collapse of the mixing measure.}
Since the barycenter map is continuous on the uniform second-moment class,
\[
\bar\pi
=
\int \mu\,Q_\infty(d\mu)
=
\int \mu\,\delta_{\pi_\infty}(d\mu)
=
\pi_\infty.
\]
Thus every convergent subsequence of $\{\pi_{N,1}\}_{N\geq2}$ has the same
limit, so the whole sequence converges:
\[
\pi_{N,1}\Rightarrow \pi_\infty.
\]
\end{proof}

\begin{corollary}[Thermodynamic limit for bounded observables]\label{cor:obs}
For every bounded continuous observable $\phi:\R^{2d}\to\R$,
\[
\lim_{N\to\infty}
\E_{\pi_N}\!\left[\frac1N\sum_{i=1}^N \phi(X_i,V_i)\right]
=
\int_{\R^{2d}} \phi\,d\pi_\infty.
\]
\end{corollary}

\begin{proof}
Exchangeability gives
\[
\E_{\pi_N}\!\left[\frac1N\sum_{i=1}^N \phi(X_i,V_i)\right]
=
\int_{\R^{2d}} \phi\,d\pi_{N,1}.
\]
Now apply Theorem~\ref{thm:stationary-poc} and the definition of weak
convergence.
\end{proof}

\begin{definition}[Backbone hypocoercive coefficients for the mean-field LSI]\label{def:mf-lsi-coeffs}
Set
\[
M:=L_U,
\qquad
\varepsilon:=\min\left\{1,\frac{\gamma}{16(M+\gamma)^2}\right\},
\]
and define
\[
a:=\frac{\gamma\varepsilon}{4},
\qquad
b:=\varepsilon,
\qquad
c:=\varepsilon^2.
\]
These are exactly the coefficients from
Theorem~\ref{thm:backbone-detailed}, recorded here so that the mean-field
statement is self-contained.
\end{definition}

\begin{proposition}[Applicability of Huang--Kopfer--Ren to the frozen mean-field generator]\label{prop:hkr-applicability}
Let \(a_\infty\in \mathcal A_{R_*}\) be the unique fixed point from
Proposition~\ref{prop:mf-unique}, and let
\[
\mathcal L^{a_\infty}\phi
:=
v\cdot \nabla_x \phi
+\bigl(-\nabla U(x)-(\gamma+\nu_{\mathrm{visc}})v+\nu_{\mathrm{visc}}a_\infty(x)\bigr)
\cdot \nabla_v \phi
+\frac{\sigma^2}{2}\Delta_v \phi.
\]
Then, in the notation of Huang--Kopfer--Ren
\cite[Theorem~3.5]{huang2026logsobolev}, the coefficients of
\(\mathcal L^{a_\infty}\) satisfy:
\begin{enumerate}[label=\textnormal{(\roman*)},leftmargin=1.5em]
\item the measure-interaction constant is \(K_I=0\);
\item the diffusion bounds \cite[(57), (63)]{huang2026logsobolev} hold with
\(\delta_1=\delta_2=\sigma^2\);
\item the dissipativity condition \cite[(C) with \(R=0\)]{huang2026logsobolev}
holds.
\end{enumerate}
Consequently, Huang--Kopfer--Ren's frozen-measure logarithmic Sobolev theorem
applies to the invariant law \(\pi_\infty=\pi^{a_\infty}\).
\end{proposition}

\begin{proof}
Item \textnormal{(i)} is immediate: the field \(a_\infty\) is frozen, so the
generator \(\mathcal L^{a_\infty}\) is an ordinary kinetic diffusion with no
remaining dependence on the current law. Hence the interaction parameter
\(K_I\) in \cite{huang2026logsobolev} is exactly zero.

For item \textnormal{(ii)}, the diffusion acts only in the velocity variables
and is constant:
\[
\Sigma
:=
\begin{pmatrix}
0\\
\sigma I_d
\end{pmatrix},
\qquad
\Sigma\Sigma^\top
=
\begin{pmatrix}
0 & 0\\
0 & \sigma^2 I_d
\end{pmatrix}.
\]
Therefore the non-degenerate block in the diffusive coordinates is exactly
\(\sigma^2 I_d\), so the upper and lower ellipticity bounds required in
\cite[(57), (63)]{huang2026logsobolev} hold with the same constant
\(\sigma^2\).

For item \textnormal{(iii)}, consider two solutions of the frozen-field SDE
driven by the same Brownian motion and write their difference as
\((\Delta X_t,\Delta V_t)\). The proof of
Lemma~\ref{lem:mf-map-contraction}, specialized to the case \(a=b=a_\infty\),
shows that the forcing term disappears and the quadratic form
\[
\mathscr Q(\xi,\eta)=\alpha\|\xi\|^2+2\beta\langle \xi,\eta\rangle+\|\eta\|^2
\]
satisfies
\[
\frac{d}{dt}\mathscr Q(\Delta X_t,\Delta V_t)
\le
-\lambda_{\mathrm{cpl}}\mathscr Q(\Delta X_t,\Delta V_t)
\]
for some \(\lambda_{\mathrm{cpl}}>0\) depending only on the structural
constants. Equivalently, the drift of \(\mathcal L^{a_\infty}\) is uniformly
dissipative with respect to the positive definite metric induced by
\(\mathscr Q\). This is exactly the content of
\cite[(C) with \(R=0\)]{huang2026logsobolev} in the frozen, interaction-free
case.

Thus all hypotheses of \cite[Theorem~3.5]{huang2026logsobolev} are satisfied
for the invariant law of \(\mathcal L^{a_\infty}\), namely \(\pi_\infty\).
\end{proof}

\begin{theorem}[Mean-field logarithmic Sobolev inequality in the small-viscosity regime]\label{thm:mf-lsi}
Let
\[
\Gamma_{\mathrm{hypo}}^\infty(g,g)
:=
a\|\nabla_v g\|^2
+2b\,\nabla_x g\cdot\nabla_v g
+c\|\nabla_x g\|^2,
\qquad
ac>b^2,
\]
with the coefficients \(a,b,c\) from
Definition~\ref{def:mf-lsi-coeffs}. Assume in addition that
\[
0\leq \nu_{\mathrm{visc}}\leq \nu_*,
\]
with \(\nu_*\) chosen small enough that Proposition~\ref{prop:mf-unique}
applies. Then the stationary mean-field law \(\pi_\infty\) satisfies
\[
\Ent_{\pi_\infty}(g^2)
\leq
C_{\mathrm{LSI}}
\int_{\R^{2d}} \Gamma_{\mathrm{hypo}}^\infty(g,g)\,\dd\pi_\infty
\]
for every smooth compactly supported \(g\), for some constant
\(C_{\mathrm{LSI}}<\infty\) depending only on the structural constants.
\end{theorem}

\begin{proof}
Proposition~\ref{prop:hkr-applicability} shows that the frozen-field kinetic
generator associated with \(\pi_\infty\) satisfies the exact hypotheses of
Huang--Kopfer--Ren's frozen-measure logarithmic Sobolev theorem
\cite[Theorem~3.5]{huang2026logsobolev}. Hence there exists a constant
\(C_{\mathrm E}<\infty\), depending only
on the structural bounds of the frozen field ball \(\mathcal A_{R_*}\), such
that
\[
\Ent_{\pi_\infty}(g^2)
\le
C_{\mathrm E}
\int_{\R^{2d}} \bigl(\|\nabla_x g\|^2+\|\nabla_v g\|^2\bigr)\,\dd\pi_\infty
\]
for every smooth compactly supported \(g\).

Now let
\[
M_\Gamma:=
\begin{pmatrix}
c & b\\
b & a
\end{pmatrix}.
\]
Since \(ac>b^2\), the matrix \(M_\Gamma\) is positive definite. Its smallest
eigenvalue is
\[
\lambda_\Gamma
:=
\frac{a+c-\sqrt{(a-c)^2+4b^2}}{2}>0.
\]
Consequently,
\[
\Gamma_{\mathrm{hypo}}^\infty(g,g)
\ge
\lambda_\Gamma\bigl(\|\nabla_x g\|^2+\|\nabla_v g\|^2\bigr).
\]
Combining the two inequalities gives
\[
\Ent_{\pi_\infty}(g^2)
\le
\frac{C_{\mathrm E}}{\lambda_\Gamma}
\int_{\R^{2d}}\Gamma_{\mathrm{hypo}}^\infty(g,g)\,\dd\pi_\infty.
\]
Renaming \(C_{\mathrm E}/\lambda_\Gamma\) as \(C_{\mathrm{LSI}}\) proves the
claimed mean-field logarithmic Sobolev inequality.
\end{proof}

\subsection{Translated confinement and the wall-to-infinity limit}

\begin{definition}[Translated confining family]\label{def:translated-family}
Fix a reference potential \(U_*:\R^d\to\R\) satisfying
Assumption~\ref{ass:potential}. For each translation vector \(a\in\R^d\), define
\[
U_a(x):=U_*(x-a).
\]
Let \(P_t^{N,a}\), \(Q_\tau^{N,a}\), \(\pi_N^a\), and \(\pi_\infty^a\) denote,
respectively, the continuous finite-particle semigroup, the squashed observation
chain, the invariant law, and the stationary mean-field law for the kinetic core
with confining potential \(U_a\).

Define the translation maps
\[
t_a(x,v):=(x+a,v)\in\R^{2d},
\]
\[
T_a^N(x_1,v_1,\dots,x_N,v_N)
:=
\bigl(x_1+a,v_1,\dots,x_N+a,v_N\bigr)\in(\R^{2d})^N.
\]
\end{definition}

\begin{theorem}[Translation covariance of the confining family]\label{thm:translated-family}
For every \(a\in\R^d\), the translated family from
Definition~\ref{def:translated-family} satisfies the following exact covariance
relations.
\begin{enumerate}[label=\textnormal{(\roman*)},leftmargin=1.5em]
\item If \((X_i^0(t),V_i^0(t))_{i=1}^N\) solves
\eqref{eq:finite-x}--\eqref{eq:finite-v} with potential \(U_*\), then
\[
X_i^a(t):=X_i^0(t)+a,
\qquad
V_i^a(t):=V_i^0(t)
\]
solves the same system with potential \(U_a\), driven by the same Brownian
motions.
\item The continuous semigroups and squashed observation chains satisfy
\[
P_t^{N,a}(F\circ T_a^N)=(P_t^{N,0}F)\circ T_a^N,
\qquad
Q_\tau^{N,a}(F\circ T_a^N)=(Q_\tau^{N,0}F)\circ T_a^N
\]
for every bounded measurable \(F\).
\item The invariant laws satisfy
\[
T_{a\#}^N\pi_N^0=\pi_N^a.
\]
\item If \(f_t^0\) solves the mean-field equation \eqref{eq:mf-pde} with
potential \(U_*\), then
\[
f_t^a:=t_{a\#}f_t^0
\]
solves the translated mean-field equation with potential \(U_a\). In particular,
\[
t_{a\#}\pi_\infty^0=\pi_\infty^a.
\]
\item The constants in the finite-\(N\) hypocoercive logarithmic Sobolev
inequality \eqref{eq:lsi-finite}, the entropy contraction
\eqref{eq:kl-finite}, and the mean-field logarithmic Sobolev inequality from
Theorem~\ref{thm:mf-lsi} are the same for every \(a\).
\end{enumerate}
\end{theorem}

\begin{proof}
The kernel is translation invariant:
\[
K(x_i+a,x_j+a)=K(x_i,x_j).
\]
Hence
\[
d_i(X+a\mathbf 1)=d_i(X),
\qquad
\omega_{ij}(X+a\mathbf 1)=\omega_{ij}(X),
\]
and therefore the viscous force is translation covariant:
\[
F_i(X+a\mathbf 1,V)=F_i(X,V).
\]
The confining drift transforms as
\[
\nabla U_a(x+a)=\nabla U_*(x).
\]
Substituting these identities into \eqref{eq:finite-x}--\eqref{eq:finite-v}
shows that if \((X_i^0,V_i^0)\) solves the reference system, then
\((X_i^0+a,V_i^0)\) solves the translated system. This proves
\textnormal{(i)}.

For \textnormal{(ii)}, let \(Z_t^{N,0}\) be the reference process and
\(Z_t^{N,a}:=T_a^N Z_t^{N,0}\). By \textnormal{(i)}, \(Z_t^{N,a}\) has the law of
the translated process started from \(T_a^N z\) whenever \(Z_t^{N,0}\) starts
from \(z\). Therefore
\[
P_t^{N,a}(F\circ T_a^N)(z)
=
\E_z\!\bigl[F(T_a^N Z_t^{N,a})\bigr]
=
\E_{T_a^N z}\!\bigl[F(Z_t^{N,0})\bigr]
=
(P_t^{N,0}F)(T_a^N z).
\]
The observation map \(S(x,v)=(x,\psi(v))\) commutes with translations in the
position variables:
\[
S\circ T_a^N=T_a^N\circ S.
\]
Hence the same covariance relation holds for the observed chain \(Q_\tau^{N,a}\).

For \textnormal{(iii)}, let \(\mu:=T_{a\#}^N\pi_N^0\). Then for every bounded
measurable \(F\),
\[
\int P_t^{N,a}F\,d\mu
=
\int P_t^{N,a}F(T_a^N z)\,\pi_N^0(dz)
=
\int P_t^{N,0}(F\circ T_a^N)(z)\,\pi_N^0(dz)
=
\int F(T_a^N z)\,\pi_N^0(dz)
=
\int F\,d\mu,
\]
where the third equality uses invariance of \(\pi_N^0\) under \(P_t^{N,0}\).
Thus \(\mu\) is invariant for the translated dynamics. By the uniqueness part of
Theorem~\ref{thm:finite-harris}, \(\mu=\pi_N^a\).

For \textnormal{(iv)}, let \(f_t^0\) be a weak solution of \eqref{eq:mf-pde}
with potential \(U_*\), and define \(f_t^a=t_{a\#}f_t^0\). Because \(K\) is
translation invariant, the normalized kernel mass transforms by
\[
m_{f_t^a}(x+a)=m_{f_t^0}(x),
\]
and therefore
\[
F[f_t^a](x+a,v)=F[f_t^0](x,v).
\]
Writing the translated mean-field SDE with potential \(U_a\),
\[
d\bar X_t^a=\bar V_t^a\,dt,
\qquad
d\bar V_t^a=
\bigl(-\nabla U_a(\bar X_t^a)-\gamma \bar V_t^a+F[f_t^a](\bar X_t^a,\bar V_t^a)\bigr)\,dt
+\sigma\,dW_t,
\]
and setting \((\bar X_t^a,\bar V_t^a):=(\bar X_t^0+a,\bar V_t^0)\) gives exactly
the reference SDE with potential \(U_*\). Hence \(f_t^a=t_{a\#}f_t^0\) is the
translated mean-field solution. Applying this to the unique stationary law from
Theorem~\ref{thm:stationary-poc} yields \(t_{a\#}\pi_\infty^0=\pi_\infty^a\).

For \textnormal{(v)}, the pointwise derivatives are unchanged by translation in
the position variables. Consequently, for every smooth \(g\),
\[
\Gamma_{\mathrm{hypo}}^N(g\circ T_a^N,g\circ T_a^N)
=
\Gamma_{\mathrm{hypo}}^N(g,g)\circ T_a^N,
\]
and similarly for the mean-field carré du champ
\(\Gamma_{\mathrm{hypo}}^\infty\).
Using the pushforward identities from \textnormal{(iii)} and \textnormal{(iv)},
\[
\Ent_{\pi_N^a}(g^2)
=
\Ent_{\pi_N^0}\bigl((g\circ T_a^N)^2\bigr),
\qquad
\int \Gamma_{\mathrm{hypo}}^N(g,g)\,d\pi_N^a
=
\int \Gamma_{\mathrm{hypo}}^N(g\circ T_a^N,g\circ T_a^N)\,d\pi_N^0,
\]
and likewise
\[
\Ent_{\pi_\infty^a}(g^2)
=
\Ent_{\pi_\infty^0}\bigl((g\circ t_a)^2\bigr),
\qquad
\int \Gamma_{\mathrm{hypo}}^\infty(g,g)\,d\pi_\infty^a
=
\int \Gamma_{\mathrm{hypo}}^\infty(g\circ t_a,g\circ t_a)\,d\pi_\infty^0.
\]
Therefore the finite-\(N\) and mean-field logarithmic Sobolev inequalities for
the translated family are identical to those of the reference family, with the
same constants. The entropy contraction constants are the same because the
semigroup covariance in \textnormal{(ii)} conjugates the translated dynamics to
the reference dynamics.
\end{proof}

\begin{corollary}[Thermodynamic limit followed by wall-to-infinity recentering]\label{cor:wall-infinity}
Let \((a_R)_{R\ge1}\subset\R^d\) satisfy \(\|a_R\|\to\infty\). For each \(R\),
consider the translated confining potential \(U_{a_R}\). Then for every bounded
continuous observable \(\Phi:\R^{2d}\to\R\),
\[
\int_{\R^{2d}} \Phi(x-a_R,v)\,\pi_{N,1}^{a_R}(dx,dv)
=
\int_{\R^{2d}} \Phi(x,v)\,\pi_{N,1}^0(dx,dv)
\]
for every \(N\), and
\[
\int_{\R^{2d}} \Phi(x-a_R,v)\,\pi_\infty^{a_R}(dx,dv)
=
\int_{\R^{2d}} \Phi(x,v)\,\pi_\infty^0(dx,dv)
\]
for every \(R\). Consequently the two limits commute:
\[
\lim_{R\to\infty}\lim_{N\to\infty}
\int \Phi(x-a_R,v)\,\pi_{N,1}^{a_R}(dx,dv)
=
\lim_{N\to\infty}\lim_{R\to\infty}
\int \Phi(x-a_R,v)\,\pi_{N,1}^{a_R}(dx,dv)
=
\int \Phi(x,v)\,\pi_\infty^0(dx,dv).
\]
\end{corollary}

\begin{proof}
The first identity is exactly the pushforward relation
\((t_{-a_R})_\#\pi_{N,1}^{a_R}=\pi_{N,1}^0\), which follows from
Theorem~\ref{thm:translated-family}\textnormal{(iii)} by taking first marginals.
The second identity is
\((t_{-a_R})_\#\pi_\infty^{a_R}=\pi_\infty^0\), which is
Theorem~\ref{thm:translated-family}\textnormal{(iv)}.
The commutation of the \(N\to\infty\) and \(R\to\infty\) limits is then
immediate, because after recentering the integral is independent of \(R\), and
the \(N\to\infty\) limit is Theorem~\ref{thm:stationary-poc}.
\end{proof}

\section{Discussion}

The preceding analysis separates three logically distinct ingredients.

\begin{enumerate}[leftmargin=1.5em]
\item The kinetic backbone supplies the hypocoercive entropy gap.
\item The row-normalized viscous coupling is an $N$-uniform, dissipative
perturbation on relative velocities.
\item The velocity squashing map is a deterministic observation-level
post-processing step that preserves all probability metrics relevant to the
convergence theory.
\end{enumerate}

This decomposition is exactly what makes the kinetic core easier than the full
fractal gas. Once cloning and adaptive diffusion are removed, the only genuinely
new interaction is the viscous velocity alignment, and its role is stabilizing
rather than destabilizing.

Two further extensions are recorded separately.

\begin{itemize}[leftmargin=1.5em]
\item The adaptive geometric extension, in which isotropic noise is replaced by
state-dependent anisotropic diffusion.
\item The full algorithmic operator, in which the kinetic stage is coupled to
cloning, companion selection, and the fitness potential.
\end{itemize}

Those additions require a separate perturbative argument. The paper is
the clean kinetic backbone on which such extensions are built. At the numerical
level, the continuous kinetic stage can be paired with splitting
integrators for Langevin dynamics, while the smooth cap is applied only at
observation times \cite{leimkuhler2015molecular}.

\bibliographystyle{plain}
\bibliography{references}

\end{document}
